% concatenated from main.tex
\documentclass[11pt]{article}
\usepackage[T1]{fontenc}
\usepackage{lmodern}
\usepackage[margin=1in]{geometry}
\usepackage{amsmath,amssymb,amsthm,mathtools}
\usepackage{booktabs,graphicx}
\usepackage{float}
\usepackage{microtype}
\usepackage{xcolor}
\usepackage[colorlinks=true,allcolors=blue!55!black]{hyperref}
\usepackage{enumitem}
\setlist{nosep}
\newtheorem{theorem}{Theorem}[section]
\newtheorem{proposition}[theorem]{Proposition}
\newtheorem{lemma}[theorem]{Lemma}
\newtheorem{corollary}[theorem]{Corollary}
\theoremstyle{definition}
\newtheorem{assumption}[theorem]{Assumption}
\newtheorem{remark}[theorem]{Remark}
\newcommand{\R}{\mathbb R}
\newcommand{\E}{\mathbb E}
\newcommand{\dd}{\mathrm d}
\newcommand{\KL}{\operatorname{KL}}
\newcommand{\Id}{I}
\newcommand{\cH}{\mathcal H}
\newcommand{\cI}{\mathcal I}
\newcommand{\eps}{\varepsilon}
\DeclareMathOperator{\osc}{osc}
\title{Global Convergence of Third-Order Langevin Dynamics for Non-Convex Optimization via Simulated Annealing}
\author{
Yingli Wang\thanks{School of Mathematical Sciences, Fudan University, Shanghai, People's Republic of China; \texttt{yingliwang@fudan.edu.cn}}
\and
Lingjiong Zhu\thanks{Department of Mathematics, Florida State University, Tallahassee, Florida, United States of America; \texttt{zhu@math.fsu.edu}}}
\date{\today}

\begin{document}
\maketitle

\begin{abstract}
We study global convergence guarantees 
of third-order Langevin dynamics for non-convex optimization via simulated annealing
with fixed friction and decreasing noise.  An explicit three-block distorted entropy
transfers dissipation from the noisy auxiliary variable to the full state.
Under dissipativity, regularity, and low-temperature functional-inequality
assumptions, logarithmic cooling drives the objective values to the global
minimum in probability at the barrier-controlled kinetic rate.  For the
exact-force-integral and midpoint three-stage discretizations, polynomially
decreasing steps preserve this rate on the physical time scale.  The cubic
local endpoint estimate gives a less restrictive sufficient step-size condition
than the available frozen-force kinetic result.  A comparison with the
one-gradient UBU integrator shows how its centered stochastic local error leads,
under the same strong-coupling analysis, to a smaller sufficient iteration
exponent. Numerical experiments
are conducted to illustrate our theory. For a double well objective, third-order Langevin terminal-success point
estimates are higher than UBU at both a common horizon and an equal gradient
budget. For a high-dimensional nonconvex neural-network objective using synthetic data,
independently tuned UBU and third-order Langevin schemes both outperform overdamped
Langevin dynamics; the third-order Langevin point estimate is higher.  
For the same neural-network objective on real data, we show the same point-estimate
ordering for best-basin probability and post-quench test accuracy.\footnote{Numerical
code and associated experiment results are publicly available at
\url{https://github.com/gagawjbytw/simulated-annealing-third-order-langevin}.}
\end{abstract}

\section{Introduction}
Consider the non-convex optimization problem:
\begin{equation}\label{opt:problem}
\min_{x\in\mathbb{R}^{d}}U(x),
\end{equation}
where $U:\mathbb{R}^{d}\rightarrow\mathbb{R}$ is the given objective. 
The optimization problem \eqref{opt:problem} appears
in many applications in machine learning, including
empirical risk minimizations in deep learning where $U$ is typically non-convex.
Many algorithms have been used to solve the optimization problem \eqref{opt:problem}. 
Among these, gradient descent, stochastic gradient and their variance-reduced or momentum-based variants come with guarantees for finding a local minimizer or a stationary point for non-convex problems. 
While convergence to a local minimum can be satisfactory in some settings 
\cite{ge2017learning,du2017gradient}, 
methods with global convergence guarantees are also desirable and preferable in many other settings \cite{hazan2016graduated,simsekli-async-MCMC-18}.

\paragraph{Simulated annealing.}
Simulated annealing replaces direct descent in a nonconvex landscape by a
time-inhomogeneous stochastic dynamics whose instantaneous equilibrium is a Gibbs
law. The most classical one is based on the \textit{overdamped Langevin dynamics} (OLD): 
\begin{align}\label{eqn:OLD}
\dd X_{t}=-\nabla U(X_{t})\dd t+\sqrt{2\varepsilon_{t}}\dd B_{t},
\end{align}
where $B_{t}$ is a standard Brownian motion and $\varepsilon_{t}$ is a temperature parameter
that depends on $t$.
The simulated annealing method was introduced as a computational principle in
\cite{KirkpatrickGelattVecchi1983}. Its diffusion formulation was developed in
\cite{GemanHwang1986,ChiangHwangSheu1987,Royer1989}: at a fixed temperature
$\eps_{t}\equiv\eps$, the stationary position law is proportional to $e^{-U/\eps}$, while a decreasing
temperature must retain enough thermal energy for the process to cross the deepest
relevant barriers. This competition leads to the logarithmic cooling scale. Spectral
gap and Sobolev methods identify the corresponding low-temperature obstruction
\cite{HolleyStroock1988,HolleyKusuokaStroock1989}, and Miclo's free-energy analysis
\cite{Miclo1992} gives quantitative convergence. Later work extended the theory to
weaker confinement and slowly growing potentials
\cite{Zitt2008,FournierTardif2021}. Tang and Zhou \cite{TangZhou2023} obtained
convergence rates for continuous-time overdamped annealing using
Eyring--Kramers estimates. Tang, Wu, and Zhou \cite{TangWuZhou2024} established
corresponding rates and step-size conditions for its discrete-time counterpart.
These results identify the landscape-dependent critical cooling constant and
give discrete-time convergence guarantees that control accumulated
discretization error on the same physical-time scale.

\paragraph{Kinetic annealing and hypocoercivity.}
There have been many variants of overdamped Langevin dynamics, 
and among these, the momentum-based kinetic or \textit{underdamped Langevin dynamics} (ULD) is
one of the most popular:
\begin{align}\label{eqn:ULD}
 \dd X_t&=V_t\dd t,\qquad
 \dd V_t=(-\nabla U(X_t)-\gamma V_t)\dd t
                 +\sqrt{2\gamma\eps_t}\dd B_t,
\end{align}
where $\gamma>0$ is the friction coefficient, $\varepsilon_{t}$ is a temperature parameter
that depends on $t$,
and $B_{t}$ is a standard $d$-dimensional Brownian motion.
At a fixed temperature
$\eps_{t}\equiv\eps$, the stationary law is proportional to $e^{-\frac{1}{\eps}(U(x)+\frac{1}{2}|v|^{2})}$
whose position marginal law is proportional to $e^{-U/\eps}$, 
the same as in the overdamped case.
Adding momentum can improve mixing at a fixed temperature, but it makes the noise
degenerate: entropy is dissipated directly only in the velocity variable. The
hypocoercive framework transfers this dissipation to the full state through the
transport structure; see
\cite{ArnoldEtAl2001,Herau2007,Villani2009}. Lyapunov and coupling methods give
ergodicity and quantitative contraction for kinetic Langevin dynamics under broad
conditions \cite{MattinglyStuartHigham2002,EberleGuillinZimmer2019}, while
nonasymptotic sampling bounds and discretizations are developed in
\cite{ChengEtAl2018,MaEtAl2021}.

For simulated annealing, Monmarch\'e \cite{Monmarche2018} introduced a distorted
entropy adapted to a kinetic process in which the friction grows as the temperature
decreases. Journel and Monmarch\'e \cite{JournelMonmarche2022} treated the
fixed-friction, vanishing-noise formulation under substantially weaker potential
conditions and also proved failure results for cooling schedules that are too fast.
He, Tan, and Wu \cite{HeTanWu2024} obtained quantitative convergence rates for that
formulation and for a particular time discretization. Their calculation adapts the
matrix-valued entropic method of Ma et al.\ \cite{MaEtAl2021}; related discrete
Hamiltonian entropy estimates are given in \cite{Monmarche2024}. At low temperature,
the rate is controlled by a logarithmic Sobolev constant. Landscape decompositions
such as those of Menz and Schlichting \cite{MenzSchlichting2014} provide its
exponential scale, while the
Holley--Stroock perturbation principle \cite{HolleyStroock1987} provides a simple
explicit bound for bounded perturbations of a quadratic potential. Time-dependent
Fisher-information techniques offer another perspective on inhomogeneous diffusions
\cite{FengZuoLi2024}.

\paragraph{Third-order and generalized Langevin dynamics.}
Generalized Langevin equations introduce memory or auxiliary variables while
preserving the desired Gibbs marginal. Their Markovian embeddings and scaling
limits have been studied, for example, in
\cite{OttobrePavliotis2011,PavliotisStoltzVaes2021}. The \textit{third-order Langevin dynamics} of
Mou et al. \cite{MouEtAl2021}, equation (8), uses two auxiliary variables and places
noise only in the last component:  
\begin{equation}\label{eqn:third:order}
\begin{aligned}
 \dd X_t&=V_t\dd t,\\
 \dd V_t&=(-\nabla U(X_t)+\lambda Z_t)\dd t,\\
 \dd Z_t&=(-\lambda V_t-\gamma Z_t)\dd t
                 +\sqrt{2\gamma\eps_t}\,\dd B_t,
\end{aligned}
\end{equation}
where $\lambda,\gamma>0$, $\varepsilon_{t}$ is a temperature parameter
that depends on $t$, and 
$B_{t}$ is a standard $d$-dimensional Brownian motion.
The resulting smoother force integration supports
accelerated sampling guarantees for strongly log-concave targets and a specific
three-stage numerical scheme. Those guarantees apply under strong convexity and a
positive lower Hessian bound. The third-order Langevin dynamics \eqref{eqn:third:order} 
and its generalizations have been previously studied in the context
of sampling \cite{MouEtAl2021,DangGurbuzbalabanIslamYaoZhu2025,high-order-Liu-2025}. 
The nonconvex global-minimization setting \eqref{opt:problem} considered
in this paper using simulated annealing analysis to the best of our knowledge is novel
for the third-order Langevin dynamics \eqref{eqn:third:order}.

\paragraph{Global optimization.}
Several related lines of work pursue global optimization with Langevin-type
processes through fixed-temperature, discretized, and stochastic gradient
formulations.
Nonasymptotic guarantees for fixed-temperature or discretized diffusions appear in
\cite{RaginskyRakhlinTelgarsky2017,ErdogduMackeyShamir2018,xu2018global,Chau2019,Zhang2019}; momentum-based
stochastic-gradient Hamiltonian methods are analyzed in
\cite{GaoGurbuzbalabanZhu2022,Chau2022}. Temperature selection has also been formulated as
a state-dependent stochastic control problem \cite{GaoXuZhou2022}. More recently,
Herty and Zanella \cite{HertyZanella2026} proposed a feedback-controlled temperature
for an interacting kinetic particle model and proved entropy decay for its
Boltzmann-type mean-field description. Together, these approaches clarify the
broader role of thermal exploration. The present model follows a logarithmically
cooled diffusion with its associated physical time scale.

\paragraph{Eyring--Kramers laws and metastability.}
At a fixed low temperature, the Arrhenius exponent determines the dominant
transition scale, whereas the Eyring--Kramers law also identifies the sharp
prefactor. For reversible overdamped Langevin dynamics, the potential-theoretic
approach of Bovier et al.\ \cite{BovierEtAl2004} gives the classical formula
for capacities and mean exit times. Bouchet and Reygner
\cite{BouchetReygner2016} treat non-reversible Langevin dynamics, while the
non-selfadjoint capacity principles of Landim, Mariani, and Seo
\cite{LandimMarianiSeo2019} and the Eyring--Kramers formula of Lee and Seo
\cite{LeeSeo2022} develop the corresponding nonreversible Gibbs theory.

On the
spectral side, Bony, Le Peutrec, and Michel \cite{BonyLePeutrecMichel2025}
establish Eyring--Kramers formulas for a broad class of Fokker--Planck type
operators with Gibbs stationary measures.
The hypoelliptic underdamped process requires additional ideas because the
standard elliptic potential theory does not apply directly. Lee, Ramil, and Seo
\cite{LeeRamilSeo2026} recently derived the mean-transition-time
Eyring--Kramers law for underdamped Langevin dynamics in a double-well
landscape. This result provides the direct second-order benchmark for the
third-order metastability theory described next.

Recent work develops a fixed-temperature metastability theory for the same
hypoelliptic chain. He et al.\ \cite{HeLiWangZhu2026} construct weak
equilibrium measures and capacities and prove capacity--hitting identities. Using
this framework, Wang and Zhu \cite{WangZhu2026} prove an Eyring--Kramers law for a
double-well potential with a unique index-one saddle. The usual Arrhenius exponent
is retained, but under matched kinetic normalizations the sharp transition-time
prefactor is strictly smaller than its underdamped counterpart. Their argument
extends the underdamped strategy of Lee, Ramil, and Seo \cite{LeeRamilSeo2026} to the step-three chain.
Together, these results provide fixed-temperature metastable transition
asymptotics. We develop the
corresponding time-inhomogeneous simulated annealing convergence theory below.

Standard diffusion annealing lowers the temperature by reducing noise while keeping
the drift fixed \cite{GemanHwang1986,ChiangHwangSheu1987,Royer1989}. For kinetic
Langevin annealing, relative entropy directly dissipates only the velocity component
of the Fisher information \cite[Section~2.2]{HeTanWu2024}. The same degeneracy is
present in the higher-order chain of Mou et al.\ \cite{MouEtAl2021}, where noise acts
only on the last component; when its amplitude is cooled, the coefficient of this
direct dissipation vanishes.

Chak, Kantas, and Pavliotis
\cite{ChakKantasPavliotis2023}, equation (1.3) and Theorem 2.7, instead analyze
\begin{equation}\label{eq:ckp}
 \dd X_t=Y_t\dd t,\qquad
 \dd Y_t=\left(-\nabla U(X_t)+\Lambda^\top Z_t\right)\dd t,\qquad
 \dd Z_t=\left(-\Lambda Y_t-T_t^{-1}A Z_t\right)\dd t+\Sigma\dd B_t,
\end{equation}
with $\Sigma\Sigma^\top=2A$, where $B_{t}$ is a standard $d$-dimensional Brownian motion.
This parametrization keeps the noise matrix fixed and increases the effective
friction $T_t^{-1}A$ as the temperature decreases. At each frozen temperature it
has the desired Gibbs law. The authors adopt this parametrization to follow the
hypocoercive framework of Monmarch\'e \cite{Monmarche2018}; the additional degeneracy
of the third-order chain is handled by a tailored distorted entropy and approximation
by nondegenerate diffusions. They prove a continuous-time annealing rate and present
numerical experiments.

We combine the fixed-friction annealing strategy of He, Tan, and Wu
\cite{HeTanWu2024} with the third-order diffusion and three-stage scheme of Mou et
al.\ \cite{MouEtAl2021}. 
He, Tan, and Wu \cite{HeTanWu2024} return to the fixed-friction formulation with
vanishing noise. They replace relative entropy by a temperature-dependent distorted
entropy whose entropy weight is of order $\eps^{-1}$. A matrix hypocoercivity
calculation adapted from Ma et al.\ \cite{MaEtAl2021} transfers the
velocity-direction dissipation to the position variable, while the low-temperature
\textit{logarithmic Sobolev inequality} (LSI) converts the resulting Fisher-information bound into frozen-temperature
contraction. Uniform
moment bounds control the temperature-variation term, so logarithmic cooling
preserves the contraction needed for annealing. We apply this strategy to the
additional variable $Z$: an explicit
three-block Fisher-information matrix transfers dissipation from $Z$ through $V$
to $X$ while $\lambda$ and $\gamma$ remain fixed and the noise amplitude decreases.
We also establish convergence for the associated time-discrete schemes.
Higher-order annealed Langevin methods have been used for posterior sampling in
linear inverse problems \cite{ZilbersteinEtAl2024}; here we study the nonconvex
zero-temperature limit.

\paragraph{Contribution.}
Our main contributions can be summarized as follows.

\begin{enumerate}[label=(\roman*)]
\item An explicit three-block dissipation certificate
(Lemma~\ref{lem:matrix}) yields the continuous-time annealing rate in
Theorem~\ref{thm:main}.  Lemma~\ref{lem:closure} extends the entropy argument to
rougher initial laws, and Proposition~\ref{prop:restart} provides the uniform
restart estimate needed for discrete comparison.

\item Local coupling, moment control, grid-point regularization, and
exact-transition smoothing yield Theorem~\ref{thm:discrete} for the exact-force
and midpoint schemes.  Corollary~\ref{cor:iterationgain} gives the resulting
rate-preserving iteration gain.  Applying the same recent-window argument to
UBU splitting scheme in Theorem~\ref{thm:ubu} shows how the additional auxiliary level delays
noise propagation to position and gives the third-order midpoint scheme a
sharper sufficient discretization bound.

\item Three reproducible annealing experiments compare overdamped, underdamped,
and third-order Langevin dynamics after validation-set tuning and independent
evaluation.  A tilted double well uses theory-driven cooling and step grids; an
$88$-dimensional regularized two-layer tanh objective on synthetic data uses independent two-stage
parameter tuning; and a $248$-dimensional example evaluates the same
nonconvex objective on real data and a fixed held-out split.  Both kinetic
schemes outperform overdamped Langevin dynamics in the synthetic and real-data
neural-network examples.  The third-order Langevin best-basin point estimates are higher
than UBU in both examples; see Section~\ref{sec:numerics}.
\end{enumerate}

The resulting exponents retain the same landscape barrier $D$ that appears in the
low-temperature LSI. At the continuous-time level, Theorem~\ref{thm:main} matches
the polynomial exponent of the kinetic annealing rate. The distinct quantitative
gain proved here appears in the discretization: the $\mathcal{O}(h^3)$ local endpoint
comparison permits a less restrictive sufficient step-size threshold than the
kinetic result of He, Tan, and Wu \cite{HeTanWu2024}. Sharp metastable prefactors offer a
complementary way to distinguish the dynamics beyond the common barrier exponent.
Table~\ref{tab:introdiscretecomparison} summarizes the comparison for discretized algorithms.
Write $\vartheta:=D/E<1$ and
\[
 r_{\rm c}(\delta):=\min\left\{\frac{\delta}{E},
                         \frac{1-\vartheta}{2}\right\},
\]
the common continuous-time exponent. For polynomial grids
$h_k\asymp s_k^{-a}$, where $s_k=\tau_k+t_0$ is shifted accumulated physical
time, the currently available sufficient conditions give the following
rate-preserving thresholds
and iteration exponents.
The third column in Table~\ref{tab:introdiscretecomparison} follows from
$\tau_k=\sum_{j=0}^{k-1}h_j$ and $h_j\asymp s_j^{-a}$, which imply
$s_k\asymp k^{1/(1+a)}$. Hence a physical-time estimate
$\mathbb P(U(X_k)>\delta)\lesssim s_k^{-r_{\rm c}(\delta)+\alpha}$
becomes, after an arbitrarily small adjustment of $\alpha$,
\[
 \mathbb P(U(X_k)>\delta)
 \lesssim k^{-r_{\rm c}(\delta)/(1+a)+\alpha},
\]
where $X_{k}$ is the $k$-th iterate of the corresponding algorithm and without loss of generality we assume that
$U$ has minimum zero.
The fourth column in Table~\ref{tab:introdiscretecomparison} uses the smallest value of $a$ allowed by the
rate-preserving condition in the second column, since this maximizes the
resulting iteration exponent.

\begin{table}[ht]
\centering
\small
\renewcommand{\arraystretch}{1.25}
\begin{tabular}{@{}p{0.27\textwidth}p{0.27\textwidth}p{0.18\textwidth}p{0.18\textwidth}@{}}
\toprule
Method and source & Sufficient rate-preserving grid condition
& Physical time after\newline $k$ steps & Proved iteration\newline exponent \\
\midrule
OLD (Euler)\newline
\cite{TangWuZhou2024}
& $h_ks_k=\mathcal{O}(1)$; equivalently $a\geq1$
& $s_k\asymp k^{1/2}$
& $r_{\rm c}(\delta)/2$ \\
ULD (frozen-force)\newline
\cite[Theorem~2.8]{HeTanWu2024}
& $h_k\sqrt{s_k}=\mathcal{O}(1)$; equivalently $a\geq1/2$
& $s_k\asymp k^{2/3}$
& $2r_{\rm c}(\delta)/3$ \\
ULD (UBU)\newline
(Corollary~\ref{cor:ubuiteration})
& $a\geq(1+\vartheta)/3$
& $s_k\asymp k^{3/(4+\vartheta)}$
& $3r_{\rm c}(\delta)/(4+\vartheta)$ \\
Third-order LD (midpoint)\newline
(Corollary~\ref{cor:iterationgain})
& $a\geq a_*=(1+\vartheta)/4$
& $s_k\asymp k^{4/(5+\vartheta)}$
& $4r_{\rm c}(\delta)/(5+\vartheta)$ \\
\bottomrule
\end{tabular}
\caption{Comparison of available and presently proved sufficient grid conditions
for discrete annealing. Arbitrarily small losses in the
iteration exponents are suppressed.}
\label{tab:introdiscretecomparison}
\end{table}

Since $\vartheta<1$, the third-order Langevin midpoint exponent is strictly larger
than each of the other three exponents in
Table~\ref{tab:introdiscretecomparison}. The two underdamped bounds are not
uniformly ordered: the UBU exponent is larger than the frozen-force exponent
when $\vartheta<1/2$, equal to it when $\vartheta=1/2$, and smaller when
$\vartheta>1/2$. This comparison concerns proved sufficient conditions rather
than lower bounds or an intrinsic ordering of the integrators.

\paragraph{Organization of the paper.}
Section~\ref{sec:model} introduces the model and assumptions.
Section~\ref{sec:continuous} proves the continuous-time result, and
Section~\ref{sec:discrete} treats the numerical schemes and their convergence.
Section~\ref{sec:numerics} gives a numerical illustration, and
Section~\ref{sec:conclusion} concludes. Appendix~\ref{sec:ubuappendix} contains the UBU
comparison theorem and its proof.

%%%%%%%%%%%%%%%%%%%%%%%%%%%%%%%%%%%%%%%%%%%%%%%%%%%%%%%%
\section{Model Setup and Assumptions}\label{sec:model}

For fixed $\lambda,\gamma>0$, we recall from \eqref{eqn:third:order} the \textit{third-order Langevin dynamics} 
$S_t=(X_t,V_t,Z_t)\in\R^{3d}$, that satisfies:
\begin{equation}\label{eq:model}
\begin{aligned}
 \dd X_t&=V_t\dd t,\\
 \dd V_t&=(-\nabla U(X_t)+\lambda Z_t)\dd t,\\
 \dd Z_t&=(-\lambda V_t-\gamma Z_t)\dd t
                 +\sqrt{2\gamma\eps_t}\,\dd B_t,
\end{aligned}
\end{equation}
where $\varepsilon_{t}$ is a temperature parameter
that depends on $t$ and $B_{t}$ is a standard $d$-dimensional Brownian motion.
Define the energy and invariant Gibbs law by
\begin{equation}\label{eq:gibbs}
 H(x,v,z):=U(x)+\frac12|v|^2+\frac12|z|^2,
 \qquad \pi_\eps(\dd x\,\dd v\,\dd z)
   :=\mathcal Z_\eps^{-1}e^{-H(x,v,z)/\eps}\dd x\,\dd v\,\dd z.
\end{equation}
For $\eps>0$, define the frozen-temperature generator
\begin{equation}\label{eq:generator}
 L_\eps:=v\cdot\nabla_x+(-\nabla U+\lambda z)\cdot\nabla_v
       +(-\lambda v-\gamma z)\cdot\nabla_z+\gamma\eps\Delta_z.
\end{equation}
The dynamics \eqref{eq:model} is time-inhomogeneous, with generator
$L_{\eps_t}$ at time $t$; when $\eps_t\equiv\eps$ is constant, it is generated by $L_\eps$.
The transport part preserves $H$ and has zero divergence; the last friction/noise
pair preserves $\mathcal N(0,\eps\Id_d)$. Hence $\pi_\eps$ is invariant for every
fixed $\eps>0$.

To recover the notation of Mou et al. \cite{MouEtAl2021}, write their potential as $F$ and their
smoothness parameter as $L_0$. Set $U=F/L_0$, $\lambda=\gamma_{\rm Mou}$,
$\gamma=\xi_{\rm Mou}$, and $\eps_t=T_t/L_0$. Equation \eqref{eq:model} then has
force $-\nabla F/L_0$ and noise $\sqrt{2\xi_{\rm Mou}T_t/L_0}$, with position
law proportional to $e^{-F/T_t}$ at frozen temperature. This makes the
normalization of both the potential and the cooling constant explicit.
Replacing $F$ by $F/T_t$ while keeping $L_0$ fixed would define a different
time-dependent force.
Likewise, Chak, Kantas, and Pavliotis
\cite[Theorem~2.7]{ChakKantasPavliotis2023}, applied to \eqref{eq:ckp}, assume a fixed
matrix $A$, whereas $A=\gamma\eps_t\Id_d$ varies with time under the present
cooling convention.

\begin{assumption}[Potential and cooling]\label{ass:potential}
The function $U\in C^4(\R^d)$ has minimum zero, bounded Hessian
$\|\nabla^2U\|_{\rm op}\leq K$, and satisfies
$x\cdot\nabla U(x)\geq r|x|^2-m$ for some $r>0,m\geq0$.
Its third and fourth derivatives have at most polynomial growth.
The temperature $\eps_t\in C^1([0,\infty);(0,\infty))$ is non-increasing in $t$. For
some $t_0>1$ and $E>0$,
\begin{equation}\label{eq:coolingassumption}
 \eps_t\log(t+t_0)\longrightarrow E,
 \qquad
 |\eps_t'|(1+\eps_t^{-3})
 \leq C\frac{\log(t+t_0)}{t+t_0}
 \quad\text{for all sufficiently large }t.
\end{equation}
\end{assumption}

The canonical schedule $\eps_t=E/\log(t+t_0)$ satisfies
\eqref{eq:coolingassumption}.

\begin{assumption}[Continuous-time initial law]\label{ass:initial}
The initial law $\nu_0$ of \eqref{eq:model} is absolutely continuous, has a finite
second moment, and finite
relative Fisher information with respect to $\pi_{\eps_0}$:
\begin{equation}\label{eq:initialfisher}
 \cI(\nu_0\mid\pi_{\eps_0})
 :=\int_{\R^{3d}}
 \left|\nabla\log\frac{\dd\nu_0}{\dd\pi_{\eps_0}}\right|^2\dd\nu_0
 <\infty.
\end{equation}
\end{assumption}

\begin{assumption}[LSI barrier]\label{ass:lsi}
Let $Z_\eps:=\int_{\R^d}e^{-U(y)/\eps}\dd y$ and
$\mu_\eps(\dd x)=Z_\eps^{-1}e^{-U(x)/\eps}\dd x$ be the position
Gibbs measure. There is $D\geq0$ such that for every $\eta>0$ there are
$C_\eta<\infty$ and $\eps_\eta>0$ for which
\begin{equation}\label{eq:lsi}
 \KL(\nu\mid\mu_\eps)\leq C_{\rm LS}^{\rm pos}(\eps)
       \cI(\nu\mid\mu_\eps),
 \qquad C_{\rm LS}^{\rm pos}(\eps)\leq C_\eta e^{(D+\eta)/\eps},
 \quad 0<\eps\leq\eps_\eta.
\end{equation}
Here $\cI(\nu\mid\mu):=\int_{\R^d}|\nabla\log(\dd\nu/\dd\mu)|^2\dd\nu$, and we assume
$D<E$ where $E$ is the cooling
constant from Assumption~\ref{ass:potential}.
\end{assumption}

\begin{assumption}[Additional discrete regularity]\label{ass:discrete}
For the discrete-time results, we further assume that
$\|\nabla^3U\|_\infty<\infty$ and that the initial law for the numerical
scheme has a finite fourth moment. The initial law may be singular, including
a point mass corresponding to deterministic initialization.
\end{assumption}

In Assumption~\ref{ass:lsi}, $D$ denotes the LSI barrier. The velocity and
auxiliary-variable factors $\mathcal N(0,\eps\Id_d)(\dd v)$ and
$\mathcal N(0,\eps\Id_d)(\dd z)$
satisfy the LSI with
constant $\eps/2$. For each $\eta>0$ and
$0<\eps\leq\eps_\eta$, tensorization therefore gives the full Gibbs law
$\pi_\eps$ an admissible LSI constant
\begin{equation}\label{eq:tensorizedlsi}
 C_{\rm LS}(\eps)=\max\{C_{\rm LS}^{\rm pos}(\eps),\eps/2\},
\end{equation}
which has the same exponential bound as in \eqref{eq:lsi}. Under the landscape
hypotheses of Menz and Schlichting \cite{MenzSchlichting2014}, $D$ can be identified
with the corresponding critical depth. The bounded-Hessian and dissipativity conditions
support the dynamical estimates, while the landscape hypotheses determine the
low-temperature Gibbs mixing across multiple wells.

\begin{remark}[Comparison with related assumptions in the literature]\label{rem:assumptioncomparison}
He, Tan, and Wu \cite[Assumptions~2.1 and 2.3]{HeTanWu2024} assume that
$U\in C^\infty$, all derivatives of $U$ have at most polynomial growth,
$\nabla U$ is globally Lipschitz, and $U$ is dissipative. They also require
finitely many local minima with nondegenerate Hessians. These potential
conditions imply the potential regularity in Assumption~\ref{ass:potential};
their LSI estimate \cite[Proposition~3.1]{HeTanWu2024} gives
Assumption~\ref{ass:lsi} with $D$ equal to their critical depth $E_*$. Our
direct LSI formulation replaces finiteness and nondegeneracy conditions on the
local minima by the functional inequality itself. Their initial-law condition
includes a smooth density, finite raw Fisher information, and moments of every
order. The analogous condition on our three-component state implies
Assumption~\ref{ass:initial}, since bounded $\nabla^2U$ and a second moment
turn finite raw Fisher information into finite relative Fisher information.
Assumption~\ref{ass:initial} thus allows less regular initial laws. Their
eventually exact schedule $\eps_t=E/\log t$ satisfies the asymptotic and
derivative bounds in \eqref{eq:coolingassumption}; we additionally require
global $C^1$ regularity of the schedule. For their discrete-time theorem,
He, Tan, and Wu \cite[Theorem~2.7]{HeTanWu2024} additionally assume that
$\nabla^2U$ is globally Lipschitz, which, under their smoothness assumption and
up to the choice of tensor norm, corresponds to
$\|\nabla^3U\|_\infty<\infty$ in Assumption~\ref{ass:discrete}. Their
requirement of moments of every order implies our finite-fourth-moment
condition. However, their smooth-density assumption excludes deterministic
initialization, whereas Assumption~\ref{ass:discrete} allows singular initial
laws.

Chak, Kantas, and Pavliotis
\cite[Assumptions~1--4]{ChakKantasPavliotis2023} assume a smooth potential with
bounded Hessian, dissipativity, and a quadratic gradient bound, together with
either a two-sided quadratic comparison or a nondegenerate Morse landscape with
additional saddle structure. These alternatives yield an LSI barrier
$\widehat{E}$, which corresponds to $D=\widehat{E}$ in
Assumption~\ref{ass:lsi}. The two formulations emphasize different structures:
their landscape assumptions specify Morse and saddle geometry, while ours uses
$C^4$ regularity together with a direct LSI. Their initial density is bounded and
smooth, whereas Assumption~\ref{ass:initial} accommodates nonsmooth and unbounded
initial densities. The cooling conditions likewise emphasize different controls.
They allow any $T_t\to0$ satisfying $T_t\geq E/\log t$ and
$|T_t'|\leq C/t$, while ours combines monotonicity, the asymptotic relation
$\eps_t\log(t+t_0)\to E$, and the sharper derivative control in
\eqref{eq:coolingassumption}. Their dynamics cools through
temperature-dependent friction with fixed noise, whereas \eqref{eq:model} keeps
the friction fixed and decreases the noise amplitude.
\end{remark}

%%%%%%%%%%%%%%%%%%%%%%%%%%%%%%%%%%%%%%%%%%%%%%%%%%%%
\section{Continuous-Time Annealing}\label{sec:continuous}

\subsection{Dissipation and initial-density approximation}
For $h:=\dd\nu/\dd\pi_\eps$, let $w:=\nabla\log h$, and write
\begin{equation}\label{eq:functional}
 \cH_\eps(\nu):=\cI_M(\nu\mid\pi_\eps)
                  +\frac{A_0}{\eps}\KL(\nu\mid\pi_\eps),
 \qquad \cI_M:=\int_{\R^{3d}} w^\top M w\,\dd\nu,
\end{equation}
where $\KL(\nu\mid\pi_\eps)$ denotes the Kullback-Leibler (KL) divergence:
\begin{equation*}
\KL(\nu\mid\pi_\eps)
:=\int_{\mathbb{R}^{3d}}\log h\,\dd\nu
=\int_{\mathbb{R}^{3d}}h\log h\,\dd\pi_\eps,
\end{equation*}
and
\begin{align}\label{eq:matrix}
A_0:=\frac{2\lambda b+2(q_1^2+q_2^2)+1}{\gamma},
\qquad
M:=\begin{pmatrix}1&1&0\\1&2&b\\0&b&\varsigma\end{pmatrix}\otimes\Id_d,
\end{align}
with
\begin{align}
 b:=\frac{K+1+(K+2)^2}{\lambda},
 \quad \varsigma:=b^2+1,
 \quad q_1:=b-\lambda,
 \quad q_2:=\lambda\varsigma-2\lambda+\gamma b.
\end{align}

The coefficients in \eqref{eq:functional}--\eqref{eq:matrix} are chosen for the
frozen-temperature identities \eqref{eq:fisherentropyderivatives}. In particular,
\eqref{eq:entropyderivative} shows that KL dissipation acts only on $w_z$ and
carries a factor $\gamma\eps$. The weight $A_0/\eps$ removes this temperature
factor, while the off-diagonal entries of $M$ transfer dissipation through
$Z\to V\to X$. The choice $\varsigma=b^2+1$ ensures $M\succ0$, and $b$ leaves the
positive margin $2\lambda b-2K-(K+2)^2=2+(K+2)^2$ in the $(x,v)$ block for
every $\|Q\|_{\rm op}\leq K$. The quantities $q_1,q_2$ are the resulting
$(x,z)$ and $(v,z)$ mixed coefficients in the matrix of
\eqref{eq:certificate}; $A_0$ is chosen so that
$\gamma A_0-2\lambda b-2(q_1^2+q_2^2)=1$ after these terms are absorbed.
These are explicit sufficient choices for \eqref{eq:certificate}, rather than
optimized constants.

The next lemma provides the coercivity used in the frozen-temperature
dissipation identity \eqref{eq:frozenidentity} in the proof of
Theorem~\ref{thm:main}; it holds uniformly over all Hessians with
$\|\nabla^2U\|_{\rm op}\leq K$.

\begin{lemma}[Algebraic dissipation certificate]\label{lem:matrix}
The matrix $M$ in \eqref{eq:matrix} is positive definite. For every symmetric $Q$ with
$\|Q\|_{\rm op}\leq K$, define
\[
 J_Q:=\begin{pmatrix}0&-\Id_{d}&0\\Q&0&-\lambda\Id_{d}\\0&\lambda\Id_{d}&-\gamma\Id_{d}\end{pmatrix},
 \qquad P_z:=\operatorname{diag}(0,0,\Id_d).
\]
Then
\begin{equation}\label{eq:certificate}
 -J_QM-MJ_Q^\top+A_0\gamma P_z\succeq\frac12\Id_{3d}.
\end{equation}
\end{lemma}
\begin{proof}
The scalar matrix defining $M$ has leading principal minors
$1$, $1$, and $\varsigma-b^2=1$, so that Sylvester's criterion gives $M\succ0$.
Denote the matrix on the left of \eqref{eq:certificate} by $B_Q$. Direct
matrix multiplication gives
\[
 B_Q=\begin{pmatrix}
 2\Id_{d}&2\Id_{d}-Q&q_1\Id_{d}\\
 2\Id_{d}-Q&2\lambda b\Id_{d}-2Q&q_2\Id_{d}\\
 q_1\Id_{d}&q_2\Id_d&(2\gamma\varsigma-2\lambda b+A_0\gamma)\Id_{d}
 \end{pmatrix}.
\]
For $u,v\in\R^d$, the quadratic form of its top-left block satisfies
\begin{align*}
 2|u|^2+2u^\top(2\Id_{d}-Q)v+v^\top(2\lambda b\Id_d-2Q)v
 &\geq 2|u|^2-2(K+2)|u||v|+(2\lambda b-2K)|v|^2\\
 &\geq |u|^2+\left(2\lambda b-2K-(K+2)^2\right)|v|^2\\
 &=|u|^2+\left(2+(K+2)^2\right)|v|^2.
\end{align*}
Moreover, for $z\in\R^d$,
\[
 2(q_1u+q_2v)\cdot z
 \geq-\frac12\left(|u|^2+|v|^2\right)
      -2\left(q_1^2+q_2^2\right)|z|^2.
\]
Since
\[
 2\gamma\varsigma-2\lambda b+A_0\gamma-2\left(q_1^2+q_2^2\right)=2\gamma\varsigma+1,
\]
we obtain, for every $\xi=(u,v,z)\in\R^{3d}$,
\[
 \xi^\top B_Q\xi
 \geq\frac12\left(|u|^2+|v|^2\right)+(2\gamma\varsigma+1)|z|^2
 \geq\frac12|\xi|^2.
\]
This proves \eqref{eq:certificate}.
\end{proof}

For nonsmooth initial densities, the next lemma justifies the integrated
entropy identities by approximation.

\begin{lemma}[Justification by initial-density approximation]\label{lem:closure}
Under Assumption~\ref{ass:potential}, let the annealed third-order process
$S_t=(X_t,V_t,Z_t)$ in \eqref{eq:model} be started, or restarted
at a deterministic time $s\geq0$, from
an absolutely continuous law $\nu_s$ with finite relative Fisher information with
respect to $\pi_{\eps_s}$ and finite second moment. Let $\nu_t$ denote the
law of $S_t$ for $t\geq s$, and let $F(t):=\cH_{\eps_t}(\nu_t)$. Then
$F(t)<\infty$ for every finite $t\geq s$.
Moreover, there exist smooth, strictly positive probability laws
$\nu_s^{(n)}$, Gaussian outside compact sets, such that
\[
 \mathcal W_2\left(\nu_s^{(n)},\nu_s\right)\to0,\qquad
 \cH_{\eps_s}\left(\nu_s^{(n)}\right)\to F(s),\qquad
 \sup_{n\geq1}\int_{\R^{3d}}|y|^2\nu_s^{(n)}(\dd y)<\infty.
\]
Let $\nu_u^{(n)}$ be the laws of \eqref{eq:model} started from $\nu_s^{(n)}$
at time $s$, and set $F_n(u):=\cH_{\eps_u}\left(\nu_u^{(n)}\right)$.
Fix $T>s$.
Suppose that there exist $a,b\in L^1([s,T])$, with $b\geq0$, such that
\begin{equation}\label{eq:abstractcomparison}
 F_n'(u)+a(u)F_n(u)\leq b(u)
 \quad\text{for a.e. }u\in[s,T]\text{ and every }n\geq1.
\end{equation}
Then $F$ satisfies
\begin{equation}\label{eq:abstractintegrated}
 F(t)\leq e^{-\int_s^t a(u)\dd u}F(s)
 +\int_s^t e^{-\int_v^t a(u)\dd u}b(v)\dd v,
 \qquad s\leq t\leq T.
\end{equation}
\end{lemma}

\begin{remark}[Use of the approximation lemma]
Lemma~\ref{lem:closure} is used in the proof of
Theorem~\ref{thm:main} to pass the comparison inequality
\eqref{eq:ode} from regularized initial laws to laws satisfying
Assumption~\ref{ass:initial}. Its time-shifted form is used in the proof
of Proposition~\ref{prop:restart} to extend the regular-law estimate
\eqref{eq:restartode} to general restart laws covered by that proposition.
\end{remark}

\begin{proof}[Proof of Lemma~\ref{lem:closure}]
Let $T>s$ be as in the statement. By the deterministic time shift $r=t-s$, with
$\widetilde\eps_r=\eps_{s+r}$,
$\widetilde B_r=B_{s+r}-B_s$,
$\widetilde a(r)=a(s+r)$, and $\widetilde b(r)=b(s+r)$ when the comparison
hypothesis is used, we may take $s=0$ and relabel $T-s$ as $T$ throughout
the proof. For each $y\in\R^{3d}$, let $\Phi_t(y)$ denote the solution of
\eqref{eq:model} at time $t$ started from $y$ at the shifted time $0$, using
the same Brownian path for every $y$. Let $\mathfrak b$ denote the drift
vector field, and set
\[
 N_T:=\sup_{0\leq t\leq T}\left|\int_0^t\sqrt{2\gamma\eps_r}\,\dd B_r\right|.
\]
Write $Y_t=\Phi_t(y)$. Subtracting the additive noise path gives a
pathwise nonautonomous ODE.
Indeed, with
\[
 G_t:=\left(0,0,\int_0^t\sqrt{2\gamma\eps_u}\,\dd B_u\right),
 \qquad \overline Y_t:=Y_t-G_t,
\]
the integral equation for \eqref{eq:model} becomes
\[
 Y_t=y+\int_0^t\mathfrak b(Y_u)\dd u+G_t,
 \qquad
 \dot{\overline Y}_t=\mathfrak b\left(\overline Y_t+G_t\right),
 \quad \overline Y_0=y.
\]
Thus, for every Brownian path, $\overline Y$ solves a nonautonomous ODE with
continuous forcing. Since $D\mathfrak b$ is bounded,
pathwise Gr\"{o}nwall's inequality gives
\begin{equation}\label{eq:pathwise-flow-growth}
 \sup_{0\leq t\leq T}|\Phi_t(y)|
 +\sup_{0\leq t\leq T}\left|\Phi_t^{-1}(y)\right|
 \leq C_T(1+|y|+N_T).
\end{equation}
The first-, second-, and third-order derivatives of the flow with respect to
its initial point,
$J_t=D_y\Phi_t(y)$,
$\mathcal K_t=D_y^2\Phi_t(y)$,
$L_t=D_y^3\Phi_t(y)$,
satisfy
\begin{align*}
 \dot J_t&=D\mathfrak b(Y_t)J_t,\qquad J_0=I,\\
 \dot{\mathcal K}_t[\eta_1,\eta_2]
 &=D\mathfrak b(Y_t)\mathcal K_t[\eta_1,\eta_2]
   +D^2\mathfrak b(Y_t)[J_t\eta_1,J_t\eta_2],\\
 \dot L_t[\eta_1,\eta_2,\eta_3]
 &=D\mathfrak b(Y_t)L_t[\eta_1,\eta_2,\eta_3]
   +D^3\mathfrak b(Y_t)[J_t\eta_1,J_t\eta_2,J_t\eta_3]+\sum_{\rm cyc}D^2\mathfrak b(Y_t)
       [J_t\eta_1,\mathcal K_t[\eta_2,\eta_3]],
\end{align*}
where $\sum_{\rm cyc}$ denotes the sum over the three cyclic permutations of
$(\eta_1,\eta_2,\eta_3)$; explicitly, 
\[
\begin{aligned}
 \sum_{\rm cyc}D^2\mathfrak b(Y_t)
       [J_t\eta_1,\mathcal K_t[\eta_2,\eta_3]]&=D^2\mathfrak b(Y_t)[J_t\eta_1,\mathcal K_t[\eta_2,\eta_3]]
 +D^2\mathfrak b(Y_t)[J_t\eta_2,\mathcal K_t[\eta_3,\eta_1]]\\
 &\qquad
 +D^2\mathfrak b(Y_t)[J_t\eta_3,\mathcal K_t[\eta_1,\eta_2]].
\end{aligned}
\]
with zero initial conditions for $\mathcal K_t$ and $L_t$. The polynomial bounds on
$\nabla^3U$ and $\nabla^4U$, followed successively by Gr\"{o}nwall's inequality,
give polynomial bounds for $\mathcal K_t$ and $L_t$.
More precisely, $\|J_t\|\leq e^{Ct}$ and, for finite exponents $p_2,p_3$,
\begin{align*}
 \|\mathcal K_t\|&\leq C\int_0^t\|\mathcal K_u\|\dd u
 +C_T\int_0^t(1+|Y_u|^{p_2})\|J_u\|^2\dd u,\\
 \|L_t\|&\leq C\int_0^t\|L_u\|\dd u
 +C_T\int_0^t(1+|Y_u|^{p_3})
 \left(\|J_u\|^3+\|J_u\|\|\mathcal K_u\|\right)\dd u.
\end{align*}
The pathwise estimate \eqref{eq:pathwise-flow-growth}
and two successive applications of
Gr\"{o}nwall's inequality therefore yield, for some finite $p_2',p_3'$,
\[
 \sup_{0\leq t\leq T}\|\mathcal K_t\|
 \leq C_T(1+|y|+N_T)^{p_2'},\qquad
 \sup_{0\leq t\leq T}\|L_t\|
 \leq C_T(1+|y|+N_T)^{p_3'}.
\]
Differentiating $\Phi_t^{-1}\circ\Phi_t=\Id$ gives the same bounds for the
inverse flow.
Indeed, if $\Psi_t=\Phi_t^{-1}$, then, with all derivatives of $\Phi_t$
evaluated at $\Psi_t(y)$,
\begin{equation}\label{eq:inverseflowderivatives}
\begin{aligned}
 D\Psi_t(y)&=\left[D\Phi_t(\Psi_t(y))\right]^{-1},\\
 D^2\Psi_t(y)[u,v]
 &=-D\Psi_t(y)D^2\Phi_t(\Psi_t(y))
       [D\Psi_t(y)u,D\Psi_t(y)v],\\
 D^3\Psi_t
 &=-D\Psi_t\left(
 D^3\Phi_t\left[D\Psi_t,D\Psi_t,D\Psi_t\right]
 +\sum_{\rm cyc}D^2\Phi_t\left[D^2\Psi_t,D\Psi_t\right]\right).
\end{aligned}
\end{equation}
The bound on $D\Psi_t$ follows equivalently by solving the first variational
equation backward, and \eqref{eq:inverseflowderivatives}
gives the stated polynomial
bounds for $D^2\Psi_t$ and $D^3\Psi_t$.
Finally, $\operatorname{div}\mathfrak b=-\gamma d$, so the pathwise
Jacobi--Liouville formula \cite{Kunita1990}
and
the change of variables under the random flow give
\begin{equation}\label{eq:flowdensity}
 \det D\Phi_t=e^{-\gamma dt},\qquad
 \rho_t(y)=e^{\gamma dt}\E\left[\rho_0\left(\Phi_t^{-1}(y)\right)\right].
\end{equation}
In particular, for every multi-index
$\alpha$ with $2\leq|\alpha|\leq3$,
and for any $0\leq t\leq T$,
\begin{equation}\label{eq:flowbounds}
\begin{aligned}
 &\left|\Phi_t^{\pm1}(y)\right|\leq C_T(1+|y|+N_T),
 \\
 &\left\|D\Phi_t^{\pm1}(y)\right\|\leq C_T,
 \\
 &\left|\partial^\alpha\Phi_t^{\pm1}(y)\right|
 \leq C_{T,\alpha}(1+|y|+N_T)^{q_\alpha},
\end{aligned}
\end{equation}
for some $q_\alpha<\infty$.

\smallskip
\noindent\emph{Step 1: regular initial densities.}
Assume first that $\rho_0$ is positive, smooth, and Gaussian outside a compact
set. Define
\[
 a_t(y,\omega):=e^{\gamma dt}\rho_0\left(\Phi_t^{-1}(y,\omega)\right),
 \qquad A_{\alpha,t}:=\frac{\partial^\alpha a_t}{a_t}.
\]
By \eqref{eq:flowbounds}, for suitable $m_\alpha<\infty$,
\begin{equation}\label{eq:Aalphabound}
 |A_{\alpha,t}(S_t,\omega)|
 \leq C_{T,\alpha}(1+|S_0|+N_T)^{m_\alpha},\qquad 0\leq t\leq T,
\end{equation}
and the right-hand side of \eqref{eq:Aalphabound}
has moments of every order because the regular
initial density has Gaussian tails. Differentiating
\eqref{eq:flowdensity} under the expectation and conditioning on $S_t=y$ give
\[
 \frac{\partial^\alpha\rho_t}{\rho_t}(y)
   =\E[A_{\alpha,t}(S_t,\omega)\mid S_t=y].
\]
Conditional Jensen's inequality and the polynomial identities relating
$\partial^\alpha\log\rho_t$ to
$(\partial^\beta\rho_t)/\rho_t$, $1\leq|\beta|\leq|\alpha|$, imply
\begin{equation}\label{eq:logdensitybounds}
 \sup_{0\leq t\leq T}\int_{\R^{3d}}(1+|y|^m)
 |\partial^\alpha\log\rho_t(y)|^p\rho_t(y)\dd y<\infty,
\end{equation}
for all finite $m,p$ and $1\leq|\alpha|\leq3$. Moreover,
\begin{equation}\label{eq:densitybounds}
 c_Te^{-C_T(1+|y|^2)}\leq\rho_t(y)
 \leq e^{\gamma dT}\|\rho_0\|_\infty,
 \qquad 0\leq t\leq T.
\end{equation}
Indeed, positivity and the Gaussian tails of the regular initial density give
$\rho_0(x)\geq c e^{-C(1+|x|^2)}$. Choose $R_T<\infty$ with
$p_T:=\mathbb P(N_T\leq R_T)>0$. On this event, \eqref{eq:flowbounds} yields
\[
 \rho_0\left(\Phi_t^{-1}(y)\right)
 \geq c_Te^{-C_T(1+|y|^2+R_T^2)}.
\]
Multiplication by $e^{\gamma dt}p_T$ and absorption of $R_T$ into the constants
prove the lower bound in \eqref{eq:densitybounds}.

Let $\chi_R\in C_c^\infty$ satisfy
\[
 \mathbf 1_{B_R}\leq\chi_R\leq\mathbf 1_{B_{2R}},
 \qquad \left|\partial^\alpha\chi_R\right|\leq C_\alpha R^{-|\alpha|}.
\]
The two bounds in \eqref{eq:densitybounds} imply
$|\log\rho_t(y)|\leq C_T(1+|y|^2)$. For
$w=\nabla\log(\rho/\pi_\eps)$, define the localized functional
\[
 \cH_{\eps,R}(\rho)
 :=\int_{\R^{3d}}\chi_Rw^\top Mw\,\rho\,\dd y
 +\frac{A_0}{\eps}\int_{\R^{3d}}
   \chi_R\log\frac{\rho}{\pi_\eps}\rho\,\dd y.
\]
For the frozen dynamics, differentiation and integration by parts give
\begin{align}\label{eq:localizedfrozen}
 \frac{\dd}{\dd t}\cH_{\eps,R}(\rho_t)
 =&-\int_{\R^{3d}}\chi_R
 w^\top\left(-J_{\nabla^2U}M-MJ_{\nabla^2U}^\top
             +A_0\gamma P_z\right)w\,\dd\nu_t\notag\\
 &\qquad-2\gamma\eps\sum_{i=1}^d\int_{\R^{3d}}\chi_R
  (\partial_{z_i}w)^\top M(\partial_{z_i}w)\,\dd\nu_t
  +\mathcal R_R^{\rm fr}.
\end{align}
Here $\mathcal R_R^{\rm fr}$ contains exactly the terms in which at least one
spatial derivative falls on $\chi_R$. For a fixed law $\nu=\rho\,\dd y$,
direct differentiation with respect to the temperature gives
\begin{align}\label{eq:localizedtemperature}
 \partial_\eps\cH_{\eps,R}(\rho)
 =&-\frac{2}{\eps^2}\int_{\R^{3d}}
       \chi_Rw^\top M\nabla H\,\dd\nu
 -\frac{A_0}{\eps^2}\int_{\R^{3d}}
       \chi_R\log\frac{\rho}{\pi_\eps}\,\dd\nu\notag\\
 &\qquad-\frac{A_0}{\eps^3}\int_{\R^{3d}}
       \chi_R\left(H-\pi_\eps(H)\right)\,\dd\nu.
\end{align}
The bounds \eqref{eq:logdensitybounds}--\eqref{eq:densitybounds} and dominated
convergence theorem show that the terms in \eqref{eq:localizedfrozen} and
\eqref{eq:localizedtemperature} not involving derivatives of $\chi_R$
converge, as $R\to\infty$, to the right-hand sides of
\eqref{eq:frozenidentity} and \eqref{eq:temperatureidentity}, respectively.
Each term in the remainder $\mathcal R_R^{\rm fr}$ in
\eqref{eq:localizedfrozen}
is a finite sum of terms
bounded by expressions of the form
\[
 C R^{-j}\int_{\{R\leq|y|\leq2R\}}
 (1+|y|^q)(1+|\log\rho_t|)
 \prod_{\ell=1}^m\left|\partial^{\beta_\ell}\log\rho_t\right|\,\rho_t\dd y,
\]
where $j\geq1$, $q<\infty$, $0\leq m\leq3$, and
$\sum_{\ell=1}^m|\beta_\ell|\leq3$, with the empty product interpreted as one.
H\"older's inequality and
\eqref{eq:logdensitybounds}, used with a larger polynomial weight and exponent,
make these tails converge to zero uniformly on $[0,T]$. They therefore vanish
in $L^1([0,T])$. The Fokker--Planck equation and the same integrable bounds
justify differentiation in time. Thus \eqref{eq:frozenidentity} and
\eqref{eq:temperatureidentity}, integrated on $[0,T]$, hold for this class of
initial densities.

\smallskip
\noindent\emph{Step 2: approximation of the initial density.}
Let $f=\sqrt{\rho_0}$. Since $\nabla H$ has at most linear growth, finite
relative Fisher information and finite second moment imply
\begin{equation}\label{eq:rawfisherfinite}
 \int_{\R^{3d}}|\nabla f|^2\dd y<\infty;
 \qquad f\in H^1(\R^{3d}).
\end{equation}
Indeed,
\[
\left|\nabla\log\rho_0\right|^2\leq
2\left|\nabla\log(\rho_0/\pi_{\eps_0})\right|^2+2\eps_0^{-2}\left|\nabla H\right|^2;
\]
multiplication by $\rho_0$ and the identity
$|\nabla\rho_0|^2/\rho_0=4|\nabla f|^2$ give \eqref{eq:rawfisherfinite}.
Choose cutoffs $\theta_n$ with $\theta_nf\to f$ in $H^1$ and
$\int_{\R^{3d}}\left(1+|y|^2\right)\left|\theta_n^2f^2-f^2\right|\dd y\to0$; 
the radii can be chosen
successively from the finite $H^1$ and second-moment tails. Convolution with a
nonnegative mollifier, addition of $a_ne^{-|y|^2}$ with $a_n\downarrow0$, and
$L^2$ normalization then produce positive smooth functions $f_n$, Gaussian
outside a compact set, such that
\begin{align}
 f_n\longrightarrow f \qquad\text{in }H^1(\R^{3d}),\label{eq:h1approx}
\end{align}
and
\begin{align}
 \int_{\R^{3d}}\left(1+|y|^2\right)\left|f_n^2-f^2\right|\dd y\longrightarrow0.
 \label{eq:weightedapprox}
\end{align}
Set $\rho_0^{(n)}=f_n^2$ and $\nu_0^{(n)}=\rho_0^{(n)}\dd y$.
Let $m=3d$ and $q=m/(m-2)>1$. Since $m\geq3$, Sobolev embedding and
\eqref{eq:h1approx} give $f_n\to f$ both in $L^2$ and in
$L^{2m/(m-2)}$. Hence,
\begin{align*}
 \left\|\rho_0^{(n)}-\rho_0\right\|_{L^1}
 &\leq \|f_n-f\|_{L^2}(\|f_n\|_{L^2}+\|f\|_{L^2}),\\
 \left\|\rho_0^{(n)}-\rho_0\right\|_{L^q}
 &\leq \|f_n-f\|_{L^{2q}}
          (\|f_n\|_{L^{2q}}+\|f\|_{L^{2q}}),
\end{align*}
and both right-hand sides tend to zero because $2q=2m/(m-2)$.
Thus $\rho_0^{(n)}\to\rho_0$ in $L^1\cap L^{1+\delta}$ with
$\delta=q-1=2/(m-2)$. If $G_+(u):=(u\log u)_+$, then
\[
 |G_+(u)-G_+(v)|
 \leq C_q|u-v|\left(1+u^{q-1}+v^{q-1}\right),\qquad u,v\geq0.
\]
The preceding $L^1$ and $L^q$ convergence and H\"older's inequality therefore
give $G_+\left(\rho_0^{(n)}\right)\to G_+(\rho_0)$ in $L^1$; in particular, the positive
entropy parts are uniformly integrable.

For $G_-(u):=(-u\log u)_+$, use
\begin{equation}\label{eq:negativeentropypointwise}
 u\log(1/u)\leq |y|^2u+e^{-|y|^2-1},
 \qquad 0<u\leq1,
\end{equation}
together with \eqref{eq:weightedapprox}. On every ball,
$0\leq G_-\left(\rho_0^{(n)}\right)\leq e^{-1}$ and convergence in measure follows from
the $L^1$ convergence above. Dominated convergence along almost-everywhere
convergent subsequences handles the integral on the ball. Outside a ball, the
weighted convergence gives
\[
 \lim_{R\to\infty}\sup_{n\geq1}
 \int_{\{|y|>R\}}|y|^2\rho_0^{(n)}(y)\dd y=0.
\]
The pointwise entropy bound \eqref{eq:negativeentropypointwise}
therefore gives a tail bound tending to zero
uniformly in $n$. Consequently, the negative entropy parts are uniformly
integrable and their integrals converge. Since
$|H(y)|\leq C(1+|y|^2)$, \eqref{eq:weightedapprox} also controls the energy
term. We have proved
\begin{equation}\label{eq:entropyapprox}
 \int_{\R^{3d}}\rho_0^{(n)}\log\rho_0^{(n)}\dd y
 \longrightarrow \int_{\R^{3d}}\rho_0\log\rho_0\dd y,
 \qquad
 \int_{\R^{3d}} H\left|\rho_0^{(n)}-\rho_0\right|\dd y\longrightarrow0.
\end{equation}

Define
\[
 \cI_M^{\rm raw}(\rho):=\int_{\R^{3d}}(\nabla\log\rho)^\top
 M(\nabla\log\rho)\rho\dd y.
\]
Equation \eqref{eq:h1approx} gives
\[
 \cI_M^{\rm raw}\left(\rho_0^{(n)}\right)
 =4\int_{\R^{3d}}(\nabla f_n)^\top M\nabla f_n\dd y
 \longrightarrow\cI_M^{\rm raw}(\rho_0).
\]
For any density $\rho$ under consideration,
\begin{align}\label{eq:relativefisherexpand}
 \cI_M(\rho\dd y\mid\pi_\eps)
 =\cI_M^{\rm raw}(\rho)
 +\frac{2}{\eps}\int_{\R^{3d}}(\nabla\rho)^\top M\nabla H\dd y
 +\frac{1}{\eps^2}\int_{\R^{3d}}(\nabla H)^\top M\nabla H\,\rho\dd y,
\end{align}
where cutoff integration by parts yields
\[
 \int_{\R^{3d}}(\nabla\rho)^\top M\nabla H\dd y
 =-\int_{\R^{3d}}\rho\,\operatorname{tr}\left(M\nabla^2H\right)\dd y.
\]
The Hessian of $H$ is bounded and $|\nabla H(y)|\leq C(1+|y|)$.
Consequently, \eqref{eq:weightedapprox}--\eqref{eq:relativefisherexpand}
imply
\begin{equation}\label{eq:functionalapprox}
 \cH_{\eps_0}\left(\nu_0^{(n)}\right)\longrightarrow
 \cH_{\eps_0}(\nu_0).
\end{equation}

\smallskip
\noindent\emph{Step 3: passage to the limit.}
Let $\nu_t^{(n)}$ and $\nu_t$ be the laws starting from
$\nu_0^{(n)}$ and $\nu_0$, respectively. By \eqref{eq:weightedapprox},
$\mathcal W_2\left(\nu_0^{(n)},\nu_0\right)\to0$; synchronous coupling and the Lipschitz drift
then give
\begin{equation}\label{eq:floww2}
 \sup_{0\leq t\leq T}\mathcal W_2\left(\nu_t^{(n)},\nu_t\right)
 \leq C_T\mathcal W_2\left(\nu_0^{(n)},\nu_0\right)\longrightarrow0.
\end{equation}
In particular, \eqref{eq:weightedapprox} gives
\[
 C_2:=\sup_{n\geq1}\int_{\R^{3d}}
 \left(1+|y|^2\right)\dd\nu_0^{(n)}(y)<\infty.
\]
The standard finite-time moment estimate for the globally Lipschitz drift,
uniformly bounded noise amplitude on $[0,T]$, and initial bound $C_2$ yields
\begin{equation}\label{eq:regularizedmoments}
 \sup_{n\geq1}\sup_{0\leq t\leq T}
 \int_{\R^{3d}}\left(1+|y|^2\right)\dd\nu_t^{(n)}(y)\leq C_{T,C_2}.
\end{equation}

For fixed $t$, relative entropy with respect to $\pi_{\eps_t}$ is lower
semicontinuous under the weak convergence implied by \eqref{eq:floww2}. The
same holds for the weighted relative Fisher information. Indeed, its
extended-valued variational representation is
\begin{equation}\label{eq:fisherdual}
 \cI_M(\nu\mid\pi_\eps)
 =\sup_{\phi\in C_c^\infty(\R^{3d};\R^{3d})}
 \int_{\R^{3d}}\left[-2\operatorname{div}(M\phi)
 +\frac{2}{\eps}\phi^\top M\nabla H-\phi^\top M\phi\right]\dd\nu.
\end{equation}
This follows by integrating the term containing $\nabla\log\rho$ by parts and
then completing the square in $\phi$.
Indeed, writing $w=\nabla\log(\dd\nu/\dd\pi_\eps)$, integration by parts gives
\begin{align*}
 &\int_{\R^{3d}}\left[-2\operatorname{div}(M\phi)
 +\frac{2}{\eps}\phi^\top M\nabla H-\phi^\top M\phi\right]\dd\nu\\
 &\qquad=\int_{\R^{3d}}
       \left(2\phi^\top Mw-\phi^\top M\phi\right)\dd\nu\\
 &\qquad=\cI_M(\nu\mid\pi_\eps)
       -\int_{\R^{3d}}(w-\phi)^\top M(w-\phi)\dd\nu
 \leq\cI_M(\nu\mid\pi_\eps).
\end{align*}
Truncating and smoothing $w$ shows that the supremum over
$C_c^\infty$ attains the right-hand side in the extended-valued sense. For
every fixed compactly supported $\phi$, the integrand in
\eqref{eq:fisherdual} is bounded and continuous, because $\nabla H$ is bounded
on the support of $\phi$. Hence, if $\nu_n\Rightarrow\nu$,
\[
 \int_{\R^{3d}} G_\phi\dd\nu=\lim_{n\to\infty}\int_{\R^{3d}} G_\phi\dd\nu_n
 \leq\liminf_{n\to\infty}\cI_M(\nu_n\mid\pi_\eps),
\]
where $G_\phi$ denotes the integrand in \eqref{eq:fisherdual}. Taking the
supremum over $\phi$ proves lower semicontinuity.

By hypothesis, \eqref{eq:abstractcomparison} holds with the same $a,b$ for
every $n$. The integrating-factor formula consequently gives
\[
 \cH_{\eps_t}\left(\nu_t^{(n)}\right)
 \leq e^{-\int_0^t a(u)\dd u}\cH_{\eps_0}\left(\nu_0^{(n)}\right)
 +\int_0^t e^{-\int_v^t a(u)\dd u}b(v)\dd v.
\]
Equation \eqref{eq:functionalapprox} gives convergence of the first term.
Equations \eqref{eq:floww2} and \eqref{eq:fisherdual}, together with lower
semicontinuity of relative entropy, give
\[
 \cH_{\eps_t}(\nu_t)\leq
 \liminf_{n\to\infty}\cH_{\eps_t}\left(\nu_t^{(n)}\right).
\]
Taking the lower limit proves \eqref{eq:abstractintegrated} after undoing the
time shift. On a fixed finite interval, \eqref{eq:regularizedmoments} provides
a moment bound uniform in $n$. Combining \eqref{eq:frozenidentity},
\eqref{eq:temperatureidentity}, and \eqref{eq:coolingexplicit}, and dropping
the nonpositive frozen part, therefore gives
$F_n'\leq C_T(1+F_n)$ uniformly in $n$.
Applying the integrating-factor argument and lower semicontinuity once more
with $a=-C_T$ and $b=C_T$ proves $F(t)<\infty$.
\end{proof}


%%%%%%%%%%%%%%%%%%%%%%%%%%%%%%%%%%%%%%%%%%%%%%%%%%%%%
\subsection{Annealing estimate}

With the coercive matrix and the approximation step in place, we are now ready to state the
continuous-time annealing result in the following theorem.

\begin{theorem}[Continuous-time annealing]\label{thm:main}
Under Assumptions~\ref{ass:potential}, \ref{ass:initial}, and \ref{ass:lsi}, for every
$\delta>0$ and $\alpha>0$ there is $C_{\delta,\alpha}<\infty$ such that, for
all sufficiently large $t$,
\begin{equation}\label{eq:mainrate}
 \mathbb P\left(U(X_t)>\delta\right)
 \leq C_{\delta,\alpha}(t+t_0)^{-\min\{\delta/E,\,(1-D/E)/2\}+\alpha}.
\end{equation}
Here $E$ is the cooling limit in \eqref{eq:coolingassumption}, and $D$ is the
LSI barrier in \eqref{eq:lsi}.
\end{theorem}

\begin{remark}[Interpretation of the continuous-time rate]\label{rem:noacceleration}
With $D$ identified with the critical depth $E_*$, the exponent in
\eqref{eq:mainrate} is the same as the kinetic annealing exponent of He, Tan,
and Wu \cite[Theorem~2.5]{HeTanWu2024}; the extra hypoelliptic level therefore
matches the kinetic polynomial exponent. For fixed $\lambda$ and $\gamma$, the
definitions in \eqref{eq:matrix} give
$b=\mathcal{O}(K^2)$, $\varsigma=\mathcal{O}(K^4)$, and $A_0=\mathcal{O}(K^8)$ as $K\to\infty$.
Thus, Theorem~\ref{thm:main} gives the same continuous-time annealing exponent
as underdamped Langevin dynamics. He, Tan, and Wu
\cite[Remark~2.6]{HeTanWu2024} identify the same barrier-controlled exponent
for overdamped annealing. The quantitative distinction established here is at
the discretization level, while fixed-temperature Eyring--Kramers theory
captures the complementary prefactor improvement.
\end{remark}

\begin{proof}[Proof of Theorem~\ref{thm:main}]
We separate moment control, frozen dissipation, and the cooling contribution.

\emph{Moment control.}
For $a>0$ small, let $b_a=2a/\lambda$, $c_a=a\lambda/\gamma$, and
\[
 V_a=1+H+a x\cdot v+b_a v\cdot z+c_a x\cdot z
                       +\frac12c_a\lambda|x|^2+C_a.
\]
Choose the constant $C_a$ so that $V_a\geq1$. Dissipativity and bounded Hessian
give quadratic lower and upper bounds on $U$.
To make them explicit, set
$R_0:=\max\{1,\sqrt{2m/r}\}$. If $x=\rho\theta$ with $|\theta|=1$ and
$\rho\geq R_0$, then
\[
 \frac{\dd}{\dd\rho}U(\rho\theta)
 =\theta\cdot\nabla U(\rho\theta)
 \geq r\rho-\frac m\rho\geq\frac r2\rho.
\]
After integration from $R_0$ to $\rho$ and adjustment on the ball $B_{R_0}$,
this gives $U(x)\geq r|x|^2/4-C_-$. On the other hand, Taylor's formula and
$\|\nabla^2U\|_{\rm op}\leq K$ give
\[
 U(x)\leq U(0)+|\nabla U(0)||x|+\frac K2|x|^2
 \leq C_+(1+|x|^2).
\]
Thus,
\begin{equation}\label{eq:quadraticU}
 \frac r4|x|^2-C_-\leq U(x)\leq C_+(1+|x|^2).
\end{equation}
Consequently, for small $a$,
$V_a$ is comparable to $1+|x|^2+|v|^2+|z|^2$.
Direct calculation, with cancellation of the $x\cdot v$ and $x\cdot z$ terms,
gives
\begin{align*}
 L_\eps V_a={}&-a x\cdot\nabla U-a|v|^2+(2a-\gamma)|z|^2
 +(c_a-\gamma b_a)v\cdot z-b_a\nabla U\cdot z+\gamma\eps d.
\end{align*}
For clarity, let $c_0=|\nabla U(0)|$ and
$c_1=|\lambda/\gamma-2\gamma/\lambda|$. Young's inequality gives
\begin{align}\label{eq:lyapunovyoung}
 a c_1|v||z|&\leq\frac a4|v|^2+a c_1^2|z|^2,\\
 \frac{2a}{\lambda}(K|x|+c_0)|z|
 &\leq\frac{ar}{4}|x|^2
       +a\left(\frac{4K^2}{r\lambda^2}+1\right)|z|^2
       +\frac{a c_0^2}{\lambda^2}.
\end{align}
In addition to the condition ensuring quadratic comparability of $V_a$, take
$a(3+c_1^2+4K^2/(r\lambda^2))\leq\gamma/2$.
Using $x\cdot\nabla U(x)\geq r|x|^2-m$ together with
\eqref{eq:lyapunovyoung} gives
\[
 L_{\eps_t}V_a
 \leq-\frac{3ar}{4}|x|^2-\frac{3a}{4}|v|^2-\frac{\gamma}{2}|z|^2
       +am+\frac{ac_0^2}{\lambda^2}+\gamma d\sup_{u\geq0}\eps_u
 \leq-cV_a+C.
\]
Dynkin's formula and Gr\"{o}nwall's inequality now yield
\begin{equation}\label{eq:moments}
 \sup_{t\geq0}\E\left[|X_t|^2+|V_t|^2+|Z_t|^2\right]<\infty.
\end{equation}
Moreover $|\nabla_z V_a|^2\leq C V_a$, and hence
\begin{align*}
 L_{\eps_t}V_a^p
 =pV_a^{p-1}L_{\eps_t}V_a
   +\gamma\eps_t p(p-1)V_a^{p-2}|\nabla_zV_a|^2\leq-c_pV_a^p+C_p,
 \qquad p\in\mathbb N.
\end{align*}
In particular, $p=1$ proves the required second-moment bound. Higher moments
are propagated whenever the corresponding initial moment is finite. The
second-moment estimate is uniform for restarted processes with uniformly
bounded initial second moments.
For the Gibbs position marginal, integration by parts gives
$\pi_\eps(x\cdot\nabla U)=d\eps$. Dissipativity therefore bounds its
second moment uniformly for $0<\eps\leq\eps_0$. The quadratic upper bound
on $U$ and the two Gaussian factors then bound $\pi_\eps(H)$ uniformly.
Indeed,
\[
 0=\int_{\R^d}\nabla\cdot\left(xe^{-U(x)/\eps}\right)\dd x
 =dZ_\eps-\eps^{-1}\int_{\R^d}
 x\cdot\nabla U(x)e^{-U(x)/\eps}\dd x,
\]
and hence
\[
 r\mu_\eps(|x|^2)-m\leq d\eps,
 \qquad
 \mu_\eps(|x|^2)\leq\frac{d\eps+m}{r}.
\]
Together with \eqref{eq:quadraticU} and
$\pi_\eps(|v|^2)=\pi_\eps(|z|^2)=d\eps$, this yields
\begin{equation}\label{eq:gibbsmoments}
 \sup_{0<\eps\leq\eps_0}
 \left\{\pi_\eps\left(|x|^2+|v|^2+|z|^2\right)+\pi_\eps(H)\right\}<\infty.
\end{equation}

\emph{Frozen dissipation.}
Recall from \eqref{eq:generator} that
\[
 L_\eps=v\cdot\nabla_x+(-\nabla U+\lambda z)\cdot\nabla_v
 +(-\lambda v-\gamma z)\cdot\nabla_z+\gamma\eps\Delta_z.
\]
Its adjoint in $L^2(\pi_\eps)$
has drift
$b^\dagger=(-v,\nabla U-\lambda z,\lambda v-\gamma z)$, with Jacobian
$J_{\nabla^2U}$. Relative-density differentiation is governed by this adjoint
drift. Let $\nu_t^\eps$ denote the
law evolving with temperature fixed at $\eps$, and let
$h:=\dd\nu_t^\eps/\dd\pi_\eps$, $g:=\log h$, $w:=\nabla g$, and
$\mathcal S:=\gamma\eps\Delta_z-\gamma z\cdot\nabla_z$. Then
\[
 \partial_t g=L_\eps^\dagger g+\gamma\eps|\nabla_z g|^2,
 \qquad L_\eps+L_\eps^\dagger=2\mathcal S.
\]
Differentiate $\int_{\R^{3d}} h w^\top M w\dd\pi_\eps$, move $L_\eps^\dagger$
off $h$, and use
$\nabla(L_\eps^\dagger g)=L_\eps^\dagger w+J_{\nabla^2U}^\top w$.
Integration by parts with $\mathcal S$ cancels the term containing
$\nabla|\nabla_z g|^2$, leaving
\begin{subequations}\label{eq:fisherentropyderivatives}
\begin{align}
 \frac{\dd}{\dd t}\cI_M(\nu_t^\eps\mid\pi_\eps)
 &=\int_{\R^{3d}} w^\top\left(J_{\nabla^2U}M+MJ_{\nabla^2U}^\top\right)w\dd\nu_t^\eps
   -2\gamma\eps\sum_{i=1}^{d}\int_{\R^{3d}}(\partial_{z_i}w)^\top M(\partial_{z_i}w)\dd\nu_t^\eps,
   \label{eq:fisherderivative}\\
 \frac{\dd}{\dd t}\KL(\nu_t^\eps\mid\pi_\eps)
 &=-\gamma\eps\int_{\R^{3d}}|w_z|^2\dd\nu_t^\eps.
   \label{eq:entropyderivative}
\end{align}
\end{subequations}
These operations are justified first for the regular initial laws in
Lemma~\ref{lem:closure}. Multiplying the entropy identity
\eqref{eq:entropyderivative}
by $A_0/\eps$ and adding it to the weighted Fisher-information identity
\eqref{eq:fisherderivative}
gives
\begin{align}\label{eq:frozenidentity}
 \frac{\dd}{\dd t}\cH_\eps(\nu_t^\eps)
 =&-\int_{\R^{3d}} w^\top\left(-J_{\nabla^2U}M-MJ_{\nabla^2U}^\top
                         +A_0\gamma P_z\right)w\,\dd\nu_t^\eps\notag\\
 &\qquad\qquad-2\gamma\eps\sum_{i=1}^d\int_{\R^{3d}}
           (\partial_{z_i}w)^\top M(\partial_{z_i}w)\,\dd\nu_t^\eps.
\end{align}
Lemma~\ref{lem:matrix} implies
\[
 \frac{\dd}{\dd t}\cH_\eps(\nu_t^\eps)
 \leq-\frac12\int_{\R^{3d}}|w|^2\dd\nu_t^\eps.
\]
On the other hand, \eqref{eq:lsi} and tensorization as in
\eqref{eq:tensorizedlsi} give
\[
 \cH_\eps(\nu_t^\eps)
 \leq\left(\|M\|+\frac{A_0C_{\rm LS}(\eps)}{\eps}\right)
       \int_{\R^{3d}}|w|^2\dd\nu_t^\eps.
\]
Consequently,
\begin{equation}\label{eq:decay}
 \frac{\dd}{\dd t}\cH_\eps(\nu_t^\eps)
 \leq-\kappa(\eps)\cH_\eps(\nu_t^\eps),\qquad
 \kappa(\eps):=\frac{1}{2(\|M\|+A_0C_{\rm LS}(\eps)/\eps)}.
\end{equation}
To record the loss in this estimate explicitly, fix $\eta>0$ and first choose
$\eta_{\rm LSI}>0$ sufficiently small such that
$(D+\eta_{\rm LSI})/E<D/E+\eta/3$. Equations
\eqref{eq:lsi} and \eqref{eq:decay} give, for small $\eps$,
\begin{align}\label{exp:as:lower:bound}
 \kappa(\eps)\geq c\eps e^{-(D+\eta_{\rm LSI})/\eps}.
\end{align}
Since $\eps_t\log(t+t_0)\to E$, for all sufficiently large $t$,
\[
 \frac{D+\eta_{\rm LSI}}{\eps_t}
 \leq\left(\frac DE+\frac{2\eta}{3}\right)\log(t+t_0).
\]
Consequently, the exponential in \eqref{exp:as:lower:bound} satisfies
\[
 e^{-(D+\eta_{\rm LSI})/\eps_t}
 \geq(t+t_0)^{-D/E-2\eta/3}.
\]
Finally,
$\eps_t\asymp1/\log(t+t_0)\geq c_\eta(t+t_0)^{-\eta/3}$, and hence
\begin{equation}\label{eq:kappalower}
 \kappa(\eps_t)\geq c_\eta(t+t_0)^{-D/E-\eta}.
\end{equation}

\emph{Cooling contribution.}
For a fixed law $\nu$ with density $\rho$, write
$w:=\nabla\log(\rho/\pi_\eps)$. Then
$\partial_\eps\log(\rho/\pi_\eps)=-(H-\pi_\eps(H))/\eps^2$ and
$\partial_\eps w=-\nabla H/\eps^2$. Therefore,
\begin{align}\label{eq:temperatureidentity}
 \partial_\eps\cH_\eps(\nu)
 ={}&-\frac{2}{\eps^2}\int_{\R^{3d}} w^\top M\nabla H\,\dd\nu
 -\frac{A_0}{\eps^2}\KL(\nu\mid\pi_\eps)
 -\frac{A_0}{\eps^3}\left(\nu(H)-\pi_\eps(H)\right).
\end{align}
For the actual annealing law $\nu_t$, apply \eqref{eq:frozenidentity}
instantaneously at $\eps=\eps_t$ and add $\eps_t'$ times
\eqref{eq:temperatureidentity}. Set $F(t)=\cH_{\eps_t}(\nu_t)$. The moment
bounds \eqref{eq:moments} for $\nu_t$ and \eqref{eq:gibbsmoments} for
$\pi_{\eps_t}$, together with $|\nabla H(y)|\leq C(1+|y|)$ and the
Cauchy--Schwarz inequality, give
\[
 \left|\int_{\R^{3d}} w^\top M\nabla H\dd\nu_t\right|\leq C\sqrt{F(t)},
 \quad \KL(\nu_t\mid\pi_{\eps_t})\leq\eps_t F(t)/A_0,
 \quad |\nu_t(H)-\pi_{\eps_t}(H)|\leq C.
\]
Since $\sqrt F\leq1+F$, the second condition in
\eqref{eq:coolingassumption} bounds the cooling term by
\begin{equation}\label{eq:coolingexplicit}
 C|\eps_t'|\left(1+\eps_t^{-3}\right)(1+F(t))
 \leq C\frac{\log(t+t_0)}{t+t_0}(1+F(t)),
\end{equation}
for all sufficiently large $t$. The constant depends on the initial law
only through its second-moment bound. On every fixed finite interval the
first coefficient in \eqref{eq:coolingexplicit} is bounded. Dropping frozen
dissipation and applying Gr\"{o}nwall's inequality therefore also bounds $F$ before the
low-temperature LSI becomes available, uniformly over the approximations
in Lemma~\ref{lem:closure}.
More explicitly, for each $T<\infty$ there is $C_T<\infty$, depending only
on the common second-moment bound and the model parameters, such that, with
$F_n(t):=\cH_{\eps_t}\left(\nu_t^{(n)}\right)$,
\[
 F_n'(t)\leq C_T(1+F_n(t)),\qquad 0\leq t\leq T.
\]
Thus, for $0\leq s\leq t\leq T$,
\[
 F_n(t)
 \leq e^{C_T(t-s)}\left(F_n(s)+C_T(t-s)\right).
\]
The same $C_T$ applies to every $n$.

For any $0<\eta<1-D/E$, the ratio of
$\log(t+t_0)/(t+t_0)$ to $(t+t_0)^{-D/E-\eta}$ tends to zero.
Absorb the coefficient of $F$ in half of the frozen decay and use
$\log(t+t_0)\leq C_\eta(t+t_0)^\eta$. This gives
\begin{equation}\label{eq:ode}
 F'(t)\leq-c_\eta(t+t_0)^{-D/E-\eta}F(t)
              +C_\eta(t+t_0)^{-1+\eta},
 \qquad F(t)=\cH_{\eps_t}(\nu_t).
\end{equation}
For any $r_0<1-D/E$, choose $\eta$ so that
$r_0<1-D/E-2\eta$. A sufficiently large multiple of
$(t+t_0)^{-r_0}$ is a supersolution of \eqref{eq:ode}. Indeed, if
$\overline F(t)=C(t+t_0)^{-r_0}$, then
\begin{equation}\label{eq:supersolutioncalculation}
 \overline F'(t)+c_\eta(t+t_0)^{-D/E-\eta}\overline F(t)
 ={C}\left[-r_0(t+t_0)^{-r_0-1}
 +c_\eta(t+t_0)^{-r_0-D/E-\eta}\right].
\end{equation}
Because $r_0+D/E+\eta<1-\eta$ and $D/E+\eta<1$, the right-hand side is at
least $C_\eta(t+t_0)^{-1+\eta}$ for all large $t$ after increasing $C$.
Choose $t_1$ sufficiently large and increase $C$ once more so that
$F(t_1)\leq\overline F(t_1)$. For $G=F-\overline F$, subtracting the
supersolution inequality obtained from
\eqref{eq:supersolutioncalculation} from \eqref{eq:ode}
gives
\[
 G'(t)+c_\eta(t+t_0)^{-D/E-\eta}G(t)\leq0,
 \qquad t\geq t_1.
\]
Multiplication by the corresponding integrating factor yields
\[
 G(t)\leq G(t_1)
 \exp\left(-\int_{t_1}^tc_\eta(u+t_0)^{-D/E-\eta}\dd u\right)\leq0.
\]
Therefore
\[
F(t)\leq C_{r_0}(t+t_0)^{-r_0}.
\]

Finally, $\KL(\nu_t\mid\pi_{\eps_t})\leq\eps_tF(t)/A_0$.
Pinsker's inequality gives
\[
 \mathbb P(U(X_t)>\delta)
 \leq\pi_{\eps_t}(U>\delta)
                   +\sqrt{\frac12\KL(\nu_t\mid\pi_{\eps_t})}.
\]
The Gibbs tail estimate used here is
\begin{equation}\label{eq:gibbstail}
 \pi_\eps(U>\delta)\leq C_{\delta,\eta}e^{-(\delta-\eta)/\eps},
 \qquad 0<\eta<\delta.
\end{equation}
To verify \eqref{eq:gibbstail}, choose a measurable set $A_\eta$ of positive
Lebesgue measure on which $U\leq\eta/2$. For $0<\eta<\delta$,
\begin{align}\label{integral:bounds}
 \int_{\{U>\delta\}}e^{-U/\eps}\dd x
 \leq e^{-(\delta-\eta/2)/\eps}
       \int_{\R^d}e^{-\eta U/(2\delta\eps)}\dd x,
\end{align}
and
\begin{align*}
 Z_\eps\geq |A_\eta|e^{-\eta/(2\eps)}.
\end{align*}
The last integral in \eqref{integral:bounds} is uniformly bounded for all sufficiently
small $\eps$ by quadratic confinement and $U\geq0$. Dividing the two bounds
gives \eqref{eq:gibbstail}.
The first condition in \eqref{eq:coolingassumption} converts this exponential
bound into $(t+t_0)^{-(\delta-\eta)/E+o(1)}$. Passing to the limiting initial
law by Lemma~\ref{lem:closure} and taking $r_0$ arbitrarily close to $1-D/E$
proves \eqref{eq:mainrate}.
\end{proof}

The discrete analysis will restart the exact process from the law of a numerical
iterate. The following proposition isolates the continuous estimate with constants
uniform over such restarting laws.

\begin{proposition}[Uniform comparison after a restart]\label{prop:restart}
Fix a bound $M_2$ on the initial second moment and
$0<\eta<1-D/E$. There are $t_*,c_\eta,C_\eta>0$, depending on $M_2$
and the model parameters and uniform in the restarting time and initial
Fisher information, with the following property. Start the exact annealed
process at $s\geq t_*$ from any law satisfying
$\int_{\R^{3d}}|y|^2\nu_s(\dd y)\leq M_2$ and
$F(s)=\cH_{\eps_s}(\nu_s)<\infty$. Then, for
$t\geq s$,
\begin{equation}\label{eq:restartintegral}
 F(t)\leq e^{-\int_s^t a_\eta(u)\dd u}F(s)
 +C_\eta\int_s^t e^{-\int_v^t a_\eta(u)\dd u}
                         (v+t_0)^{-1+\eta}\dd v,
 \quad a_\eta(u):=c_\eta(u+t_0)^{-D/E-\eta}.
\end{equation}
\end{proposition}
\begin{proof}
The restarted version of \eqref{eq:moments} gives a constant
$\overline M_2:=C(1+M_2)$ such that
\[
 \sup_{u\geq s}\int_{\R^{3d}}|y|^2\nu_u(\dd y)\leq\overline M_2.
\]
Hence the constant in \eqref{eq:coolingexplicit} depends only on $M_2$ and
the model parameters. We may therefore choose $t_*$, $c_\eta$, and $C_\eta$
uniformly over the laws in the proposition so that, for every regular initial
density and almost every $u\geq s\geq t_*$,
\begin{equation}\label{eq:restartode}
 F'(u)+a_\eta(u)F(u)
 \leq C_\eta(u+t_0)^{-1+\eta}.
\end{equation}
Set $A_s(t):=\int_s^t a_\eta(u)\dd u$. Multiplication of
\eqref{eq:restartode} by $e^{A_s(u)}$ yields
\[
 \frac{\dd}{\dd u}\left(e^{A_s(u)}F(u)\right)
 \leq C_\eta e^{A_s(u)}(u+t_0)^{-1+\eta}.
\]
Integration over $[s,t]$ and division by $e^{A_s(t)}$ give
\[
 F(t)\leq e^{-A_s(t)}F(s)
 +C_\eta\int_s^t e^{-[A_s(t)-A_s(v)]}
                   (v+t_0)^{-1+\eta}\dd v,
\]
which is \eqref{eq:restartintegral}.

For a general initial law, let $\nu_s^{(n)}$ be the regular approximations
from Lemma~\ref{lem:closure}, chosen so that their second moments are bounded
by $M_2+1$. If
$F_n(u)=\cH_{\eps_u}\left(\nu_u^{(n)}\right)$, then
$F_n(s)\to F(s)$ by Lemma~\ref{lem:closure}, while
\eqref{eq:floww2} and \eqref{eq:fisherdual} give
$F(t)\leq\liminf_{n\to\infty}F_n(t)$. Applying to $F_n$ the
integrating-factor estimate obtained from \eqref{eq:restartode}, namely
\eqref{eq:restartintegral} for regular initial laws,
and passing to the lower limit proves \eqref{eq:restartintegral}. The
constants remain independent of $F(s)$.
\end{proof}

\begin{remark}[Temperature-dependent degeneracy]\label{rem:degeneracy}
The diffusion matrix has rank $d$ in a $3d$-dimensional space even at positive
temperature, and cooling sends its nonzero eigenvalues to zero. The weight
$A_0/\eps$ compensates the factor $\eps$ in entropy dissipation, while
\eqref{eq:decay} and Proposition~\ref{prop:restart} retain the remaining
temperature dependence needed in the discrete comparison.
\end{remark}

%%%%%%%%%%%%%%%%%%%%%%%%%%%%%%%%%%%%%%%%%%%%%%%%%%%%%%%%%
\section{Discrete-Time Annealing}\label{sec:discrete}
\subsection{Discretization schemes and convergence result}

We first define the annealed scheme $(X_{k},V_{k},Z_{k})_{k\geq 0}$ 
based on the numerical scheme of Mou et al. \cite{MouEtAl2021}. 
Write $\tau_0=0$,
$\tau_{k+1}=\tau_k+h_k$, and $s_k=\tau_k+t_0$.
Given $S_k=(X_{k},V_{k},Z_{k})=(x,v,z)$, define $g_u:=\nabla U(x+uv)$ and
$\bar g:=h_k^{-1}\int_0^{h_k}g_u\dd u$. On $0\leq u\leq h_k$ construct
\begin{equation}\label{eq:stages}
\begin{aligned}
 \dd\widehat Z_u&=\left(-\lambda v-\gamma\widehat Z_u\right)\dd u
                         +\sqrt{2\gamma\eps_{\tau_k}}\dd B_u,\\
 \dd\widetilde V_u&=\left(-g_u+\lambda\widehat Z_u\right)\dd u,
 &\dd\widehat V_u&=\left(-\bar g+\lambda\widehat Z_u\right)\dd u,\\
 \dd\widetilde X_u&=\widehat V_u\dd u,
 &\dd\widetilde Z_u&=\left(-\lambda\widehat V_u-\gamma\widetilde Z_u\right)\dd u
                         +\sqrt{2\gamma\eps_{\tau_k}}\dd B_u.
\end{aligned}
\end{equation}
All variables start from the corresponding component of $(x,v,z)$. The two
$Z$ equations use the same Brownian motion. Then
$S_{k+1}=\left(\widetilde X_{h_k},\widetilde V_{h_k},\widetilde Z_{h_k}\right)$ in distribution, 
where $S_{k+1}=\left(X_{k+1},V_{k+1},Z_{k+1}\right)$ is 
generated through the Gaussian transition
\begin{equation}\label{eq:moustep}
 S_{k+1}\mid S_k=s\sim
 \mathcal N\left(m_{h_k}(s),\eps_{\tau_k}\Sigma_{h_k}\right),
\end{equation}
whose mean and covariance we now give explicitly. For a generic step length
$h>0$, set
\begin{align*}
 a_0(h)&:=\frac{1-e^{-\gamma h}}{\gamma},&
 a_1(h)&:=\frac{h}{\gamma}
          -\frac{1-e^{-\gamma h}}{\gamma^2},\\
 a_2(h)&:=\frac{h^2}{2\gamma}-\frac{h}{\gamma^2}
          +\frac{1-e^{-\gamma h}}{\gamma^3},&
 b_0(h)&:=\frac{1-e^{-\gamma h}-\gamma h e^{-\gamma h}}{\gamma^2},\\
 b_1(h)&:=\frac{h(1+e^{-\gamma h})}{\gamma^2}
          -\frac{2(1-e^{-\gamma h})}{\gamma^3}.
\end{align*}
For $s=(x,v,z)$ and
$\Delta U=\int_0^h\nabla U(x+uv)\dd u$, write
$m_h:=\left(m_h^x,m_h^v,m_h^z\right)$. Direct integration of
\eqref{eq:stages} yields
\begin{align}\label{eq:explicitmean}
 m_h^x(s)
 &:=x-\frac h2\Delta U+\left(h-\lambda^2a_2(h)\right)v
      +\lambda a_1(h)z,\notag\\
 m_h^v(s)
 &:=-\Delta U+\left(1-\lambda^2a_1(h)\right)v
      +\lambda a_0(h)z,\notag\\
 m_h^z(s)
 &:=\left(\frac{\lambda}{\gamma}
   -\frac{\lambda(1-e^{-\gamma h})}{\gamma^2h}\right)\Delta U
   +\left(-\lambda a_0(h)+\lambda^3b_1(h)\right)v
      +\left(e^{-\gamma h}-\lambda^2b_0(h)\right)z.
\end{align}
To describe the covariance, define the column vector
\begin{equation}\label{eq:noisekernels}
\mathbf k(r):=
\begin{pmatrix}
 \displaystyle\frac{\lambda}{\gamma}
 \left(r-\frac{1-e^{-\gamma r}}{\gamma}\right)\\[5pt]
 \displaystyle\frac{\lambda}{\gamma}(1-e^{-\gamma r})\\[5pt]
 \displaystyle e^{-\gamma r}
 -\frac{\lambda^2}{\gamma^2}
  \left(1-e^{-\gamma r}-\gamma r e^{-\gamma r}\right)
\end{pmatrix}.
\end{equation}
Equivalently, $\mathbf k(r)=\mathbf u_0+r\mathbf u_1
+e^{-\gamma r}(\mathbf v_0+r\mathbf v_1)$, where
\begin{align*}
 \mathbf u_0:=\left(-\frac{\lambda}{\gamma^2},
                       \frac{\lambda}{\gamma},
                      -\frac{\lambda^2}{\gamma^2}\right)^\top,
 \quad\mathbf u_1:=\left(\frac{\lambda}{\gamma},0,0\right)^\top,
 \quad\mathbf v_0:=\left(\frac{\lambda}{\gamma^2},
                      -\frac{\lambda}{\gamma},
                       1+\frac{\lambda^2}{\gamma^2}\right)^\top,
 \quad\mathbf v_1:=\left(0,0,\frac{\lambda^2}{\gamma}\right)^\top.
\end{align*}
Then the covariance in \eqref{eq:moustep} can be evaluated without numerical
quadrature as follows: 
\begin{equation}\label{eq:explicitcovariance}
 \Sigma_h
 :=2\gamma\int_0^h\mathbf k(r)\mathbf k(r)^\top\dd r
   \otimes\Id_d
 =2\gamma K_h\otimes\Id_d,
\end{equation}
where
\begin{align*}
 K_h:=&h\mathbf u_0\mathbf u_0^\top
 +\frac{h^2}{2}\left(\mathbf u_0\mathbf u_1^\top
                      +\mathbf u_1\mathbf u_0^\top\right)
 +\frac{h^3}{3}\mathbf u_1\mathbf u_1^\top\\
 &+J_0(\gamma,h)\left(\mathbf u_0\mathbf v_0^\top
                      +\mathbf v_0\mathbf u_0^\top\right)+J_1(\gamma,h)\left(\mathbf u_0\mathbf v_1^\top
                      +\mathbf v_1\mathbf u_0^\top
                      +\mathbf u_1\mathbf v_0^\top
                      +\mathbf v_0\mathbf u_1^\top\right)\\
 &\qquad+J_2(\gamma,h)\left(\mathbf u_1\mathbf v_1^\top
                      +\mathbf v_1\mathbf u_1^\top\right)
 +J_0(2\gamma,h)\mathbf v_0\mathbf v_0^\top\\
 &\qquad\qquad+J_1(2\gamma,h)\left(\mathbf v_0\mathbf v_1^\top
                       +\mathbf v_1\mathbf v_0^\top\right)
 +J_2(2\gamma,h)\mathbf v_1\mathbf v_1^\top,
\end{align*}
with
\[
 J_n(q,h):=\frac{n!}{q^{n+1}}
 \left(1-e^{-qh}\sum_{j=0}^n\frac{(qh)^j}{j!}\right),
\]
for $n=0,1,2$ and $q>0$.
The force enters the endpoint only through $\Delta U$. Besides this exact-integral
version, we consider the implementable midpoint version obtained by replacing
$\Delta U$ in that same endpoint mean by
\begin{equation}\label{eq:midpoint}
 \Delta U_{\rm mid}:=h_k\nabla U\left(x+\frac12h_kv\right).
\end{equation}
The covariance is unchanged. This selects the exact-integral and midpoint
members of the integration family of Mou et al. \cite{MouEtAl2021} analyzed 
in the next theorem below.

\begin{theorem}[Discrete annealing for two schemes of Mou et al.]\label{thm:discrete}
Assume Assumptions~\ref{ass:potential}, \ref{ass:lsi}, and \ref{ass:discrete}.
For either \eqref{eq:stages} or its midpoint endpoint version, take
\[
 h_k=c_hs_k^{-a},\qquad c_h>0,\quad a>0,
\]
with $c_h$ small enough that every step is in a fixed sufficiently small
interval $(0,h_*]$. Let $\vartheta=D/E$, where $D$ and $E$ are defined in
\eqref{eq:lsi} and \eqref{eq:coolingassumption}, respectively.
If $a>\vartheta/2$, then for every
$\delta,\alpha>0$ and all large $k$,
\begin{equation}\label{eq:discreterate}
 \mathbb P\left(U(X_k)>\delta\right)
 \leq C_{\delta,\alpha}s_k^{-r(\delta,a)+\alpha},\qquad
 r(\delta,a):=\min\left\{\frac\delta E,
                   \frac{1-\vartheta}{2},2a-\vartheta\right\}.
\end{equation}
In particular, $a\geq(1+\vartheta)/4$ retains the continuous-time exponent.
Moreover $s_k\asymp k^{1/(1+a)}$, so the corresponding iteration exponent is
$r(\delta,a)/(1+a)$, again up to an arbitrarily small loss.
\end{theorem}

\begin{proof}
    We provide the proof in Section~\ref{sec:proof_of_theorem_discrete}.
\end{proof}

The next corollary chooses the coarsest polynomial grid that preserves the
continuous-time rate and records the resulting iteration exponent.

\begin{corollary}[Rate-preserving iteration exponent]
\label{cor:iterationgain}
Under the assumptions of Theorem~\ref{thm:discrete}, let
\[
 r_{\rm c}(\delta):=\min\left\{\frac\delta E,
                    \frac{1-\vartheta}{2}\right\},
 \qquad a_*:=\frac{1+\vartheta}{4}.
\]
For either scheme in that theorem, the choice $a=a_*$ gives, for every
$\delta,\alpha>0$ and all sufficiently large $k$,
\[
 \mathbb P\left(U(X_k)>\delta\right)
 \leq C_{\delta,\alpha}
 k^{-4r_{\rm c}(\delta)/(5+\vartheta)+\alpha}.
\]
\end{corollary}
\begin{proof}
At $a=a_*$, the discretization exponent satisfies
$2a-\vartheta=(1-\vartheta)/2$, so \eqref{eq:discreterate} has exponent
$r_{\rm c}(\delta)$. Since
$s_k\asymp k^{1/(1+a_*)}=k^{4/(5+\vartheta)}$, substitution in
\eqref{eq:discreterate}, followed by a relabeling of the arbitrarily small
loss, proves the claim.
\end{proof}

For comparison, Appendix~\ref{sec:ubuappendix} analyzes the one-gradient
underdamped UBU splitting in the same recent-window framework. Its leading
local stochastic error is $\mathcal{O}(\sqrt\eps\,h^{5/2})$ and yields the
rate-preserving sufficient condition $a\geq(1+\vartheta)/3$ and iteration
exponent $3r_{\rm c}(\delta)/(4+\vartheta)$; see
Theorem~\ref{thm:ubu} and Corollary~\ref{cor:ubuiteration}.

\begin{remark}[Step-size comparison with kinetic annealing]
\label{rem:stepsizecomparison}
Theorem~2.8 of He, Tan, and Wu \cite{HeTanWu2024} for the (kinetic) underdamped Langevin dynamics 
assumes the sufficient grid
condition
\[
 \limsup_{k\to\infty}\Delta t_k\sqrt{T_k}<\infty.
\]
For polynomial steps $\Delta t_k\asymp T_k^{-a}$, this corresponds to
$a\geq1/2$. For UBU, Theorem~\ref{thm:ubu} gives the rate-preserving
sufficient condition $a\geq(1+\vartheta)/3$ and iteration exponent
$3r_{\rm c}(\delta)/(4+\vartheta)$. For the third-order midpoint scheme
analyzed here,
Theorem~\ref{thm:discrete} requires only
$a>\vartheta/2$ for convergence and retains the continuous-time exponent once
\[
 a\geq a_*:=\frac{1+\vartheta}{4}<\frac12.
\]
Indeed, the $\mathcal{O}(h^3)$ local coupling bound for this scheme accumulates over the
recent comparison window as $\sum_{j=m}^{k-1}h_j^3$; see
\eqref{eq:telescoping}.
This produces the term $2a-\vartheta$ in
\eqref{eq:discreterate}. If
\[
 r_{\rm c}(\delta)=\min\left\{\frac\delta E,
                    \frac{1-\vartheta}{2}\right\},
\]
then the choice $a=a_*$ gives iteration exponent
$4r_{\rm c}(\delta)/(5+\vartheta)$, whereas the kinetic sufficient condition
$a=1/2$ gives $2r_{\rm c}(\delta)/3$. The former is larger because
$\vartheta<1$. He, Tan, and Wu \cite[Remark~2.9(iv)]{HeTanWu2024} also record
the sufficient overdamped Euler condition
$\limsup_k\Delta t_kT_k<\infty$, citing an earlier preprint by Tang and Zhou.
The published analysis of Tang, Wu, and Zhou \cite{TangWuZhou2024} uses the
same condition, corresponding to $a\geq1$ and iteration exponent
$r_{\rm c}(\delta)/2$. Thus the presently proved
sufficient-condition comparisons give
\[
 \frac{4r_{\rm c}(\delta)}{5+\vartheta}
 >\max\left\{
       \frac{3r_{\rm c}(\delta)}{4+\vartheta},
       \frac{2r_{\rm c}(\delta)}3
      \right\},
 \qquad
 \min\left\{
       \frac{3r_{\rm c}(\delta)}{4+\vartheta},
       \frac{2r_{\rm c}(\delta)}3
      \right\}
 >\frac{r_{\rm c}(\delta)}2.
\]
The UBU and frozen-force exponents cross at $\vartheta=1/2$: the UBU
exponent is larger for $\vartheta<1/2$ and smaller for $\vartheta>1/2$.
For the numerical landscape in Section~\ref{sec:numerics}, identifying the LSI
barrier with the critical depth gives
$\vartheta=D_{\rm crit}/E=0.250836/0.28\approx0.896$. Hence the UBU
coefficient $3/(4+\vartheta)\approx0.613$ is smaller than the frozen-force
coefficient $2/3$. This crossover separates the regimes favored by the two
sufficient bounds: UBU for $\vartheta<1/2$ and frozen force for
$\vartheta>1/2$. The third-order result improves both discrete sufficient
exponents, while the continuous-time exponents agree as explained in
Remark~\ref{rem:noacceleration}. UBU and the third-order midpoint
scheme each use one interior force evaluation and exact linear Gaussian
subflows. Their different bounds arise because UBU has a centered
$\mathcal{O}(\sqrt\eps\,h^{5/2})$ stochastic local error, while noise reaches the position
component one integration later in the third-order chain and the local endpoint
bound is $\mathcal{O}(h^3)$. The kinetic frozen-force scheme also uses one force evaluation
per step, whereas the cost of the exact-integral scheme depends on the chosen
quadrature.
\end{remark}

Throughout the remainder of this section, constants may depend on the
dimension, the potential, $\lambda,\gamma$, and $c_h$.
Assumption~\ref{ass:discrete} enters through the control estimate and midpoint
quadrature. We also record the grid asymptotics used below. The grid covers an
unbounded time horizon, since finite accumulation would keep $h_k$ bounded away
from zero. Thus $s_k\to\infty$, $h_k\to0$, and $s_{k+1}/s_k\to1$. Applying the
mean value theorem to $s_k^{1+a}$ gives
$s_k^{1+a}\sim(1+a)c_h k$.

\subsection{Local comparison and smoothing}
The first ingredient controls the one-step numerical error and provides the
moments needed to average this error over the numerical law.

For $t\geq0$ and $h>0$, let $Q_{t,h}^{\rm int}$ and $Q_{t,h}^{\rm mid}$
denote the one-step kernels obtained from \eqref{eq:moustep} and
\eqref{eq:midpoint}, respectively, after replacing $(\tau_k,h_k)$ by $(t,h)$.
Let
\begin{align}\label{P:t:t:h}
 P_{t,t+h}f(s):=\E\left[f\left(S_{t+h}^{t,s}\right)\right]
\end{align}
denote the transition operator of the exact time-inhomogeneous process
\eqref{eq:model} started from $s$ at time $t$.

\begin{lemma}[Local coupling and moments]\label{lem:local}
Let $Q_{t,h}\in\left\{Q_{t,h}^{\rm int},Q_{t,h}^{\rm mid}\right\}$, and let
$P_{t,t+h}$ be the exact transition in \eqref{P:t:t:h}. Couple the two using the same
Brownian path and start from $s\in\R^{3d}$. For $h\leq h_*$ and large $t$,
\begin{equation}\label{eq:local}
 \left\|S_h^{\rm num}-S_{t+h}^{\rm exact}\right\|_{L^2(\mathbb P_s)}
 \leq C h^3\left(1+|s|^2\right)
       +\frac{C h^{3/2}}{(t+t_0)\log^{3/2}(t+t_0)}.
\end{equation}
For the exact-integral scheme the first factor $1+|s|^2$ can be replaced by
$1+|s|$. Both schemes satisfy $\sup_{k\geq0}\E|S_k|^p<\infty$ for
$1\leq p\leq4$.
\end{lemma}

\begin{proof}
Throughout this proof,
$\|Y\|_{L^p(\mathbb P_s)}:=(\E_s|Y|^p)^{1/p}$ denotes the norm under the
synchronous coupling conditional on the initial state at $s$.
First freeze the exact temperature at $\eps_t$. For the exact-integral
interpolation \eqref{eq:stages}, elementary integrations, the Lipschitz bound
on $\nabla U$, and the Gaussian moment bounds for $\widehat Z$ give, uniformly
for $0\leq u\leq h$ and any fixed $p\geq2$,
\[
 \widehat Z_u=e^{-\gamma u}z
 -\frac{\lambda(1-e^{-\gamma u})}{\gamma}v
 +\sqrt{2\gamma\eps_t}\int_0^u e^{-\gamma(u-r)}\dd B_r,
\]
and hence
\[
\sup_{u\leq h}\left\|\widehat Z_u\right\|_{L^p(\mathbb P_s)}\leq C(1+|s|).
\]
Moreover,
\[
 |g_u|\leq |\nabla U(0)|+K(|x|+h|v|),\qquad
 \widehat V_u-v=-u\bar g+\lambda\int_0^u\widehat Z_r\dd r.
\]
These identities give
\begin{align*}
 &\left\|\widehat V_u-v\right\|_{L^p(\mathbb P_s)}\leq Ch(1+|s|),\\
 &\left\|\widetilde X_u-x-uv\right\|_{L^p(\mathbb P_s)}
 +\left\|\widehat Z_u-\widetilde Z_u\right\|_{L^p(\mathbb P_s)}
 \leq Ch^2(1+|s|),\\
 &\left|\widehat V_u-\widetilde V_u\right|
 =\left|\int_0^ug_r\dd r-u\bar g\right|\leq CKh^2|v|,
\end{align*}
where the last inequality is obtained by using
$|g_r-\bar g|\leq Kh|v|$ and integrate over $[0,u]$, 
and the second line above is established 
by using
\[
\widetilde X_u-x-uv=\int_0^u\left(\widehat V_r-v\right)\dd r, 
\]
and
\[
\widehat Z_u-\widetilde Z_u
=\lambda\int_0^u e^{-\gamma(u-r)}\left(\widehat V_r-v\right)\dd r.
\]
Relative to the exact drift at
$\widetilde S_u=\left(\widetilde X_u,\widetilde V_u,\widetilde Z_u\right)$, the three
interpolation residuals are
\begin{equation}\label{eq:residuals}
 \begin{pmatrix}
 \widehat V_u-\widetilde V_u\\
 \nabla U(\widetilde X_u)-g_u+\lambda(\widehat Z_u-\widetilde Z_u)\\
 -\lambda(\widehat V_u-\widetilde V_u)
 \end{pmatrix}.
\end{equation}
Their $L^p$ norm is at most $Ch^2(1+|s|)$. The noises cancel under synchronous
coupling. Let $S_u^{\rm frozen}=\left(X_u^{\rm frozen},V_u^{\rm frozen},
Z_u^{\rm frozen}\right)$ be the solution, started from $s$ and driven by the same
Brownian path, of
\begin{equation}\label{eq:frozenprocess}
\begin{aligned}
 \dd X_u^{\rm frozen}&=V_u^{\rm frozen}\dd u,\\
 \dd V_u^{\rm frozen}&=\left(-\nabla U\left(X_u^{\rm frozen}\right)
                      +\lambda Z_u^{\rm frozen}\right)\dd u,\\
 \dd Z_u^{\rm frozen}&=\left(-\lambda V_u^{\rm frozen}
                      -\gamma Z_u^{\rm frozen}\right)\dd u
                      +\sqrt{2\gamma\eps_t}\dd B_u.
\end{aligned}
\end{equation}
If $E_u=\widetilde S_u-S_u^{\rm frozen}$, the
Lipschitz drift and \eqref{eq:residuals} imply
\[
 \|E_u\|_{L^p(\mathbb P_s)}
 \leq C\int_0^u\|E_r\|_{L^p(\mathbb P_s)}\dd r
      +C\int_0^u h^2(1+|s|)\dd r,
 \qquad 0\leq u\leq h.
\]
Gr\"{o}nwall's inequality therefore gives
\[
\|E_h\|_{L^p(\mathbb P_s)}\leq Ch^3(1+|s|).
\]

For midpoint quadrature,
\[
 \left|\Delta U_{\rm mid}-\Delta U\right|
 \leq \frac{\|\nabla^3U\|_\infty}{24}h^3|v|^2.
\]
The coefficients multiplying this difference in the $(X,V,Z)$ endpoint mean
are, respectively, 
\begin{equation}\label{eq:forcecoefficients}
 -\frac h2,\qquad -1,\qquad
 \frac\lambda\gamma-\frac{\lambda(1-e^{-\gamma h})}{\gamma^2h}.
\end{equation}
Indeed, writing $q=\Delta U_{\rm mid}-\Delta U$, the two force treatments
give $\delta\widehat V_u=-uq/h$. Consequently,
\begin{align}\label{provide:1}
 \delta V_h=-q,\qquad
 \delta X_h=\int_0^h\delta\widehat V_u\dd u=-\frac h2q,
\end{align}
and
\begin{align}\label{provide:2}
 \delta Z_h
 =-\lambda\int_0^he^{-\gamma(h-u)}\delta\widehat V_u\dd u
 =\frac{\lambda q}{h}\int_0^hu e^{-\gamma(h-u)}\dd u
 =\left(\frac\lambda\gamma
 -\frac{\lambda(1-e^{-\gamma h})}{\gamma^2h}\right)q,
\end{align}
Therefore, \eqref{provide:1} and \eqref{provide:2} provide the coefficients in \eqref{eq:forcecoefficients}.
They are bounded for $h\leq h_*$. Together with the exact-integral estimate,
this proves the first term in \eqref{eq:local} for the midpoint scheme.
We will also use the
coarser bound $C K h^2|v|$, which follows using only the Lipschitz gradient.

The actual temperature changes within a step. By
\eqref{eq:coolingassumption}, $\eps_t\asymp1/\log(t+t_0)$ and
\[
 |(\sqrt{\eps_t})'|
 \leq C\eps_t^{5/2}\frac{\log(t+t_0)}{t+t_0}
 \leq\frac{C}{(t+t_0)\log^{3/2}(t+t_0)}.
\]
Hence the difference between the frozen and exact noise coefficients is
\[
 \sigma_{t+u}-\sigma_t
 =\sqrt{2\gamma}\left(\sqrt{\eps_{t+u}}-\sqrt{\eps_t}\right),
\]
where
\[
 \left|\sqrt{\eps_{t+u}}-\sqrt{\eps_t}\right|
 \leq\int_t^{t+u}|(\sqrt{\eps_r})'|\dd r
 \leq\frac{Cu}{(t+t_0)\log^{3/2}(t+t_0)}.
\]
The It\^o isometry gives
\[
 \left\|\int_0^h
   \left(\sqrt{\eps_{t+u}}-\sqrt{\eps_t}\right)\dd B_u\right\|_{L^2}
 \leq \frac{C h^{3/2}}{(t+t_0)\log^{3/2}(t+t_0)}.
\]
The Lipschitz drift bound and Gr\"{o}nwall's inequality propagate this perturbation
to the full state and give the second term of \eqref{eq:local}.

We next establish the uniform moment bounds
$\sup_{k\geq0}\E|S_k|^p<\infty$, $1\leq p\leq4$, for both numerical schemes.
The Lyapunov function $V_a$ used in the continuous proof can here be denoted
$\mathcal V$ to avoid confusion with the step exponent. It satisfies
\[
L_\eps\mathcal V^p\leq-c_p\mathcal V^p+C_p, 
\]
uniformly for
$0<\eps\leq\eps_0$, and
\[
|\nabla\mathcal V^p(s)|\leq C_p\left(1+|s|^{2p-1}\right).
\]
The frozen exact process \eqref{eq:frozenprocess} consequently satisfies
\[
 \E_s\left[\mathcal V^p\left(S_h^{\rm frozen}\right)\right]
 \leq e^{-c_ph}\mathcal V^p(s)+C_ph.
\]
The residual argument in \eqref{eq:residuals} was stated for every fixed
$L^q$, so the exact-integral endpoint satisfies
\[
 \left\|S_h^{\rm int}-S_h^{\rm frozen}\right\|_{L^q(\mathbb P_s)}
 \leq C_qh^3(1+|s|),\qquad q\geq2.
\]
For moment stability of the midpoint endpoint we use the coarser estimate
\[
 |\Delta U_{\rm mid}-\Delta U|
 \leq K|v|\int_0^h|u-h/2|\dd u
 \leq \frac14Kh^2|v|.
\]
The three endpoint coefficients multiplying this force increment are bounded
by $Ch$, $1$, and $Ch$, respectively, for small $h$.
Indeed, by \eqref{eq:forcecoefficients},
$\left|h/2\right|\leq Ch$, 
while Taylor's formula gives, uniformly for $0<h\leq h_*$,
\[
 \left|\frac\lambda\gamma
 -\frac{\lambda(1-e^{-\gamma h})}{\gamma^2h}\right|
 =\frac\lambda\gamma
 \left|1-\frac{1-e^{-\gamma h}}{\gamma h}\right|
 \leq Ch.
\]
Hence, by the triangle inequality,
\begin{equation}\label{eq:coarsemomentlocal}
 \left\|S_h^{\rm num}-S_h^{\rm frozen}\right\|_{L^q(\mathbb P_s)}
 \leq C_qh^2(1+|s|),
\end{equation}
for either scheme. Their conditional means have linear growth in $s$ and their
Gaussian covariances are uniformly bounded for $h\leq h_*$; consequently both
endpoints have conditional $q$-moment bounded by $C_q(1+|s|^q)$.
We use these statements with $q=2p$ below.
The mean value inequality gives
\[
 \left|\mathcal V^p(y)-\mathcal V^p(y')\right|
 \leq C_p\left(1+|y|^{2p-1}+|y'|^{2p-1}\right)|y-y'|.
\]
Applying H\"older's inequality to the frozen endpoint coupling therefore yields
\[
 \left|Q_{t,h}\mathcal V^p(s)
       -\E_s\left[\mathcal V^p\left(S_h^{\rm frozen}\right)\right]\right|
 \leq C_ph^2\left(1+|s|^{2p}\right)
 \leq C_ph^2\mathcal V^p(s).
\]
Consequently,
\[
 Q_{t,h}\mathcal V^p(s)
 \leq \left(e^{-c_ph}+C_ph^2\right)\mathcal V^p(s)+C_ph.
\]
Since $h_k\to0$, the coefficient is at most $1-c_ph_k/2$ for all sufficiently
large $k$. Thus, for $p=1,2$,
\begin{equation}\label{eq:discrete-lyapunov-recursion}
\E\left[\mathcal V^p(S_{k+1})\right]
 \leq\left(1-\frac{c_p}{2}h_k\right)\E\left[\mathcal V^p(S_k)\right]+C_ph_k.
\end{equation}
The initial fourth moment, the linear growth of the numerical mean, and the
Gaussian noise give, for $p=1,2$,
\[
 \E\left[\mathcal V^p(S_{k+1})\mid S_k=s\right]
 \leq C_p(1+\mathcal V^p(s)).
\]
The conditional moment bound and the finite fourth moment of the initial law
imply by induction that $Y_k:=\E[\mathcal V^p(S_k)]$ is finite for every $k$.
For all sufficiently large $k$, \eqref{eq:discrete-lyapunov-recursion}
also shows that $Y_k$ decreases whenever $Y_k\geq2C_p/c_p$.
After including the finite initial segment, this
gives $\sup_{k\geq0}Y_k<\infty$ for $p=1,2$.
Since
$\mathcal V(s)\asymp1+|s|^2$, the case $p=2$ yields
$\sup_{k\geq0}\E|S_k|^4<\infty$,
and, for $1\leq q\leq4$,
\[
 \E|S_k|^q\leq\left(\E|S_k|^4\right)^{q/4}.
\]
This proves all the remaining moment bounds.
\end{proof}

Grid-point initialization of the distorted entropy additionally requires
regularity of the numerical law. The Gaussian covariance of one step provides
this regularity and a quantitative entropy bound.

\begin{lemma}[Gaussian smoothing at numerical grid points]\label{lem:gaussian}
For the exact-integral scheme \eqref{eq:stages}--\eqref{eq:moustep} and its
midpoint version \eqref{eq:midpoint}, the numerical law at every positive grid
point has a smooth
positive density. For $h\leq h_*$,
\begin{equation}\label{eq:covbound}
 \lambda_{\min}(\Sigma_h)\geq c h^5,\qquad
 \det\Sigma_h\geq c h^{9d}.
\end{equation}
Consequently, for all sufficiently large $k$,
\begin{equation}\label{eq:gridfisher}
 \cH_{\eps_{\tau_k}}(\mathcal L(S_k))
              \leq C s_k^{5a}(\log s_k)^2.
\end{equation}
\end{lemma}
\begin{proof}
The endpoint noise kernels and their covariance are given explicitly in
\eqref{eq:noisekernels}--\eqref{eq:explicitcovariance}; they follow by
integrating successively the $\widehat Z$, $\widehat V$, $\widetilde X$, and
$\widetilde Z$ equations with the same Brownian driver.
The kernels in \eqref{eq:noisekernels} have leading terms
$(\lambda r^2/2,\lambda r,1)$, with
remainders $\mathcal{O}(r^3),\mathcal{O}(r^2),\mathcal{O}(r)$, respectively.
Let $D_h=\operatorname{diag}\left(h^{5/2},h^{3/2},h^{1/2}\right)\otimes\Id_d$.
Directly integrating these kernels gives, as $h\downarrow0$,
\[
 D_h^{-1}\Sigma_hD_h^{-1}\longrightarrow
 2\gamma\begin{pmatrix}
 \lambda^2/20&\lambda^2/8&\lambda/6\\
 \lambda^2/8&\lambda^2/3&\lambda/2\\
 \lambda/6&\lambda/2&1
 \end{pmatrix}\otimes\Id_d.
\]
The scalar matrix is the Gram matrix of the linearly independent polynomials
$(\lambda u^2/2,\lambda u,1)$ in $L^2(0,1)$, hence is positive definite.
The scaled covariance is therefore bounded below by a fixed positive
multiple of the identity for all sufficiently small $h$. Consequently
$\Sigma_h\succeq c\operatorname{diag}(h^5,h^3,h)\otimes\Id_d$;
taking $h_*\leq1$ and taking determinants proves \eqref{eq:covbound}.
The covariance is identical for midpoint quadrature. In particular, the
actual transition covariance has lower bound $c\eps_t h^5\Id_{3d}$, which records
the vanishing-temperature scale explicitly.
Indeed, the midpoint replacement changes only the conditional mean, so both
schemes have transition covariance
$C_{t,h}=\eps_t\Sigma_h$.
Since $h\leq1$,
\[
 \Sigma_h\succeq c\operatorname{diag}\left(h^5,h^3,h\right)\otimes\Id_d
 \succeq ch^5\Id_{3d},
\]
and therefore
\[
 C_{t,h}\succeq c\eps_t h^5\Id_{3d},\qquad
 \det C_{t,h}=\eps_t^{3d}\det\Sigma_h>0.
\]

A mixture of translates of a Gaussian with covariance $C_k$ has raw Fisher
information at most $\operatorname{tr}C_k^{-1}$: its score is the conditional
expectation of the Gaussian score, so Jensen's inequality applies.
More precisely, write $Y=A+\xi$, where, conditionally on the random mean $A$,
$\xi\sim\mathcal N(0,C_k)$. If $\mu_A$ is the law of $A$ and
$\varphi_{C_k}$ is the density of $\mathcal N(0,C_k)$, then
\begin{equation}\label{eq:mixturedensity}
 \rho_Y(y)=\int_{\R^{3d}}\varphi_{C_k}(y-a)\,\mu_A(\dd a).
\end{equation}
Differentiation under the mixture integral \eqref{eq:mixturedensity}
gives
\[
 \nabla\log\rho_Y(y)
 =-\E\left[C_k^{-1}(Y-A)\mid Y=y\right].
\]
Conditional Jensen's inequality then yields
\begin{align*}
 \cI^{\rm raw}(\rho_Y)
 =\E\left|\E\left[C_k^{-1}\xi\mid Y\right]\right|^2
 \leq\E\left|C_k^{-1}\xi\right|^2
 =\operatorname{tr}\left(C_k^{-1}\right).
\end{align*}
Here
$C_k=\eps_{\tau_{k-1}}\Sigma_{h_{k-1}}$. Thus
\[
 \cI\left(\nu_k\mid\pi_{\eps_{\tau_k}}\right)
 \leq \frac{C}{\eps_{\tau_{k-1}}h_{k-1}^5}
       +\frac{C}{\eps_{\tau_k}^2}\E\left(1+|S_k|^2\right).
\]
Convexity of $\Psi(u):=u\log u$ (equivalently, concavity of differential
entropy; see \cite[Section~8.6]{CoverThomas2006}) bounds the negative
differential entropy of the mixture by that of the Gaussian component.
Indeed, Jensen's inequality, Fubini's theorem, and translation invariance give
\begin{align*}
 \int_{\R^{3d}}\Psi(\rho_k(y))\,\dd y
 &\leq\int_{\R^{3d}}\int_{\R^{3d}}
       \Psi\!\left(\varphi_{C_k}(y-a)\right)\,\mu_A(\dd a)\dd y\\
 &=\int_{\R^{3d}}\varphi_{C_k}(z)\log\varphi_{C_k}(z)\,\dd z\\
 &=-\frac{3d}{2}\log(2\pi e)-\frac12\log\det C_k.
\end{align*}
For any $\sigma>0$, nonnegativity of the relative entropy with respect to
$\mathcal N(0,\sigma^2\Id_{3d})$ also gives
\[
 \int_{\R^{3d}}\rho_k(y)\log\rho_k(y)\,\dd y
 \geq-\frac{3d}{2}\log\left(2\pi\sigma^2\right)
      -\frac{1}{2\sigma^2}\E|S_k|^2> -\infty.
\]
Moreover, $\log\mathcal Z_\eps$ is bounded above for $\eps\leq\eps_0$, because
$H\geq0$ and $\mathcal Z_{\eps_0}<\infty$.
Since, for $\nu=\rho\,\dd y$,
\[
 \KL(\nu\mid\pi_\eps)
 =\int_{\R^{3d}}\rho(y)\log\rho(y)\,\dd y
  +\frac{\nu(H)}{\eps}+\log\mathcal Z_\eps,
\]
together with \eqref{eq:covbound}, Lemma~\ref{lem:local}, and
$h_{k-1}\asymp s_k^{-a}$, it yields that
\begin{align*}
 -\log\det C_k&\leq C\log s_k,
 &\frac{\nu_k(H)}{\eps_{\tau_k}}&\leq C\log s_k,\\
 \KL\left(\nu_k\mid\pi_{\eps_{\tau_k}}\right)&\leq C\log s_k,
 &\cI\left(\nu_k\mid\pi_{\eps_{\tau_k}}\right)
   &\leq Cs_k^{5a}(\log s_k)^2.
\end{align*}
Since $\cI_M\leq\|M\|\cI$, substitution in \eqref{eq:functional} proves
\eqref{eq:gridfisher}. Smoothness and positivity
follow from convolution with the nondegenerate Gaussian, while the required
second and fourth moments follow from Lemma~\ref{lem:local}.
\end{proof}

To propagate the local endpoint error to the terminal optimization event, we next
establish a smoothing estimate in the following lemma for the exact inhomogeneous transition.

\begin{lemma}[One-unit smoothing by a finite controlled coupling]\label{lem:smoothing}
Assume $\|\nabla^3U\|_\infty<\infty$. For the exact inhomogeneous transition
operator, any bounded measurable $f$, and any $s,s'\in\R^{3d}$,
\begin{equation}\label{eq:smoothing}
 |P_{t,t+1}f(s)-P_{t,t+1}f(s')|
 \leq C\eps_{t+1}^{-1/2}(1+|s|+|s'|)|s-s'|\|f\|_\infty.
\end{equation}
For $T\geq t+1$,
\[
 |P_{t,T}f(s)-P_{t,T}f(s')|
 \leq C\eps_{t+1}^{-1/2}(1+|s|+|s'|)|s-s'|\|f\|_\infty.
\]
For $0\leq T-t\leq1$ and globally Lipschitz $f$,
\[
\operatorname{Lip}(P_{t,T}f)\leq e^C\operatorname{Lip}(f).
\]
The constants are independent of $t,T,s,s'$.
\end{lemma}

\begin{proof}
We prove a finite-difference estimate directly at the level of bounded
measurable test functions. On a unit interval let $(X_u,V_u,Z_u)$ start from
$s=(x,v,z)$, and write $e=s'-s=(e_x,e_v,e_z)$. Let $q$ be the vector
polynomial of degree at most five determined by
\begin{align*}
 q(0)&=e_x,&q'(0)&=e_v,&
 q''(0)&=-\nabla U(x+e_x)+\nabla U(x)+\lambda e_z,\\
 q(1)&=0,&q'(1)&=0,&q''(1)&=0.
\end{align*}
Finite-dimensional interpolation and the bounded Hessian imply
$\max_{0\leq j\leq3}\left\|q^{(j)}\right\|_\infty\leq C|e|$.
Define the adapted shifted path
\[
 X_u^e:=X_u+q(u),\qquad V_u^e:=V_u+q'(u),\qquad
 Z_u^e:=Z_u+\lambda^{-1}
       \left(q''(u)+\nabla U(X_u+q(u))-\nabla U(X_u)\right).
\]
It starts from $s'$ and has exactly the same endpoint as the unshifted path.
Its first two equations agree with \eqref{eq:model}; its last equation has
an additional drift
\begin{align}\label{eq:control}
 c_e(u)=&\lambda^{-1}\left(q'''
       +\left[\nabla^2U(X_u+q)-\nabla^2U(X_u)\right]V_u
       +\nabla^2U(X_u+q)q'\right)\notag\\
 &+\lambda q'+\frac\gamma\lambda
          \left(q''+\nabla U(X_u+q)-\nabla U(X_u)\right).
\end{align}
Since $X_u$ has finite variation, differentiating $\nabla U(X_u)$ produces
only its ordinary chain-rule term. The bound
$\|\nabla^3U\|_\infty<\infty$
gives
\begin{equation}\label{eq:finitecontrolbound}
 |c_e(u)|\leq C|e|(1+|V_u|).
\end{equation}

We check the change of measure locally in $e$, with an explicit temperature
dependence. Let
$N_u:=\int_0^u\sqrt{2\gamma\eps_{t+r}}\dd B_r$. Pathwise Gr\"{o}nwall's inequality and
the uniform upper bound on the temperature give
\[
 \sup_{u\leq1}|S_{t+u}|\leq C\left(1+|s|+\sup_{u\leq1}|N_u|\right),
\]
and
\[
 \E_s \left[e^{c\sup_{u\leq1}|S_{t+u}|^2}\right]\leq C e^{C|s|^2},
\]
for some $c>0$ independent of $t,s$. The second estimate follows from
the Gaussian tail of the maximum of a Brownian motion with bounded
deterministic time change, coordinate by coordinate.
Let $b_e(u):=c_e(u)/\sqrt{2\gamma\eps_{t+u}}$. By
\eqref{eq:finitecontrolbound}, Novikov's condition holds whenever
$|e|\leq c_*\sqrt{\eps_{t+1}}$, with $c_*>0$ independent of $s,t$.
Consequently,
\[
 R_e:=\exp\left(-\int_0^1 b_e\cdot\dd B
                         -\frac12\int_0^1|b_e|^2\dd u\right)
\]
has expectation one and defines an equivalent measure
$\dd\mathbb Q_e=R_e\dd\mathbb P$. Under $\mathbb Q_e$,
$B_u+\int_0^u b_e(r)\dd r$ is a standard $d$-dimensional Brownian motion, so that the shifted path solves the
uncontrolled SDE starting at $s'$. The drift
\[
 \mathfrak b(x,v,z)=(v,-\nabla U(x)+\lambda z,-\lambda v-\gamma z)
\]
is globally Lipschitz because $\nabla^2U$ is bounded; this ensures pathwise
uniqueness and hence uniqueness in law.
Endpoint equality therefore yields, for every bounded
measurable $f$,
\[
 P_{t,t+1}f(s')=\E_s[R_e f(S_{t+1})].
\]
Use the relative entropy in the direction $\mathbb P\mid\mathbb Q_e$:
the stochastic integral in $-\log R_e$ has zero mean, and hence
\begin{equation}\label{eq:pathentropy}
 \KL(\mathbb P\mid\mathbb Q_e)
 =\frac12\E_s\left[\int_0^1|b_e(u)|^2\dd u\right]
 \leq\frac{C|e|^2}{\eps_{t+1}}(1+|s|^2).
\end{equation}
Pinsker's inequality on the underlying probability space gives
\[
 \left|P_{t,t+1}f(s')-P_{t,t+1}f(s)\right|
 \leq\|f\|_\infty\sqrt{2\KL(\mathbb P\mid\mathbb Q_e)}
 \leq C\eps_{t+1}^{-1/2}|e|(1+|s|)\|f\|_\infty
\]
for these small increments. For arbitrary $s,s'$, partition their straight
line segment into increments of length at most
$c_*\sqrt{\eps_{t+1}}$. Apply the local estimate at each point and sum;
all intermediate norms are at most $|s|+|s'|$.
More precisely, take
\[
 N=\max\left\{1,
 \left\lceil\frac{|s'-s|}{c_*\sqrt{\eps_{t+1}}}\right\rceil\right\},
 \qquad s_i=s+\frac{i}{N}(s'-s),\quad 0\leq i\leq N.
\]
Then $|s_{i+1}-s_i|\leq c_*\sqrt{\eps_{t+1}}$,
$|s_i|\leq |s|+|s'|$, and
\begin{align*}
 \left|P_{t,t+1}f(s')-P_{t,t+1}f(s)\right|
 &\leq\sum_{i=0}^{N-1}
   \left|P_{t,t+1}f(s_{i+1})-P_{t,t+1}f(s_i)\right|\\
 &\leq C\eps_{t+1}^{-1/2}(1+|s|+|s'|)
       \sum_{i=0}^{N-1}|s_{i+1}-s_i|\|f\|_\infty\\
 &=C\eps_{t+1}^{-1/2}(1+|s|+|s'|)|s'-s|\|f\|_\infty.
\end{align*}
This proves
\eqref{eq:smoothing} globally for all $t$, with the factor
$\eps_{t+1}^{-1/2}$ quantifying the temperature dependence.

For $T\geq t+1$ apply \eqref{eq:smoothing} to
$P_{t+1,T}f$, whose supremum norm is at most $\|f\|_\infty$.
For $T-t\leq1$, synchronous coupling and the uniform Lipschitz drift bound
give 
\[
\left|S_T^s-S_T^{s'}\right|\leq e^C\left|s-s'\right|. 
\]
This proves the last assertion.
\end{proof}

We now combine the three preceding lemmas with the uniform restart estimate to complete
the proof of Theorem~\ref{thm:discrete}.

\subsection{Proof of Theorem~\ref{thm:discrete}}\label{sec:proof_of_theorem_discrete}
\begin{proof}[Proof of Theorem~\ref{thm:discrete}]
Fix $T=\tau_k$, and choose a number
$\beta\in(\vartheta,\min\{1,2a\})$, to be taken arbitrarily close to
$\vartheta$. This interval is nonempty because $a>\vartheta/2$.
Choose $m$ so that $\tau_m\leq T-T^\beta<\tau_{m+1}$. Then
$T-\tau_m\asymp T^\beta$ and $s_j\asymp T$ for $m\leq j\leq k$.
Restart the exact annealed process at time $\tau_m$ with law $\nu_m=\mathcal L(S_m)$.
Its initial distorted entropy is polynomial in $T$ by
Lemma~\ref{lem:gaussian}, and its moments are uniformly bounded by
Lemma~\ref{lem:local}. Proposition~\ref{prop:restart} therefore applies
with a common $M_2$ and common constants for every $m$. It gives, for small
$\eta>0$ with $\vartheta+\eta<\beta$,
\begin{equation}\label{eq:restart}
 \cH_{\eps_T}(\nu_mP_{\tau_m,T})
 \leq C T^{5a}(\log T)^2
             e^{-cT^{\beta-\vartheta-\eta}}
       +C_\eta T^{-1+\vartheta+2\eta}.
\end{equation}
Indeed, on this interval the dissipation rate is at least
$c_\eta T^{-\vartheta-\eta}$ and the additive cooling term is at most
$C_\eta T^{-1+\eta}$. In particular,
\[
 \int_{\tau_m}^T a_\eta(u)\dd u
 \geq c_\eta T^{-\vartheta-\eta}(T-\tau_m)
 \geq c_\eta T^{\beta-\vartheta-\eta}.
\]
The source integral in
\eqref{eq:restartintegral} is bounded by
\[
 C_\eta T^{-1+\eta}\int_0^\infty
          e^{-c_\eta T^{-\vartheta-\eta}u}\dd u
 \leq C_\eta T^{-1+\vartheta+2\eta}.
\]
This proves \eqref{eq:restart}, uniformly in the restarting numerical law.

Fix $0<\zeta<\delta$ and take a smooth function $f:\R^{3d}\to[0,1]$
depending only on $U(x)$, equal to zero when $U(x)\leq\delta-\zeta$ and equal
to one when $U(x)\geq\delta$. Quadratic confinement makes the transition
region compact in $x$, so $f$ is globally Lipschitz. For all $j$ in the window,
Lemma~\ref{lem:smoothing} shows that
$g_j=P_{\tau_{j+1},T}f$ satisfies
\begin{equation}\label{eq:weightedlip}
 |g_j(y)-g_j(z)|
 \leq C_{\delta,\zeta}\sqrt{\log T}\,
              \left(1+|y|+|z|\right)|y-z|.
\end{equation}
This also holds for the final intervals of length less than one, using the
short-time part of that lemma instead of its smoothing part.

The definitions in Lemma~\ref{lem:local} give
$\nu_{j+1}=\nu_jQ_{\tau_j,h_j}$, while the Chapman--Kolmogorov identity for
the exact process \eqref{eq:model} gives
$P_{\tau_j,T}=P_{\tau_j,\tau_{j+1}}P_{\tau_{j+1},T}$. Telescoping these
numerical and exact transitions gives the exact identity
\[
 \nu_k f-\nu_mP_{\tau_m,T}f
 =\sum_{j=m}^{k-1}\nu_j
           \left(Q_{\tau_j,h_j}-P_{\tau_j,\tau_{j+1}}\right)P_{\tau_{j+1},T}f.
\]
Write $\left(S_j^{\rm num},S_j^{\rm ex}\right)$ for the conditional endpoint coupling
from Lemma~\ref{lem:local} given $S_j=s$. Its two marginals satisfy
\[
 \E_s\left[\left(1+\left|S_j^{\rm num}\right|+\left|S_j^{\rm ex}\right|\right)^2\right]
 \leq C\left(1+|s|^2\right).
\]
Therefore \eqref{eq:weightedlip}, Cauchy--Schwarz inequality, and \eqref{eq:local} give
\[
 \left|\left(Q_{\tau_j,h_j}-P_{\tau_j,\tau_{j+1}}\right)g_j(s)\right|
 \leq C\sqrt{\log T}(1+|s|)
       \left(h_j^3\left(1+|s|^2\right)
                    +\frac{h_j^{3/2}}{T\log^{3/2}T}\right).
\]
Integrating against $\nu_j$ uses the uniform third moment already proved.
Finally $\sum_{j=m}^{k-1}h_j=T-\tau_m\asymp T^\beta$ and
$h_j\asymp T^{-a}$ throughout the window. Thus
\begin{align}\label{eq:telescoping}
 |\nu_k f-\nu_mP_{\tau_m,T}f|
 &\leq C_{\delta,\zeta}\sqrt{\log T}
       \sum_{j=m}^{k-1}\left(h_j^3+
                      \frac{h_j^{3/2}}{T\log^{3/2}T}\right)\notag\\
 &\leq C_{\delta,\zeta}\left(
       \sqrt{\log T}\cdot T^{\beta-2a}
              +\frac{T^{\beta-1-a/2}}{\log T}\right).
\end{align}
This calculation retains the temperature variation within each numerical step.

We next bound $\nu_mP_{\tau_m,T}f$. Since
$0\leq f(x,v,z)\leq\mathbf 1_{\{U(x)>\delta-\zeta\}}$, Pinsker's inequality and
the bound
$\KL(\nu_mP_{\tau_m,T}\mid\pi_{\eps_T})
\leq\eps_T\cH_{\eps_T}(\nu_mP_{\tau_m,T})/A_0$, which follows from
\eqref{eq:functional} and $\cI_M\geq0$, give
\[
 \nu_mP_{\tau_m,T}f
 \leq\pi_{\eps_T}(U>\delta-\zeta)
      +\sqrt{\frac{\eps_T}{2A_0}\cH_{\eps_T}(\nu_mP_{\tau_m,T})}.
\]
For any small $\zeta'>0$, the Gibbs tail estimate \eqref{eq:gibbstail},
applied with $(\delta-\zeta,\zeta'/2)$ in place of $(\delta,\eta)$, and the
first relation in \eqref{eq:coolingassumption} give, for all large $T$,
\[
 \pi_{\eps_T}(U>\delta-\zeta)\leq CT^{-(\delta-\zeta-\zeta')/E}.
\]
For the second term, we insert \eqref{eq:restart} and use
$\sqrt{x+y}\leq\sqrt x+\sqrt y$ and $\eps_T\leq\eps_0$:
\begin{equation}\label{eq:stretchedexponential}
 \sqrt{\frac{\eps_T}{2A_0}\cH_{\eps_T}(\nu_mP_{\tau_m,T})}
 \leq CT^{5a/2}\log T\,e^{-cT^{\beta-\vartheta-\eta}/2}
      +C_\eta T^{-(1-\vartheta)/2+\eta}.
\end{equation}
The first term on the right-hand side of \eqref{eq:stretchedexponential} is
the stretched-exponential term; it comes from the square root of the first
term on the right-hand side of \eqref{eq:restart}.
Since $\beta-\vartheta-\eta>0$, for
every fixed $N$ there is $C_N<\infty$ such that
\[
 CT^{5a/2}\log T\,e^{-cT^{\beta-\vartheta-\eta}/2}\leq C_NT^{-N},
 \qquad T\geq1.
\]
Consequently,
\[
 \nu_mP_{\tau_m,T}f
 \leq C T^{-(\delta-\zeta-\zeta')/E}
       +C_\eta T^{-(1-\vartheta)/2+\eta}
       +C_NT^{-N}.
\]
Since $\mathbb P(U(X_k)>\delta)\leq\nu_kf$, combine
this with \eqref{eq:telescoping} to obtain
\begin{align}\label{eq:fourratebound}
 \mathbb P(U(X_k)>\delta)
 \leq C\bigg(&T^{-(\delta-\zeta-\zeta')/E}
 +T^{-(1-\vartheta)/2+\eta}
 +\sqrt{\log T}\,T^{-(2a-\beta)}\notag\\
 &+(\log T)^{-1}T^{-(1+a/2-\beta)}+T^{-N}\bigg).
\end{align}
The four polynomial exponents in \eqref{eq:fourratebound} before the
parameter limits
are
\[
 \frac{\delta-\zeta-\zeta'}{E},\qquad
 \frac{1-\vartheta}{2}-\eta,\qquad
 2a-\beta,\qquad 1+\frac a2-\beta.
\]
Letting $\beta\downarrow\vartheta$ and
$\zeta,\zeta',\eta\downarrow0$ shows that the first three converge to the
three exponents in \eqref{eq:discreterate}. The fourth does not restrict the
rate because
\[
 1+\frac a2-\vartheta-\frac{1-\vartheta}{2}
 =\frac{1+a-\vartheta}{2}>0.
\]
The factors involving $\log T$ and the parameter differences are absorbed by
$T^\alpha$. This proves \eqref{eq:discreterate}.
The condition $a>\vartheta/2$ makes every relevant exponent positive.
The iteration statement follows from the grid asymptotics.
\end{proof}

%%%%%%%%%%%%%%%%%%%%%%%%%%%%%%%%%%%%%%%%%%%%%%
\section{Numerical Illustration}\label{sec:numerics}

In this section, we conduct numerical experiments to compare overdamped Langevin dynamics (OLD) \eqref{eqn:OLD}, 
underdamped Langevin dynamics (ULD) \eqref{eqn:ULD}, 
together with the third-order Langevin
dynamic \eqref{eq:model}. Thus every frozen process has position marginal
proportional to $e^{-U/\eps}$. We fix OLD mobility and the kinetic mass at one;
this choice defines the reference physical clock, and any other constant OLD
mobility corresponds to a rescaling of that clock. The remaining kinetic
parameters are selected on validation paths and then held fixed on independent
evaluation paths, since finite-time annealing performance can change substantially
with friction.
The fixed-temperature distinction is already provided by the sharp
Eyring--Kramers theory: Wang and Zhu \cite{WangZhu2026} show that, in an explicit
matched-parameter regime, the third-order transition-time prefactor is strictly
smaller than its underdamped counterpart. We use this fixed-temperature result as
the theoretical benchmark and focus the experiment on whether the advantage
persists under the genuinely time-inhomogeneous cooling schedule proved in this
paper. 

For OLD we use Euler--Maruyama. For ULD we use the UBU splitting
\eqref{eq:ubustep}, with numerical friction denoted by $\eta$ and step
temperature $\eps_k=\eps_{\tau_k}$. Here $\eta$ is the numerical value of
$\gamma_{\rm u}$ in Appendix~\ref{sec:ubuappendix}, rather than the small
exponent-loss parameter in the proofs. We choose
UBU to align the numerical
treatment of the two kinetic dynamics: like the midpoint third-order scheme,
it uses one interior gradient evaluation per step and exactly resolves its
linear Gaussian subdynamics. This controls for gradient budget and force
quadrature, so the comparison more directly reflects the underlying dynamics.
For third-order Langevin dynamics, we use Algorithm~1 of Mou et al.\
\cite{MouEtAl2021}, namely the Gaussian endpoint \eqref{eq:moustep} with the
midpoint force integral \eqref{eq:midpoint}. The strong-order-two local
structure of UBU is discussed in Lemma~\ref{lem:ubulocal}; see
\cite{SanzSernaZygalakis2021,PaulinWhalley2024}.

We tune the parameters over the following validation grids:
\[
 \eta\in\{0.1,0.25,0.5,1,2,4\},\qquad
 \lambda\in\{0.25,0.5,1,1.5,2\},\qquad
 \gamma\in\{0.25,0.5,1,2,4,8\}.
\]
The score is terminal success at physical time $300$, with ever-reached success
and restricted mean hitting time as tie-breakers. Figure~\ref{fig:tuning} records
the scores on $1{,}200$ validation paths; the evaluation uses independent paths.

\begin{figure}[H]
\centering
\includegraphics[width=\textwidth]{figures/parameter_tuning.png}
\caption{Annealing parameter tuning on the validation set. Left: UBU terminal
success as a function of ULD friction. Right: the third-order product grid; the
cross marks the selected pair. Evaluation uses independent paths.}
\label{fig:tuning}
\end{figure}

%%%%%%%%%%%%%%%%%%%%%%%%%%%%%%%%%%%%%%%%%%%%%%%%%%%%%%%%%%%%%%%%%%
\subsection{Annealing with the theory-driven cooling schedule}\label{sec:toy:example}

To test annealing under decreasing temperature, we consider a toy example in this section
by using the smooth confining
tilted double well
\begin{equation}\label{eq:numericalpotential}
 U(x):=\frac{x^2}{20}-0.65e^{-(x+2)^2}-0.45e^{-(x-2)^2}.
\end{equation}
To match the normalization $\min U=0$ in
Assumption~\ref{ass:potential}, the theoretical potential for this experiment
is understood to be $U-U(x_\star)$. We retain the unshifted values below for
readability. This additive shift changes neither the gradient and the dynamics
nor the normalized Gibbs law, objective gaps, or critical depths.
Its global minimum is at $x_\star=-1.854300$, its nonglobal minimum is at
$x_{\rm loc}=1.792080$, and the intervening saddle is at $x_s=0.082433$.
Their objective values are $-0.464426$, $-0.270384$, and $-0.019547$,
respectively. Hence the one-dimensional landscape has critical depth
\[
 D_{\rm crit}=U(x_s)-U(x_{\rm loc})=0.250836.
\]
We take the canonical schedule dictated by Theorem~\ref{thm:main},
\[
 \eps_t=\frac{E}{\log(t+e)},\qquad E=0.28>D_{\rm crit},
\]
and start every path at $x_{\rm loc}$ with zero auxiliary variables. The target
$G_{0.08}=\{x:U(x)-U(x_\star)\leq0.08\}$ excludes the local well, whose gap is
$0.194042$. At $T=300$ the temperature has fallen to $0.049013$.

Validation selects $\eta=0.25$ for ULD and $(\lambda,\gamma)=(1,4)$ for third
order. We then use $4{,}000$ new paths per method and record terminal success,
ever-reached success, the restricted mean first-passage time, the terminal
objective gap, and completed left--right crossings. A crossing is counted only
after the path moves from $x>1$ to $x<-1$, or conversely, to avoid rapid
recrossings near the saddle.

The first protocol holds the physical-time discretization fixed: all methods use
$h_k=0.05e^{1/2}(\tau_k+e)^{-1/2}$ and the same horizon $T=300$. The second
protocol holds the computational budget at $40{,}000$ new gradient evaluations.
Its grids use $h_k=0.05e^a(\tau_k+e)^{-a}$, so $h_0=0.05$ for every method.
OLD and UBU use the shared baseline exponent $1/2$, while the third-order
method uses exponent
$a_\star=(1+D_{\rm crit}/E)/4=0.473961$, the smallest exponent in
Theorem~\ref{thm:discrete} that retains the continuous-time rate. This tests the
iteration-count mechanism suggested by Remark~\ref{rem:stepsizecomparison};
because it changes both the attained physical time and terminal temperature, it
evaluates complete method--grid pairs, incorporating both the attained physical
time and terminal temperature. For these parameters, the sufficient convergence
condition in Theorem~\ref{thm:ubu} is
$a>2D_{\rm crit}/(3E)=0.597229$, and the rate-preserving threshold is
$a_{\rm u,*}=(1+D_{\rm crit}/E)/3=0.631948$. Thus the shared baseline
$a=1/2$ lies outside the asymptotic coverage of the present UBU theorem and
gives UBU a slower-decaying step schedule than its sufficient threshold requires.
At the equal gradient budget, UBU reaches $T=287.762$, while the third-order
grid reaches $T=309.029$.

\begin{table}[H]
\centering
\caption{Theory-guided annealing on \eqref{eq:numericalpotential}, with standard
errors in parentheses. ``Ever'' means that $G_{0.08}$ was reached by time $T$,
and RMST is $\E[\min\{\tau,T\}]$. Each row uses $4{,}000$ evaluation paths.}
\label{tab:annealing}
\scriptsize
\resizebox{\textwidth}{!}{%
\begin{tabular}{@{}llccccccc@{}}
\toprule
Protocol & Method & $N_\nabla$ & $T$ & Terminal & Ever
& RMST & $\E[U(X_T)-U_\star]$ & Mean crossings \\
\midrule
Common horizon & OLD & $42{,}557$ & $300.000$ & $0.324\;(0.007)$
& $0.356\;(0.008)$ & $216.9\;(1.9)$ & $0.1540\;(0.0016)$ & $0.371\;(0.008)$ \\
& ULD (UBU) & $42{,}557$ & $300.000$ & $0.402\;(0.008)$
& $0.459\;(0.008)$ & $195.1\;(2.0)$ & $0.1364\;(0.0017)$ & $0.492\;(0.009)$ \\
& Third order & $42{,}557$ & $300.000$ & $0.421\;(0.008)$
& $0.474\;(0.008)$ & $191.5\;(2.0)$ & $0.1334\;(0.0017)$ & $0.505\;(0.009)$ \\
\addlinespace
Equal budget & OLD & $40{,}000$ & $287.762$ & $0.326\;(0.007)$
& $0.356\;(0.008)$ & $209.3\;(1.8)$ & $0.1547\;(0.0016)$ & $0.379\;(0.008)$ \\
& ULD (UBU) & $40{,}000$ & $287.762$ & $0.405\;(0.008)$
& $0.457\;(0.008)$ & $186.5\;(1.9)$ & $0.1363\;(0.0017)$ & $0.489\;(0.009)$ \\
& Third order & $40{,}000$ & $309.029$ & $0.415\;(0.008)$
& $0.469\;(0.008)$ & $199.7\;(2.0)$ & $0.1342\;(0.0017)$ & $0.496\;(0.009)$ \\
\bottomrule
\end{tabular}%
}
\end{table}

\begin{figure}[H]
\centering
\includegraphics[width=\textwidth]{figures/annealing_three_way_comparison.png}
\caption{Theory-guided annealing with tolerance $0.08$. Left and center: success
curves under the common-horizon protocol. Right: terminal and ever-reached
success at an equal gradient budget. Error bars in the right panel are one
Monte Carlo standard error.}
\label{fig:annealing}
\end{figure}

At the common physical horizon, all reported point estimates favor the tuned
third-order process over tuned UBU: terminal success is $0.421$ versus $0.402$,
ever-reached success is $0.474$ versus $0.459$, restricted mean hitting time is
$191.5$ versus $195.1$, and the mean terminal gap is $0.1334$ versus $0.1364$.
The terminal-success difference (third-order Langevin method minus UBU) is $0.0198$ with an
independent-sample standard error $0.0110$, giving a 95\% normal interval
approximately $[-0.002,0.041]$. At the equal gradient budget, third order again
has the higher terminal and ever-reached point estimates, $0.415$ versus $0.405$
and $0.469$ versus $0.457$, respectively; the terminal-success difference is
$0.0098$ with standard error $0.0110$. These intervals quantify the current
Monte Carlo uncertainty in the kinetic ranking. Moreover, the equal-budget
third-order grid reaches a later physical time and lower terminal temperature,
so that protocol compares complete method--grid pairs. Both kinetic methods
outperform OLD for terminal and first-passage criteria in this parameter regime.

As a discretization check, halving the initial step and using $1{,}000$ additional
paths at the common horizon gives terminal-success estimates $0.308$, $0.423$,
and $0.412$ for OLD, UBU, and third-order Langevin method, and ever-reached estimates $0.352$,
$0.481$, and $0.459$. This smaller independent sample gives higher UBU point
estimates, with the differences covered by its larger Monte Carlo uncertainty.

%%%%%%%%%%%%%%%%%%%%%%%%%%%%%%%%%%%%%%%%%%
\subsection{A neural-network example with synthetic data}\label{sec:NN:synthetic}

We next consider a high-dimensional nonconvex objective given by the
regularized empirical risk of a two-layer neural network. Let
$\widetilde x_i=(x_i,1)\in\R^{21}$ and
\begin{equation}\label{eq:nnpotential}
 f_\theta(x_i):=\frac{A}{\sqrt m}\sum_{j=1}^m
 \tanh(b_j)\tanh\left(w_j^\top\widetilde x_i\right),\qquad
 U(\theta):=\frac{1}{2n}\sum_{i=1}^n(f_\theta(x_i)-y_i)^2
 +\frac{\mu}{2}\|\theta\|^2.
\end{equation}
We take input dimension $20$, $n=32$, $m=4$, $A=2$, and $\mu=10^{-3}$,
so $\theta=(w_1,\ldots,w_m,b_1,\ldots,b_m)\in\R^{88}$.  The targets are
generated once from a fixed $12$-unit teacher network, with
$x_i\sim\mathcal N(0,I_{20}/20)$ and data seed $11$.  Both layers are
trainable; the parametrization $\tanh(b_j)$ merely bounds the effective output
weights.  The risk is nonconvex, while the bounded activation and output
parametrization make its nonquadratic derivatives bounded.  Together with the
quadratic regularizer, this also gives the smoothness and dissipativity
structure assumed in our analysis.

A preliminary search from $96$ random starts identifies three recurrent
local-minimum levels, with
objective values $0.012778835$, $0.012810789$, and $0.012842665$, reached by
$39$, $29$, and $27$ starts, respectively.  All annealing paths begin at a
representative of the highest of these three levels.  We then apply the same
deterministic L-BFGS-B ``quench'' to every stochastic endpoint. The lowest
level found by the preliminary search defines the reference basin, and a run
is classified as successful when its quenched value is within $10^{-6}$ of
that level.

On $96$ separate validation paths per method, a common grid
$E\in\{0.004,0.008,0.012,0.016,0.024\}$ selects $E=0.024$ by the
method-symmetric rule of making pooled best-basin success closest to $0.4$.
We then tune the kinetic parameters specifically for this neural-network
problem.  A first stage evaluates
\[
 \eta\in\{0.1,0.25,0.5,1,2,4\},\qquad
 \lambda\in\{0.25,0.5,1,1.5,2\},\qquad
 \gamma\in\{0.25,0.5,1,2,4,8\}
\]
on $64$ validation paths per candidate, using common random numbers within
each method.  The three best UBU and third-order candidates are then compared
on $192$ new validation paths per candidate.  Maximizing confirmed best-basin
success, with the mean quenched objective as tie-breaker, selects
$\eta=0.1$ and $(\lambda,\gamma)=(0.5,4)$.
The held-out experiment uses $768$ paths per method, the shared schedule
$\eps_t=0.024/\log(t+e)$, the common grid
$h_k=0.05e^{1/2}(\tau_k+e)^{-1/2}$, and horizon $T=50$.  Hence every method
uses $3{,}059$ new gradient evaluations.

\begin{table}[H]
\centering
\caption{Held-out annealing results for the $88$-dimensional neural-network
potential \eqref{eq:nnpotential}.  Parentheses contain one Monte Carlo standard
error.  The objective gap is measured from the best level found in the
preliminary multistart search and is reported after deterministic quenching.}
\label{tab:nnannealing}
\begin{tabular}{@{}l c @{\hspace{2em}} c@{}}
\toprule
Method & Best-basin probability & Mean quenched gap $(\times10^{-3})$ \\
\midrule
OLD & $0.243\;(0.015)$ & $0.0445\;(0.0016)$ \\
ULD (UBU) & $0.388\;(0.018)$ & $0.0352\;(0.0023)$ \\
Third-order LD & $0.418\;(0.018)$ & $0.0377\;(0.0027)$ \\
\bottomrule
\end{tabular}
\end{table}

\begin{figure}[H]
\centering
\includegraphics[width=\textwidth]{figures/neural_network_annealing.png}
\caption{Neural-network experiment.  Left: local
minimum levels identified by the preliminary multistart search.  Center:
held-out best-basin probabilities, with one-standard-error bars.  Right:
empirical distributions of the quenched endpoint objective gaps; curves farther
to the upper left are better.}
\label{fig:nnannealing}
\end{figure}

As a same-random-number sensitivity analysis, the parameters inherited from the
one-dimensional experiment, $\eta=0.25$ and $(\lambda,\gamma)=(1,4)$, give
best-basin probabilities $0.383$ for UBU and $0.378$ for third-order Langevin method on the
held-out noise streams.  The network-specific tuning therefore changes the
corresponding point estimates by $0.005$ and $0.040$, respectively.

Relative to OLD, the best-basin improvements are $0.145$ for UBU and $0.174$
for the third-order LD method, with independent-sample 95\% normal intervals
$[0.099,0.190]$ and $[0.128,0.221]$, respectively. The third-order best-basin
probability exceeds the UBU point estimate by $0.030$, with a 95\% normal
interval of $[-0.019,0.079]$. Halving the initial step and using $192$
new paths per method gives best-basin estimates $0.255$, $0.380$, and $0.427$
for OLD, UBU, and third-order LD method.

%%%%%%%%%%%%%%%%%%%%%%%%%%%%%%%%%%%%%%%%%%%%%%%
\subsection{A neural-network example with real data}\label{sec:NN:real}

We next retain the nonconvex neural-network potential in Section~\ref{sec:NN:synthetic} while replacing the
synthetic observations by the Connectionist Bench (Sonar, Mines versus Rocks)
data from the UCI Machine Learning Repository \cite{GormanSejnowski1988Sonar}.
The data contain $208$ observations with $60$ real-valued features and binary
rock/mine labels.  A fixed stratified split (seed $151$) assigns $160$
observations to training and $48$ to testing.  Each feature is standardized
using only the training sample and the resulting vectors are divided by
$\sqrt{60}$; labels are encoded as $-1$ and $1$.

We reuse the objective in \eqref{eq:nnpotential}, now with
$\widetilde x_i\in\R^{61}$, $n=160$, $m=4$, $A=2$, and $\mu=10^{-3}$.
Consequently, $\theta\in\R^{248}$.  The bounded tanh parametrization and
quadratic regularizer retain the smoothness and dissipativity structure, while
the empirical risk remains nonconvex.  A deterministic search from $64$ random
starts finds $39$ distinct local-minimum levels at tolerance $10^{-6}$; the
best level occurs from $12$ starts. Every annealing path begins at a recurrent
higher level. As above, the deterministic endpoint quench assigns each path
to a local-minimum basin. The lowest value found by the multistart search is
the reference level, and success means a quenched objective within $10^{-4}$
of that value.

Temperature and kinetic parameters are selected on validation paths, with
independent paths reserved for final evaluation.
On $24$ validation paths per method, the method-symmetric rule of pooled
best-basin success closest to $0.35$ selects $E=0.24$ from
$\{0.08,0.12,0.16,0.20,0.24,0.30,0.40\}$.
A first tuning stage uses $32$ paths per candidate over
$\eta\in\{0.05,0.1,0.25,0.5,1\}$ and
$(\lambda,\gamma)\in\{0.25,0.5,1\}\times\{2,4,8\}$; the two finalists for
each kinetic method are compared on $96$ fresh paths.  This selects
$\eta=0.05$ and $(\lambda,\gamma)=(1,4)$.  The held-out evaluation uses $256$
new paths per method, $\eps_t=0.24/\log(t+e)$,
$h_k=0.05e^{1/2}(\tau_k+e)^{-1/2}$, and horizon $T=20$, corresponding to $839$
gradient evaluations per path.

\begin{table}[H]
\centering
\caption{Held-out results for the $248$-dimensional UCI Sonar potential.
Parentheses contain one Monte Carlo standard error across independent
annealing paths.  Test-accuracy errors are conditional on the fixed $48$-record
test split.}
\label{tab:sonarannealing}
\begin{tabular}{@{}l c c c@{}}
\toprule
Method & Best-basin probability & Mean quenched objective & Test accuracy \\
\midrule
OLD & $0.180\;(0.024)$ & $0.10671\;(0.00035)$ & $0.824\;(0.003)$ \\
ULD (UBU) & $0.277\;(0.028)$ & $0.10427\;(0.00033)$ & $0.834\;(0.003)$ \\
Third-order LD & $0.352\;(0.030)$ & $0.10357\;(0.00034)$ & $0.843\;(0.003)$ \\
\bottomrule
\end{tabular}
\end{table}

\begin{figure}[H]
\centering
\includegraphics[width=\textwidth]{figures/uci_sonar_annealing.png}
\caption{UCI Sonar experiment.  Left: local-minimum gaps found by the
$64$-start deterministic search.  Center: held-out best-basin probabilities.
Right: mean post-quench accuracy on the fixed test split.  Error bars in the
center and right panels are one Monte Carlo standard error across annealing
paths.}
\label{fig:sonarannealing}
\end{figure}

Relative to OLD, the UBU and third-order best-basin improvements are $0.098$ and
$0.172$, respectively.  The third-order point estimate is $0.074$ higher than
UBU.  The corresponding post-quench test-accuracy point estimates are
$0.824$, $0.834$, and $0.843$, again with the third-order LD method highest.  With the initial step
halved and $64$ fresh paths per method, best-basin estimates are $0.266$,
$0.281$, and $0.344$, while test accuracies are $0.835$, $0.830$, and $0.844$.
Thus, the third-order LD method remains the highest point estimate under step refinement.

%%%%%%%%%%%%%%%%%%%%%%%%%%%%%%%%%%%%%%%%%%%%%%%%%%%%%%%%

\section{Conclusion}\label{sec:conclusion}

We established global convergence guarantees of third-order Langevin dynamics for non-convex optimization
via simulated annealing with fixed friction and decreasing noise.  A three-block distorted entropy transfers
dissipation through the full chain, and an initial-density approximation covers
nonsmooth initial laws.  Logarithmic cooling then gives convergence in
probability at the barrier-controlled kinetic rate.  For the exact-force and
midpoint three-stage schemes, a cubic local endpoint estimate yields a less
restrictive sufficient step-size condition and a better sufficient iteration
exponent than both the existing kinetic frozen-force result and the UBU bound
proved by the same strong-coupling method.
The experiments complement these rate comparisons.  The third-order Langevin method
has higher terminal-success point estimates than UBU in the theory-guided
double well.  On a high-dimensional synthetic neural-network objective using synthetic data,
independent two-stage tuning and held-out evaluation show that both kinetic
methods outperform OLD in best-basin recovery, while the third-order Langevin
method has a higher point estimate than UBU; both improvements over OLD persist
after step refinement.
Using the real data, the same best-basin ordering is observed,
and the third-order Langevin method also gives the highest post-quench test-accuracy
point estimate. These qualitative conclusions persist under step refinement.

\bibliographystyle{alpha}
\bibliography{bibtex}

\appendix

%%%%%%%%%%%%%%%%%%%%%%%%%%%%%%%%%%%%%%%%%%%%%%%%%%%%%
\section{Underdamped UBU Comparison}\label{sec:ubuappendix}

This appendix provides the UBU result summarized in
Table~\ref{tab:introdiscretecomparison} and used in the comparison in
Remark~\ref{rem:stepsizecomparison}.

\subsection{Scheme and convergence result}

We now define the underdamped comparator used in
Table~\ref{tab:introdiscretecomparison} and Section~\ref{sec:numerics}.
In particular, let us define the UBU scheme $(\bar{X}_{k},\bar{V}_{k})_{k\geq 0}$ for ULD.
Fix a friction parameter $\gamma_{\rm u}>0$ and let
$P^{\rm u}_{t,T}$ denote the transition operator of the time-inhomogeneous ULD:
\begin{equation}\label{eq:uldannealed}
 \dd X_t=V_t\dd t,\qquad
 \dd V_t=\left(-\nabla U(X_t)-\gamma_{\rm u}V_t\right)\dd t
       +\sqrt{2\gamma_{\rm u}\eps_t}\,\dd B_t.
\end{equation}
At frozen temperature $\eps$, its invariant law is
\begin{equation}\label{eq:uldgibbs}
 \pi^{\rm u}_\eps(\dd x\,\dd v)
 \propto \exp\left[-\frac{U(x)+|v|^2/2}{\eps}\right]\dd x\,\dd v.
\end{equation}
For $r>0$, let $\mathcal U_r^\eps$ be the exact Gaussian flow over time $r$ of
the force-free part of \eqref{eq:uldannealed}. Writing
$q_r:=e^{-\gamma_{\rm u}r}$ and $F_r:=(1-q_r)/\gamma_{\rm u}$, its mean is
$(x+F_rv,q_rv)$ and its covariance is
\begin{equation}\label{eq:ubuflowcov}
 \eps\begin{pmatrix}
 \displaystyle \frac{2r}{\gamma_{\rm u}}
  -\frac{4(1-q_r)-(1-q_r^2)}{\gamma_{\rm u}^2}
 &\displaystyle \frac{(1-q_r)^2}{\gamma_{\rm u}}\\[3pt]
 \displaystyle \frac{(1-q_r)^2}{\gamma_{\rm u}}
 &1-q_r^2
 \end{pmatrix}\otimes\Id_d.
\end{equation}
Given $\bar S_k=(\bar{X}_{k},\bar{V}_{k})=(x,v)$, one frozen-temperature UBU step is
\begin{equation}\label{eq:ubustep}
 (y,w)=\mathcal U_{h_k/2}^{\eps_{\tau_k}}(x,v),\qquad
 w^+=w-h_k\nabla U(y),\qquad
 \bar S_{k+1}=\mathcal U_{h_k/2}^{\eps_{\tau_k}}\left(y,w^+\right),
\end{equation}
where the two half steps use independent Brownian increments. This is the
one-gradient symmetric splitting of Sanz-Serna and Zygalakis
\cite[Example~8]{SanzSernaZygalakis2021}. Its local-error analysis must retain
the centered stochastic term of order $h^{5/2}$; see also the correction
\cite{PaulinWhalley2024}.

The following theorem turns this centered local-error structure into a
quantitative annealing rate and step-size condition for the UBU scheme.

\begin{theorem}[Underdamped annealing with UBU]\label{thm:ubu}
Assume Assumptions~\ref{ass:potential}, \ref{ass:lsi}, and
\ref{ass:discrete}. For the UBU scheme \eqref{eq:ubustep}, take
\[
 h_k=c_hs_k^{-a},\qquad c_h>0,\quad a>0,
\]
with $c_h$ sufficiently small, and let the initial numerical law have a finite
fourth moment. Write $\vartheta=D/E$. If $a>2\vartheta/3$, then for every
$\delta,\alpha>0$ and all sufficiently large $k$,
\begin{equation}\label{eq:uburate}
 \mathbb P\left(U(\bar X_k)>\delta\right)
 \leq C_{\delta,\alpha}s_k^{-r_{\rm u}(\delta,a)+\alpha},
 \qquad
 r_{\rm u}(\delta,a)
 :=\min\left\{\frac\delta E,\frac{1-\vartheta}{2},
                         \frac{3a}{2}-\vartheta\right\}.
\end{equation}
In particular, $a\geq(1+\vartheta)/3$ retains the continuous-time exponent.
Moreover $s_k\asymp k^{1/(1+a)}$, so the corresponding iteration exponent is
$r_{\rm u}(\delta,a)/(1+a)$, up to an arbitrarily small loss.
\end{theorem}

\begin{proof}
The proof is given in Section~\ref{sec:proof_ubu}.
\end{proof}

Choosing the smallest step-decay exponent in Theorem~\ref{thm:ubu} that
preserves the continuous-time rate yields the iteration exponent in the following corollary.

\begin{corollary}[Rate-preserving UBU exponent]\label{cor:ubuiteration}
Under the assumptions of Theorem~\ref{thm:ubu}, let
\[
 r_{\rm c}(\delta):=\min\left\{\frac\delta E,
                    \frac{1-\vartheta}{2}\right\},
 \qquad a_{\rm u,*}:=\frac{1+\vartheta}{3}.
\]
Then the choice $a=a_{\rm u,*}$ gives
\[
 \mathbb P\left(U(\bar X_k)>\delta\right)
 \leq C_{\delta,\alpha}
 k^{-3r_{\rm c}(\delta)/(4+\vartheta)+\alpha}
\]
for all sufficiently large $k$.
\end{corollary}

\begin{proof}
At $a=a_{\rm u,*}$,
$3a/2-\vartheta=(1-\vartheta)/2$, while
$s_k\asymp k^{1/(1+a)}=k^{3/(4+\vartheta)}$. Substitute these identities in
\eqref{eq:uburate} and absorb the arbitrarily small loss into $\alpha$.
\end{proof}

\subsection{Proof of the UBU comparison}\label{sec:proof_ubu}

We collect the underdamped analogues of the local comparison,
grid-regularization, restart, and smoothing estimates. Constants in this
subsection may depend on $d,U,\gamma_{\rm u}$, and $c_h$.

We begin with the one-step strong error and uniform moment control for the
UBU chain in the next lemma.

\begin{lemma}[UBU local error and moments]\label{lem:ubulocal}
Let $Q^{\rm UBU}_{t,h}$ be one step of \eqref{eq:ubustep}, with temperature
frozen at $\eps_t$, and couple it synchronously to \eqref{eq:uldannealed}
started from $s=(x,v)$. For $h\leq h_*$ and all sufficiently large $t$,
\begin{equation}\label{eq:ubulocal}
 \left\|\bar S_h^{\rm UBU}-\bar S_{t+h}^{\rm exact}\right\|_{L^2(\mathbb P_s)}
 \leq C\sqrt{\eps_t}\,h^{5/2}
      +Ch^3\left(1+|s|^2\right)
      +\frac{Ch^{3/2}}{(t+t_0)\log^{3/2}(t+t_0)}.
\end{equation}
The $\mathcal{O}\left(\sqrt{\eps_t}\,h^{5/2}\right)$ contribution is an It\^o integral with
conditional mean zero. The numerical chain satisfies
\begin{equation}\label{eq:ubumoments}
 \sup_{k\geq0}\E\left|\bar S_k\right|^p<\infty,
 \qquad 1\leq p\leq4.
\end{equation}
\end{lemma}

\begin{proof}
First freeze the temperature at $\eps_t$. We give the conditional local
calculation because the estimates in the fixed-temperature UBU literature are
usually stated after averaging over an equilibrium law. Write
$G=\nabla U$, $\mathsf H_U=\nabla^2U$, and let $(X_r,V_r)$ be the exact frozen ULD
started from $s$. The globally Lipschitz drift and the bounded noise amplitude
give, for $p=2,4$,
\begin{equation}\label{eq:ubuconditionalmoments}
 \sup_{0\leq r\leq h_*}
 \left(\|X_r\|_{L^p(\mathbb P_s)}+\|V_r\|_{L^p(\mathbb P_s)}
       +\|G(X_r)\|_{L^p(\mathbb P_s)}\right)
 \leq C_p(1+|s|).
\end{equation}
Let $\left(X_r^0,V_r^0\right)$ be the force-free process driven by the same Brownian
path and put $Y=X_{h/2}^0$, the UBU force point. Variation of constants gives
\begin{equation}\label{eq:ubuforcepoint}
\begin{aligned}
 X_{h/2}-Y&=-\int_0^{h/2}F_{h/2-r}G(X_r)\dd r,\\
 \|X_{h/2}-Y\|_{L^p(\mathbb P_s)}
 &\leq C_ph^2(1+|s|).
\end{aligned}
\end{equation}
The exact and UBU velocity endpoints therefore differ by
\[
 \int_0^h q_{h-r}G(X_r)\dd r-hq_{h/2}G(Y).
\]
For $\chi(r)=q_{h-r}G(X_r)$, the symmetric midpoint identity reads
\begin{equation}\label{eq:ubumidpointidentity}
 \int_0^h\chi(r)\dd r-h\chi(h/2)
 =\int_{h/2}^h\dd u\int_{h/2}^u\dd v
       \int_{h-v}^{v}\dd\dot\chi(r).
\end{equation}
For smooth paths this identity follows from two applications of the
fundamental theorem of calculus. In the present setting,
$\dot\chi(r)=q_{h-r}\{\gamma_{\rm u}G(X_r)+\mathsf H_U(X_r)V_r\}$ is a continuous
semimartingale. Since $X$ has finite variation, It\^o's formula gives
\begin{equation}\label{eq:ubuchiito}
 \dd\dot\chi(r)=a_\chi(r)\dd r+\sigma_\chi(r)\dd B_r,
\end{equation}
with finite-variation and martingale integrands
\begin{align*}
 a_\chi(r)&:=q_{h-r}\left[
 \gamma_{\rm u}^2G(X_r)+\gamma_{\rm u}\mathsf H_U(X_r)V_r
 +\nabla^3U(X_r)[V_r,V_r]-\mathsf H_U(X_r)G(X_r)\right],\\
 \sigma_\chi(r)&:=\sqrt{2\gamma_{\rm u}\eps_t}\,q_{h-r}\mathsf H_U(X_r).
\end{align*}
Since $0<q_{h-r}\leq1$, $\|\mathsf H_U\|_{\rm op}\leq K$, and
$\|\nabla^3U\|_\infty<\infty$, the conditional moment bound
\eqref{eq:ubuconditionalmoments} with $p=2,4$ gives, for $0\leq r\leq h$,
\begin{equation}\label{eq:ubuchiintegrands}
\begin{aligned}
 \|a_\chi(r)\|_{L^2(\mathbb P_s)}
 &\leq C\left(\|G(X_r)\|_{L^2(\mathbb P_s)}+\|V_r\|_{L^2(\mathbb P_s)}
       +\|V_r\|_{L^4(\mathbb P_s)}^2\right)
 \leq C\left(1+|s|^2\right),\\
 \|\sigma_\chi(r)\|_{\rm F}^2
 &\leq2\gamma_{\rm u}\eps_t\,dK^2,
\end{aligned}
\end{equation}
where $\|\cdot\|_{\rm F}$ denotes the Frobenius norm. Let
\[
 \mathcal T_h:=\left\{(u,v,r):\ \frac h2\leq v\leq u\leq h,\
 h-v\leq r\leq v\right\},
 \qquad |\mathcal T_h|=\frac{h^3}{24},
\]
be the triangular integration domain in \eqref{eq:ubumidpointidentity}. By
\eqref{eq:ubuchiintegrands}, we have:
\begin{equation}\label{eq:ubufubiniconditions}
 \int_{\mathcal T_h}\E_s|a_\chi(r)|\,\dd(u,v,r)
 \leq\frac{h^3}{24}\,C\left(1+|s|^2\right),
 \qquad
 \int_{\mathcal T_h}\E_s\|\sigma_\chi(r)\|_{\rm F}^2\,\dd(u,v,r)
 \leq\frac{\gamma_{\rm u}\eps_t dK^2}{12}h^3.
\end{equation}
The first bound justifies Fubini's theorem for the finite-variation part,
and the second is the square-integrability hypothesis of the stochastic
Fubini theorem for the martingale part. For fixed $r\in[0,h]$, the
$(u,v)$-section of $\mathcal T_h$ has area
\[
 \kappa_h(r):=\frac12\min\{r,h-r\}^2.
\]
Therefore \eqref{eq:ubumidpointidentity} holds in $L^2(\mathbb P_s)$ in the form
\begin{equation}\label{eq:ubuvelocitysplit}
 \int_0^h\chi(r)\dd r-h\chi(h/2)
 =\int_{\mathcal T_h}a_\chi(r)\,\dd(u,v,r)
  +\int_0^h\kappa_h(r)\sigma_\chi(r)\dd B_r.
\end{equation}
Thus the stochastic part of the velocity error is
\[
 \alpha_h^v:=\int_0^h\kappa_h(r)\sigma_\chi(r)\dd B_r
 =\sqrt{2\gamma_{\rm u}\eps_t}
 \int_{h/2}^h\dd u\int_{h/2}^u\dd v
       \int_{h-v}^{v}q_{h-r}\mathsf H_U(X_r)\dd B_r.
\]
It is conditionally centered. It\^o's isometry,
\eqref{eq:ubuchiintegrands}, and $\int_0^h\kappa_h(r)^2\dd r=h^5/320$ yield
\[
 \|\alpha_h^v\|_{L^2(\mathbb P_s)}^2
 =\E_s\int_0^h\kappa_h(r)^2\|\sigma_\chi(r)\|_{\rm F}^2\dd r
 \leq\frac{\gamma_{\rm u}\eps_t dK^2}{160}h^5,
 \qquad\text{so}\qquad
 \|\alpha_h^v\|_{L^2(\mathbb P_s)}\leq C\sqrt{\eps_t}\,h^{5/2}.
\]
Since $\chi(h/2)=q_{h/2}G(X_{h/2})$, the remaining part of the velocity error is
\[
 \beta_h^v:=\int_{\mathcal T_h}a_\chi(r)\,\dd(u,v,r)
 +hq_{h/2}\{G(X_{h/2})-G(Y)\}.
\]
Minkowski's integral inequality, \eqref{eq:ubuchiintegrands}, the bounded
Hessian, and the force-point estimate \eqref{eq:ubuforcepoint} give
\[
 \|\beta_h^v\|_{L^2(\mathbb P_s)}
 \leq\frac{h^3}{24}\sup_{0\leq r\leq h}\|a_\chi(r)\|_{L^2(\mathbb P_s)}
 +Kh\|X_{h/2}-Y\|_{L^2(\mathbb P_s)}
 \leq Ch^3\left(1+|s|^2\right).
\]

The position endpoints have the analogous difference
\[
 \int_0^hF_{h-r}G(X_r)\dd r-hF_{h/2}G(Y).
\]
Apply \eqref{eq:ubumidpointidentity} to $\psi(r):=F_{h-r}G(X_r)$. Since
$\frac{\dd}{\dd r}F_{h-r}=-q_{h-r}$,
\[
 \dot\psi(r)=-q_{h-r}G(X_r)+F_{h-r}\mathsf H_U(X_r)V_r,
\]
and It\^o's formula gives
$\dd\dot\psi(r)=a_\psi(r)\dd r+\sigma_\psi(r)\dd B_r$, where
\begin{align*}
 a_\psi(r)&:=F_{h-r}\left[\nabla^3U(X_r)[V_r,V_r]
              -\mathsf H_U(X_r)G(X_r)\right]
   -(1+q_{h-r})\mathsf H_U(X_r)V_r-\gamma_{\rm u}q_{h-r}G(X_r),\\
 \sigma_\psi(r)&:=\sqrt{2\gamma_{\rm u}\eps_t}\,F_{h-r}\mathsf H_U(X_r).
\end{align*}
Since $0\leq F_{h-r}\leq h-r\leq h$, the argument leading to
\eqref{eq:ubuchiintegrands} gives
\[
 \sup_{0\leq r\leq h}\|a_\psi(r)\|_{L^2(\mathbb P_s)}\leq C\left(1+|s|^2\right),
 \qquad
 \|\sigma_\psi(r)\|_{\rm F}^2\leq2\gamma_{\rm u}\eps_t\,dK^2h^2.
\]
Hence the analogue of \eqref{eq:ubufubiniconditions} holds for
$(a_\psi,\sigma_\psi)$, and, as in \eqref{eq:ubuvelocitysplit},
\[
 \int_0^h\psi(r)\dd r-h\psi(h/2)
 =\int_{\mathcal T_h}a_\psi(r)\,\dd(u,v,r)
  +\int_0^h\kappa_h(r)\sigma_\psi(r)\dd B_r.
\]
The stochastic part
$\alpha_h^x:=\int_0^h\kappa_h(r)\sigma_\psi(r)\dd B_r$ is conditionally
centered, and It\^o's isometry gives
\[
 \|\alpha_h^x\|_{L^2(\mathbb P_s)}^2
 \leq2\gamma_{\rm u}\eps_t\,dK^2h^2\int_0^h\kappa_h(r)^2\dd r
 =\frac{\gamma_{\rm u}\eps_t dK^2}{160}h^7,
\]
so that
\[
 \|\alpha_h^x\|_{L^2(\mathbb P_s)}\leq C\sqrt{\eps_t}\,h^{7/2}.
\]
The residual part of the position error is
\[
 \beta_h^x:=\int_{\mathcal T_h}a_\psi(r)\,\dd(u,v,r)
 +hF_{h/2}\{G(X_{h/2})-G(Y)\}.
\]
As for $\beta_h^v$, now also using $F_{h/2}\leq h/2$,
\[
 \|\beta_h^x\|_{L^2(\mathbb P_s)}
 \leq\frac{h^3}{24}\sup_{0\leq r\leq h}\|a_\psi(r)\|_{L^2(\mathbb P_s)}
 +\frac{Kh^2}{2}\|X_{h/2}-Y\|_{L^2(\mathbb P_s)}
 \leq Ch^3\left(1+|s|^2\right).
\]
Consequently, the full frozen one-step error has the
decomposition
\begin{equation}\label{eq:ubulocaldecomposition}
 \bar S_h^{\rm UBU}-S_h^{\rm frozen}=\alpha_h+\beta_h,\qquad
 \alpha_h:=\left(\alpha_h^x,\alpha_h^v\right),\quad
 \beta_h:=\left(\beta_h^x,\beta_h^v\right),\qquad
 \E_s\alpha_h=0,
\end{equation}
with
\[
 \|\alpha_h\|_{L^2(\mathbb P_s)}
 \leq C\sqrt{\eps_t}\left(h^{5/2}+h^{7/2}\right)
 \leq C\sqrt{\eps_t}\,h^{5/2},\qquad
 \|\beta_h\|_{L^2(\mathbb P_s)}
 \leq Ch^3(1+|s|^2).
\]
This conditional decomposition follows the local-error analysis of
Sanz-Serna and Zygalakis \cite[Section~7.6]{SanzSernaZygalakis2021}, while
the associated dimension-dependent weighted inner-product constants are
refined by Paulin and Whalley \cite{PaulinWhalley2024}.

It remains to unfreeze the temperature. The two exact linear half flows may
be realized with the restrictions of a common Brownian path to the two half
intervals. As in the proof of Lemma~\ref{lem:local},
\[
 |\sqrt{\eps_{t+u}}-\sqrt{\eps_t}|
 \leq \frac{Cu}{(t+t_0)\log^{3/2}(t+t_0)}.
\]
It\^o isometry for the velocity component, followed by one integration for
the position component, gives the last term in \eqref{eq:ubulocal}.

For the moment estimate, let
\[
 \mathcal V_{\rm u}(x,v)
 =1+U(x)+\frac12|v|^2+b\,x\cdot v+c|x|^2,
\]
where $b,c>0$ are sufficiently small. Dissipativity and the bounded Hessian
give $\mathcal V_{\rm u}\asymp1+|x|^2+|v|^2$ and, for the exact frozen ULD,
\[
 P^{{\rm u},\eps}_h\mathcal V_{\rm u}^p
 \leq e^{-c_ph}\mathcal V_{\rm u}^p+C_ph,
 \qquad p=1,2.
\]
The following coarse local bound preserves linear dependence on the initial
state:
\begin{equation}\label{eq:ubucoarselocal}
 \left\|\bar S_h^{\rm UBU}-S_h^{\rm frozen}\right\|_{L^q(\mathbb P_s)}
 \leq C_qh^2(1+|s|),\qquad q\geq2,
\end{equation}
where $S^{\rm frozen}$ is the exact ULD with temperature fixed at $\eps_t$.
Indeed, drive its force-free version $\left(X_r^0,V_r^0\right)$ and both UBU half
steps by consecutive pieces of the same Brownian path, and put
$Y=X_{h/2}^0$. The exact force-free formulas for the second half-step give
\[
 \bar V_h^{\rm UBU}=V_h^0-hq_{h/2}\nabla U(Y),\qquad
 \bar X_h^{\rm UBU}=X_h^0-hF_{h/2}\nabla U(Y).
\]
On the other hand, variation of constants for the exact frozen process gives
\[
 V_h^{\rm frozen}=V_h^0-\int_0^h q_{h-r}\nabla U(X_r)\dd r,
 \qquad
 X_h^{\rm frozen}=X_h^0-\int_0^h F_{h-r}\nabla U(X_r)\dd r.
\]
The bounded Hessian and the linear growth of $\nabla U$, together with the
standard finite-time moment estimate, yield uniformly for
$0\leq r\leq h\leq h_*$,
\[
 \left\|X_r-X_r^0\right\|_{L^q(\mathbb P_s)}\leq C_qh^2(1+|s|),\qquad
 \left\|X_r^0-Y\right\|_{L^q(\mathbb P_s)}\leq C_qh(1+|s|).
\]
Since $q_r$ and $F_r$ are Lipschitz on $[0,h_*]$, substituting these two
bounds into the preceding variation-of-constants formulas proves
\eqref{eq:ubucoarselocal}. Moreover, both coupled endpoints have linear
moment growth: for every $q\geq2$ and $0<h\leq h_*$,
\begin{equation}\label{eq:ubuendpointmoments}
 \left\|S_h^{\rm frozen}\right\|_{L^q(\mathbb P_s)}\leq C_q(1+|s|),
 \qquad
 \left\|\bar S_h^{\rm UBU}\right\|_{L^q(\mathbb P_s)}\leq C_q(1+|s|).
\end{equation}
The first bound in \eqref{eq:ubuendpointmoments} is the finite-time moment estimate behind
\eqref{eq:ubuconditionalmoments}, which holds for every $q\geq2$. The second bound in \eqref{eq:ubuendpointmoments}
follows from the first bound, \eqref{eq:ubucoarselocal}, and the triangle
inequality:
\[
 \left\|\bar S_h^{\rm UBU}\right\|_{L^q(\mathbb P_s)}
 \leq\left\|S_h^{\rm frozen}\right\|_{L^q(\mathbb P_s)}
 +\left\|\bar S_h^{\rm UBU}-S_h^{\rm frozen}\right\|_{L^q(\mathbb P_s)}
 \leq C_q(1+|s|)+C_qh^2(1+|s|).
\]
The mean-value theorem gives
\[
 \left|\mathcal V_{\rm u}^p(y)-\mathcal V_{\rm u}^p(y')\right|
 \leq C_p\left(1+|y|^{2p-1}+|y'|^{2p-1}\right)|y-y'|.
\]
Hence, H\"older's inequality, \eqref{eq:ubuendpointmoments} with $q=4p-2$,
and \eqref{eq:ubucoarselocal} with $q=2$ imply, for $p=1,2$,
\[
\begin{aligned}
 \left|Q^{\rm UBU}_{t,h}\mathcal V_{\rm u}^p(s)
       -P^{{\rm u},\eps_t}_h\mathcal V_{\rm u}^p(s)\right|
 &\leq C_p\left(1+\left\|\bar S_h^{\rm UBU}\right\|_{L^{4p-2}(\mathbb P_s)}^{2p-1}
      +\left\|S_h^{\rm frozen}\right\|_{L^{4p-2}(\mathbb P_s)}^{2p-1}\right)
      \left\|\bar S_h^{\rm UBU}-S_h^{\rm frozen}\right\|_{L^2(\mathbb P_s)}\\
 &\leq C_ph^2\left(1+|s|^{2p}\right)
 \leq C_ph^2\mathcal V_{\rm u}^p(s).
\end{aligned}
\]
Consequently,
\[
 Q^{\rm UBU}_{t,h}\mathcal V_{\rm u}^p
 \leq \left(1-c_ph+C_ph^2\right)\mathcal V_{\rm u}^p+C_ph.
\]
For all sufficiently small steps this yields the same discrete Lyapunov
recursion as in Lemma~\ref{lem:local}. The finite fourth moment at time zero
then proves \eqref{eq:ubumoments}, after absorbing the finite initial segment.
\end{proof}

Let $\cH^{\rm u}_\eps$ denote the kinetic distorted entropy inspired by
He, Tan, and Wu \cite[Eq.~(2.7)]{HeTanWu2024}. For the fixed friction
$\gamma_{\rm u}$, let
\[
 M_{\rm u}:=
 \begin{pmatrix}1&1\\1&1\end{pmatrix}\otimes\Id_d,\qquad
 A_{\rm u}:=
 2+\frac{2K+(1+\gamma_{\rm u}+K)^2+1}{\gamma_{\rm u}},
\]
and define
\begin{equation}\label{eq:ulddistorted}
 \cH^{\rm u}_\eps(\nu)
 :=\int_{\R^{2d}}
 \left|\nabla_x\log\frac{\dd\nu}{\dd\pi^{\rm u}_\eps}
       +\nabla_v\log\frac{\dd\nu}{\dd\pi^{\rm u}_\eps}\right|^2\dd\nu
   +\frac{A_{\rm u}}{\eps}\KL(\nu\mid\pi^{\rm u}_\eps).
\end{equation}
This value of $A_{\rm u}$ provides the uniform two-block dissipation
certificate used below to derive the stated convergence exponents.

The final Gaussian half-step 
in the next lemma provides grid-point smoothing and controls the
distorted entropy introduced above.

\begin{lemma}[UBU grid regularization]\label{lem:ubugaussian}
At every positive UBU grid point, the numerical law has a smooth positive
density and finite full relative Fisher information. Moreover, for all
sufficiently large $k$,
\begin{equation}\label{eq:ubugridfisher}
 \cH^{\rm u}_{\eps_{\tau_k}}\left(\mathcal L(\bar S_k)\right)
 \leq Cs_k^{3a}(\log s_k)^2.
\end{equation}
\end{lemma}

\begin{proof}
Conditionally on the state after the force kick in \eqref{eq:ubustep}, the
last $\mathcal U_{h/2}^{\eps_t}$ substep adds an independent centered Gaussian
with covariance $C^{\rm u}_{t,h}$ given by \eqref{eq:ubuflowcov} with
$r=h/2$. Taylor expansion at $h=0$ gives
\begin{equation}\label{eq:ubucovbound}
 C^{\rm u}_{t,h}\succeq
 c\eps_t\operatorname{diag}(h^3,h)\otimes\Id_d,\qquad
 \det C^{\rm u}_{t,h}\geq c\eps_t^{2d}h^{4d}.
\end{equation}
Thus the endpoint law is a mixture of Gaussians with the same nondegenerate
covariance. The mixture-score argument in Lemma~\ref{lem:gaussian} yields
\[
 \cI^{\rm raw}\left(\mathcal L(\bar S_k)\right)
 \leq\operatorname{tr}\left(\left(C^{\rm u}_{\tau_{k-1},h_{k-1}}\right)^{-1}\right)
 \leq\frac{C}{\eps_{\tau_{k-1}}h_{k-1}^3}.
\]
Since $|\nabla U(x)|\leq C(1+|x|)$, the relative Fisher information is at
most this quantity plus
$C\eps_{\tau_k}^{-2}\E(1+|\bar S_k|^2)$. The entropy-of-a-Gaussian-mixture
argument used in Lemma~\ref{lem:gaussian}, \eqref{eq:ubucovbound}, and
Lemma~\ref{lem:ubulocal} similarly give
\[
 \KL\left(\mathcal L(\bar S_k)\mid\pi^{\rm u}_{\eps_{\tau_k}}\right)
 \leq C\log s_k.
\]
Using $h_{k-1}\asymp s_k^{-a}$ and
$\eps_{\tau_k}^{-1}=\mathcal{O}(\log s_k)$ in \eqref{eq:ulddistorted} proves
\eqref{eq:ubugridfisher}. Smoothness and positivity follow from the final
Gaussian convolution.
\end{proof}

We record in the next lemma the restart decay and semigroup smoothing estimates needed to
propagate the local error over the final time window.

\begin{lemma}[ULD restart and smoothing]\label{lem:uldrestart}
Let $\rho$ be a law on $\R^{2d}$ with second moment bounded by $M_2$,
$\cH^{\rm u}_{\eps_s}(\rho)<\infty$, and finite full relative Fisher
information $\cI(\rho\mid\pi^{\rm u}_{\eps_s})<\infty$. If
$T-s\asymp T^\beta$, $s\asymp T$, and
$\beta>\vartheta$, then, for every sufficiently small $\eta>0$,
\begin{equation}\label{eq:uldrestart}
 \cH^{\rm u}_{\eps_T}(\rho P^{\rm u}_{s,T})
 \leq C\cH^{\rm u}_{\eps_s}(\rho)
          e^{-cT^{\beta-\vartheta-\eta}}
       +C_\eta T^{-1+\vartheta+2\eta},
\end{equation}
with constants uniform over such $\rho$. Moreover, for every bounded
measurable $f$ and $T\geq t+1$,
\begin{equation}\label{eq:uldsmoothing}
 \left|P^{\rm u}_{t,T}f(s)-P^{\rm u}_{t,T}f(s')\right|
 \leq C\eps_{t+1}^{-1/2}|s-s'|\|f\|_\infty.
\end{equation}
For $0\leq T-t\leq1$ and globally Lipschitz $f$,
$\operatorname{Lip}(P^{\rm u}_{t,T}f)\leq C\operatorname{Lip}(f)$.
\end{lemma}

\begin{proof}
We first prove the restart estimate. For a symmetric matrix $Q$ with
$\|Q\|_{\rm op}\leq K$, set
\[
 J_Q^{\rm u}:=
 \begin{pmatrix}0&-\Id_d\\Q&-\gamma_{\rm u}\Id_d\end{pmatrix},
 \qquad P_v:=\operatorname{diag}(0,\Id_d).
\]
Direct multiplication, using the $M_{\rm u}$ and $A_{\rm u}$ fixed above,
gives
\begin{align}\label{eq:uldcertificate}
 B_Q^{\rm u}
 &:=-J_Q^{\rm u}M_{\rm u}
       -M_{\rm u}(J_Q^{\rm u})^\top
       +A_{\rm u}\gamma_{\rm u}P_v\notag\\
 &=
 \begin{pmatrix}
 2\Id_d&(1+\gamma_{\rm u})\Id_d-Q\\
 (1+\gamma_{\rm u})\Id_d-Q&
 (A_{\rm u}+2)\gamma_{\rm u}\Id_d-2Q
 \end{pmatrix}\succeq\Id_{2d}.
\end{align}
Indeed, its quadratic form at $(r,v)\in\R^{2d}$ is bounded below by
\[
 |r|^2+
 \left((A_{\rm u}+2)\gamma_{\rm u}-2K
       -(1+\gamma_{\rm u}+K)^2\right)|v|^2
 =|r|^2+(4\gamma_{\rm u}+1)|v|^2.
\]

Let $H_{\rm u}(x,v):=U(x)+|v|^2/2$ and
$w:=\nabla\log(\dd\nu/\dd\pi^{\rm u}_\eps)$. For a regular density evolving
at frozen temperature, the same Fisher-information differentiation as in
\eqref{eq:fisherentropyderivatives} gives
\begin{align}\label{eq:uldfrozendissipation}
 \frac{\dd}{\dd t}\cH^{\rm u}_\eps(\nu_t)
 ={}&-\int_{\R^{2d}}w^\top B_{\nabla^2U}^{\rm u}w\,\dd\nu_t\notag\\
 &-2\gamma_{\rm u}\eps\sum_{i=1}^d
   \int_{\R^{2d}}(\partial_{v_i}w)^\top
       M_{\rm u}(\partial_{v_i}w)\,\dd\nu_t
 \leq-\int_{\R^{2d}}|w|^2\dd\nu_t .
\end{align}
Tensorization of Assumption~\ref{ass:lsi} with the velocity Gaussian gives
\[
 \cH^{\rm u}_\eps(\nu)
 \leq\left(\|M_{\rm u}\|
       +\frac{A_{\rm u}C_{\rm LS}(\eps)}{\eps}\right)
       \int_{\R^{2d}}|w|^2\dd\nu.
\]
Consequently, for every small $\eta>0$ and all sufficiently large $u$,
the first relation in \eqref{eq:coolingassumption} and the LSI barrier imply
\begin{equation}\label{eq:uldkappalower}
 \kappa_{\rm u}(\eps_u):=
 \left(\|M_{\rm u}\|
       +\frac{A_{\rm u}C_{\rm LS}(\eps_u)}{\eps_u}\right)^{-1}
 \geq c_\eta(u+t_0)^{-\vartheta-\eta}.
\end{equation}

The moment constant is uniform over the restarting laws. Choose $b>0$ small,
set $c=b\gamma_{\rm u}/2$, and add a constant so
that
\[
 \mathcal V_{\rm u}(x,v)
 =1+U(x)+\frac12|v|^2+b\,x\cdot v+c|x|^2
\]
is positive and comparable to $1+|x|^2+|v|^2$. Its generator satisfies
\[
 L^{\rm u}_{\eps}\mathcal V_{\rm u}
 =-(\gamma_{\rm u}-b)|v|^2-b\,x\cdot\nabla U(x)
       +\gamma_{\rm u}\eps d
 \leq-c_0\mathcal V_{\rm u}+C_0,
\]
uniformly for $0<\eps\leq\eps_0$. Thus a second-moment bound at the restart
time propagates uniformly. The same integration-by-parts argument as in
\eqref{eq:gibbsmoments} bounds
$\pi^{\rm u}_\eps(H_{\rm u})$ uniformly.

At a fixed law, differentiation in the temperature gives
\begin{align}\label{eq:uldtemperatureidentity}
 \partial_\eps\cH^{\rm u}_\eps(\nu)
 =-\frac{2}{\eps^2}\int_{\R^{2d}}
       w^\top M_{\rm u}\nabla H_{\rm u}\,\dd\nu
 -\frac{A_{\rm u}}{\eps^2}\KL(\nu\mid\pi^{\rm u}_\eps)
 -\frac{A_{\rm u}}{\eps^3}
       \left(\nu(H_{\rm u})-\pi^{\rm u}_\eps(H_{\rm u})\right).
\end{align}
The uniform moment bounds, Cauchy--Schwarz inequality, and
$\KL\leq\eps\cH^{\rm u}_\eps/A_{\rm u}$ show that the absolute value of
$\eps_u'$ times the right-hand side is at most
\[
 C|\eps_u'|(1+\eps_u^{-3})
       \left(1+\cH^{\rm u}_{\eps_u}(\nu_u)\right)
 \leq C\frac{\log(u+t_0)}{u+t_0}
       \left(1+\cH^{\rm u}_{\eps_u}(\nu_u)\right).
\]
Combining this estimate with \eqref{eq:uldfrozendissipation} and
\eqref{eq:uldkappalower}, and absorbing the coefficient of the functional
into half of the frozen dissipation, yields, after increasing the starting
time,
\begin{equation}\label{eq:uldrestartode}
 \frac{\dd}{\dd u}\cH^{\rm u}_{\eps_u}(\nu_u)
 \leq-c_\eta(u+t_0)^{-\vartheta-\eta}
             \cH^{\rm u}_{\eps_u}(\nu_u)
       +C_\eta(u+t_0)^{-1+\eta}.
\end{equation}
This is the fixed-friction version of the distorted-entropy calculation of
He, Tan, and Wu \cite[Sections~3.2--3.4]{HeTanWu2024}.

On $[s,T]$, where $s\asymp T$ and $T-s\asymp T^\beta$, integration of
\eqref{eq:uldrestartode} gives \eqref{eq:uldrestart}. The initial term is
bounded by
$\cH^{\rm u}_{\eps_s}(\rho)e^{-cT^{\beta-\vartheta-\eta}}$, while the source
convolution satisfies
\[
 C_\eta T^{-1+\eta}\int_0^\infty
 e^{-cT^{-\vartheta-\eta}r}\dd r
 \leq C_\eta T^{-1+\vartheta+2\eta}.
\]

It remains to remove the temporary regularity assumption. The assumed full
relative Fisher information and the second moment imply that the square root
of the density belongs to $H^1(\R^{2d})$. The cutoff, mollification, and
positive-tail construction in Lemma~\ref{lem:closure} therefore gives regular
laws $\rho_n$ converging in $H^1$ at the square-root level and in weighted
$L^1$ at the density level. In the borderline case $d=1$, use
$H^1(\R^2)\hookrightarrow L^p(\R^2)$ for every finite $p$ in place of the
$m\geq3$ Sobolev exponent used there. The entropy and Fisher calculations in
that lemma, with $H,M$ replaced by $H_{\rm u},M_{\rm u}$, then imply
\[
 \cH^{\rm u}_{\eps_s}(\rho_n)
 \longrightarrow\cH^{\rm u}_{\eps_s}(\rho).
\]
The directional Fisher term is lower semicontinuous under weak convergence:
its variational representation is the analogue of \eqref{eq:fisherdual} with
$M$ replaced by the positive semidefinite matrix $M_{\rm u}$. Relative
entropy is lower semicontinuous as well. The moment bound above is uniform in
the approximating sequence, so the constants in \eqref{eq:uldrestartode} are
also uniform. The cutoff localization argument of
Lemma~\ref{lem:closure} justifies the integrated differential inequality for
each $\rho_n$. Passing to the lower limit in that integrated inequality proves
\eqref{eq:uldrestart} for the law stated in the lemma.

For smoothing, write $s'-s=(e_x,e_v)$ and choose a cubic polynomial $q$ on
$[0,1]$ satisfying
\[
 q(0)=e_x,\quad q'(0)=e_v,\quad q(1)=q'(1)=0,
 \qquad \max_{0\leq j\leq2}\left\|q^{(j)}\right\|_\infty\leq C|s-s'|.
\]
If $(X,V)$ starts from $s$, set $X^e=X+q$ and $V^e=V+q'$. The shifted path
starts from $s'$ and has the same endpoint, while its velocity equation has
the additional drift
\[
 c_e(u)=q''(u)+\nabla U(X_u+q(u))-\nabla U(X_u)
                   +\gamma_{\rm u}q'(u).
\]
The bounded Hessian gives $|c_e(u)|\leq C|s-s'|$. Girsanov's theorem and
Pinsker's inequality therefore give, first for
$|s-s'|\leq c\sqrt{\eps_{t+1}}$,
\[
 |P^{\rm u}_{t,t+1}f(s)-P^{\rm u}_{t,t+1}f(s')|
 \leq C\eps_{t+1}^{-1/2}|s-s'|\|f\|_\infty.
\]
Partitioning the line segment from $s$ to $s'$ proves the estimate globally.
The Markov property gives \eqref{eq:uldsmoothing} for $T\geq t+1$, and
synchronous coupling with the globally Lipschitz drift gives the short-time
claim.
\end{proof}

\begin{proof}[Proof of Theorem~\ref{thm:ubu}]
Fix $T=\tau_k$ and choose
\[
 \beta\in(\vartheta,\min\{1,3a/2\}),
\]
which is possible because $a>2\vartheta/3$. Choose $m$ so that
$\tau_m\leq T-T^\beta<\tau_{m+1}$. Then
$T-\tau_m\asymp T^\beta$, $s_j\asymp T$, and $h_j\asymp T^{-a}$ throughout
the comparison window. Lemmas~\ref{lem:ubulocal} and
\ref{lem:ubugaussian}, including the full Fisher-information conclusion of
the latter, allow Lemma~\ref{lem:uldrestart} to be applied at $\tau_m$ and give
\begin{equation}\label{eq:uburestartbound}
 \cH^{\rm u}_{\eps_T}\left(\nu_mP^{\rm u}_{\tau_m,T}\right)
 \leq CT^{3a}(\log T)^2e^{-cT^{\beta-\vartheta-\eta}}
       +C_\eta T^{-1+\vartheta+2\eta},
\end{equation}
where $\nu_m=\mathcal L(\bar S_m)$.

Choose the same smooth event approximation $f$ as in the proof of
Theorem~\ref{thm:discrete}. Lemma~\ref{lem:uldrestart} implies that
$g_j=P^{\rm u}_{\tau_{j+1},T}f$ has Lipschitz constant at most
$C_{\delta,\zeta}\sqrt{\log T}$; the short terminal intervals are handled by
the Lipschitz regularity of $f$. Telescoping the UBU and exact ULD kernels,
then using Lemma~\ref{lem:ubulocal}, gives
\begin{align}\label{eq:ubutelescoping}
 \left|\nu_kf-\nu_mP^{\rm u}_{\tau_m,T}f\right|
 &\leq C\sum_{j=m}^{k-1}\left(
       h_j^{5/2}+\sqrt{\log T}\,h_j^3
       +\frac{h_j^{3/2}}{T\log T}\right)\notag\\
 &\leq C\left(
       T^{\beta-3a/2}
       +\sqrt{\log T}\,T^{\beta-2a}
       +\frac{T^{\beta-1-a/2}}{\log T}\right).
\end{align}
Here the factor $\sqrt{\log T}\asymp\eps_T^{-1/2}$ cancels the
$\sqrt{\eps_{\tau_j}}$ in the leading local stochastic error.

As in the proof of Theorem~\ref{thm:discrete}, Pinsker's inequality,
$\KL(\cdot\mid\pi^{\rm u}_{\eps_T})\leq\eps_T\cH^{\rm u}_{\eps_T}(\cdot)/A_{\rm u}$,
\eqref{eq:uburestartbound}, and the Gibbs tail estimate \eqref{eq:gibbstail}
give
\[
 \nu_mP^{\rm u}_{\tau_m,T}f
 \leq CT^{-(\delta-\zeta-\zeta')/E}
 +CT^{3a/2}\log T\,e^{-cT^{\beta-\vartheta-\eta}/2}
 +C_\eta T^{-(1-\vartheta)/2+\eta}.
\]
Here \eqref{eq:gibbstail} applies to $\pi^{\rm u}_{\eps_T}$ because its
position marginal is also $\mu_{\eps_T}$. The middle term is the
stretched-exponential term; it comes from the square root of the first term
on the right-hand side of \eqref{eq:uburestartbound} and, as after
\eqref{eq:stretchedexponential}, is at most $C_NT^{-N}$ for every fixed $N$. Combining these bounds with
\eqref{eq:ubutelescoping} gives
\begin{align*}
 \mathbb P\left(U(\bar X_k)>\delta\right)\leq C\bigg(&
 T^{-(\delta-\zeta-\zeta')/E}
 +T^{-(1-\vartheta)/2+\eta}
 +T^{-(3a/2-\beta)}\\
 &+\sqrt{\log T}\,T^{-(2a-\beta)}
 +(\log T)^{-1}T^{-(1+a/2-\beta)}+T^{-N}\bigg).
\end{align*}
Let $\beta\downarrow\vartheta$ and
$\zeta,\zeta',\eta\downarrow0$. The exponent $2a-\vartheta$ is larger than
$3a/2-\vartheta$. The temperature-variation exponent is also nonbinding,
since
\[
 1+\frac a2-\vartheta-\frac{1-\vartheta}{2}
 =\frac{1+a-\vartheta}{2}>0.
\]
This proves \eqref{eq:uburate}; the grid
asymptotics give the iteration statement.
\end{proof}

\begin{remark}[Local error comparison between UBU and the third-order scheme]
\label{rem:ubunoisetiming}
Although both UBU and \eqref{eq:midpoint} use one interior force evaluation,
their force arguments have different stochastic regularity. In UBU, noise
enters velocity directly and reaches position at size
$\mathcal{O}(\sqrt\eps\,h^{3/2})$ before the force kick, producing the
$\mathcal{O}(\sqrt\eps\,h^{5/2})$ term in \eqref{eq:ubulocal}. In the third-order
chain, noise enters $Z$ and reaches position only after two further ordinary
integrations, at size $\mathcal{O}(\sqrt\eps\,h^{5/2})$; its contribution to the force
integral is therefore $\mathcal{O}(\sqrt\eps\,h^{7/2})$, below the deterministic
$\mathcal{O}(h^3)$ endpoint error. The comparison is between the sufficient conditions
proved by the present strong-coupling argument. A weak-error or martingale
analysis that exploits the centering of the UBU $h^{5/2}$ term could produce
a sharper UBU bound.
\end{remark}


\end{document}
